\documentclass{article}
\usepackage{alexconfig}
\usepackage{tikz-cd}

\title{Selected Notes from An Introduction to Manifolds by Loring W. Tu}

\begin{document}
\maketitle

\section{Smooth Functions in Euclidean Space}

\subsection{$C^\infty$ Versus Analytic Functions}
\

\begin{definition}[$C^k$ Functions]
  \

  A function $f: U \to R$ is \textbf{$C^k$} at $p$ if every chain of partial derivatives $$\frac{\partial^j f}{\partial x^{i_1} \dots \partial x^{i_j}} = \frac{\partial f}{\partial x^{i_1}}[ \dots [\dots [\frac{\partial f}{\partial x^{i_j}}]]]$$exist and are continuous at $p$ for all coordinates $x^{i_1}, \dots, x^{i_j}$

  If it is true for any number of partial derivatives, we call the function \textbf{$C^\infty$} or \textbf{smooth}.
\end{definition}

\begin{example}
  \

  If $f: \mathbb{R} \to \mathbb{R}$, with $f(x) = x^\frac{1}{3}$, then $f$ is $C^0$ but not $C^1$, since $f'(0) = \frac{1}{3 \sqrt[3]{0^2}}$ is not defined.
\end{example}

\begin{definition}[Real Analytic Functions]
  \

  A function $f$ is \textbf{real analytic} at $p$ if in some neighborhood of $p$, it is equal to $$f(x) = f(p) + \sum_i \frac{\partial f}{\partial x^i}(p)(x^i - p^i) + \frac{1}{2!} \sum_{i, j} \frac{\partial^2 f}{\partial x^i \partial x^j}(p)(x^i - p^i)(x^j-p^j) + \dots$$That sum is called the \textbf{Taylor series}
\end{definition}

\begin{example}
  \

  $\sin(x)$ is real analytic.
\end{example}

\subsection{Taylor's Theorem With Remainder}
\

\begin{definition}[Star Shaped Sets]
  \

  A subset $S \subseteq \mathbb{R}^n$ is \textbf{star-shaped} with respect to some point $p \in S$ if for every $x \in S$, the line segment between $p$ and $x$ lies entirely within $S$
\end{definition}

\begin{lemma}[Taylor's Theorem With Remainder]
  \

  Let $f$ be $C^\infty$ on an open subset $U \subseteq \mathbb{R}^n$ which is star shaped, with respect to a point $p$. Then, there are $C^\infty$ functinos $g_1(x) \dots g_n(x)$ on $U$ such that $$f(x) = f(p) + \sum_{i=1}^n (x^i - p^i)g_i(x)$$with $$g_i(p) = \frac{\partial f}{\partial x^i}(p)$$
\end{lemma}

\subsection{Exercises}
\

\begin{exercise}[1.1]
  \

  Find a function $h: \mathbb{R} \to \mathbb{R}$ that is $C^2$ but not $C^3$.
\end{exercise}

\begin{solution}
  \

  Consider $f(x) = x^\frac{5}{2}$, $f'(x) = \frac{5}{2}x^\frac{3}{2}$ defined at $0$, $f''(x) = \frac{15}{4}x^\frac{1}{2}$ also defined, but $f''' = \frac{5}{2} x^{-\frac{1}{x}}$ not defined at $0$
\end{solution}

\begin{exercise}[1.7]
  \

  Let $f = x^3$. Show that $f$ is a bijective $C^\infty$ map but that $f ^{-1}$ is not $C^\infty$.
\end{exercise}

\begin{solution}
  \

  $f$ is bijective since for any $x \in \mathbb{R}$, $f(\sqrt[3]{x}) = x$, and this is indeed unique. However, $f ^{-1}(x) = \sqrt[3]{x}$ is not $C^\infty$, since ${f ^{-1}}'(x) = \frac{1}{3\sqrt[3]{x^2}}$ is not defined (and therefore not continuous) at $0$.
\end{solution}

\section{Tangent Vectors on $\mathbb{R}^n$ as Derivations}

\subsection{The Directional Derivative}
\

\begin{definition}[Tangent Space]
  \

  The \textbf{tangent space} of $\mathbb{R}^n$ at a point $p$, denoted $T_p\mathbb{R}^n$ is the vector space composed of all vectors emanating from $p$.

\end{definition}

\begin{remark}
  \

  This definition will not work for manifolds, so we will introduce a more general definition and show it is equivalent for any $\mathbb{R}^n$.
\end{remark}

\begin{proposition}
  \

  The line through the point $p \in \mathbb{R}^n$, pointing in some direction $v \in T_p \mathbb{R}^n$, is parameterized by $$c(t) = (p^1 + tv^1, \dots, p^n + tv^n)$$
\end{proposition}

\begin{definition}[Directional Derivative]
  \

  For some vector $v \in T_p \mathbb{R}^n$, and some function $f: \mathbb{R}^n \to \mathbb{R}$, we define the \textbf{directional derivative} of $f$ with respect to $v$ as $$D_v(f) = \left. \frac{d f(c(t))}{dt} \right\vert_{t=0}$$
\end{definition}

\begin{proposition}
  \

  By chain rule, we can see that $$D_v(f) = \left.\frac{df(c(t))}{dt}\right\vert_{t=0} = \sum_{k=1}^n \left.\frac{dc^k}{dt}\right\vert_{t=0} \left.\frac{\partial f}{\partial x^i}\right\vert_p$$but since $c$ is linear, $\frac{dc^k}{dt} = v^k$, so we have $$D_v(f) = \sum_{k=1}^n v^k \left.\frac{\partial f}{\partial x^k}\right\vert_p$$
\end{proposition}

\subsection{Germs of Functions}
\

\begin{definition}[Germs]
  \

  Consider pairs $(f,U)$ where $f: \mathbb{R}^n \to \mathbb{R}$ and $p \in U \subseteq \mathbb{R}^n$ for some point $p$. Then, two pairs $(f,U)$ and $(g,V)$ are equivalent if $f \vert_{U \cap V} = g \vert_{U \cap V}$, and this is an equivalence relation.

  The equivalence classes of such form are called \textbf{germs}, and the set of all germs at some point $p$ is denoted $C^\infty_p(\mathbb{R}^n)$.
\end{definition}

\subsection{Derivations of Functions at a Point}
\

\begin{proposition}
  \

  $D_v$ as defined above is $\mathbb{R}$-linear and satisfies the Leibniz rule, $$D_v(f,g) = D_v(f)g + fD_v(g)$$
\end{proposition}

\begin{definition}[Space of Derivations]
  \

  Any linear map $D: C^\infty_p \to \mathbb{R}$ satisfying the Leibniz rule at $p$ is called a \textbf{point derivation}. The space of all such point derivations is denoted $\mathcal{D}_p$ and forms a vector space.
\end{definition}

\begin{proposition}
  \

  The map $\varphi: \mathcal{D}_p(\mathbb{R}^n) \to T_p(\mathbb{R}^n)$ which takes $$\sum_{k=1}^n v^k \left.\frac{\partial}{\partial x^k}\right\vert_p \mapsto v$$ is bijective.
\end{proposition}

\begin{customproof}
  \

  For any $v \in T_p \mathbb{R}^n$, there is a unique derivation $\sum_{k=1}^n v^k \left.\frac{\partial}{\partial x^k}\right\vert_p$
\end{customproof}

\subsection{Vector Fields}
\

\begin{definition}[Vector Fields]
  \

  A \textbf{vector field} $X$ on an open subset $U \subseteq \mathbb{R}^n$ is a function which assigns to each point $p \in U$, a tangent vector $X_p \in T_p \mathbb{R}^n$. Since $T_p \mathbb{R}^n$ has a basis $\{\frac{\partial }{\partial x^i}\}_{1 \leq i \leq n}$, we can write a vector field as $$X_p = \sum_{i = 1}^n a^i(p)\left.\frac{\partial}{\partial x^i}\right\vert_p$$where $a^i(p)$ just outputs the $i$th coefficient for the vector assigned to $p$.

  If all $a^i$ are smooth on $U$, then we call $X$ a \textbf{smooth vector field}.
\end{definition}

\begin{proposition}
  \

  The set of all smooth vector fields on some $U$, denoted $\mathcal{X}(U)$, form an $\mathbb{R}$-module.
\end{proposition}

\subsection{Vector Fields as Derivatives}
\

\begin{definition}
  \

  If $X$ is a vector field on $U$, $f: U \to \mathbb{R}$, define $Xf: \mathbb{R}^n \to \mathbb{R}$ by $$Xf(p) = X_p f$$for all $p \in U$.
\end{definition}

\begin{definition}[Derivations of Algebras]
  \

  If $A$ is an algebra over a field $k$, then a \textbf{derivation} of $A$ is a $k$-linear map $D:A \to A$ such that $$D(ab) = aD(b) + D(a)b$$ for all $a,b \in A$. The set of all such derivations forms a $k$-vector space, denoted $\text{Der}(A)$
\end{definition}

\begin{proposition}
  \

  There is a bijection $$\varphi: \mathcal{X}(U) \cong \text{Der}(C^\infty(U))$$$$\varphi: X \mapsto [f \mapsto Xf]$$
\end{proposition}

\subsection{Exercises}
\

\begin{exercise}[2.1]
  \

  Let $X$ be the vector field $x \frac{\partial }{\partial x} + y \frac{\partial }{\partial y}$, let $f = x^2 + y^2 + z^2$ on $\mathbb{R}^2$. Compute $Xf$
\end{exercise}

\begin{solution}
  \

  $$(Xf)(p) = x\frac{\partial f}{\partial x} + y\frac{\partial f}{\partial y} = 2x^2 + 2y^2$$
\end{solution}

\section{Exterior Algebra of Multicovectors}

\subsection{Dual Spaces}
\

\begin{definition}[Dual Space]
  \

  If $V$ is a vector space, the \textbf{dual space} $V^*$ is the space of all linear functionals $$V^* = \text{Hom}(V, \mathbb{R})$$Elements of $V^*$ are called \textbf{$1$-covectors}.
\end{definition}

\begin{proposition}
  \

  The functions $\alpha^i(e_j) = \delta^i_j$ form a basis for $V^*$
\end{proposition}

\begin{corollary}
  \

  For a finite dimensional vector space $V$, $$\dim V^* = \dim V$$
\end{corollary}

\subsection{Permutations}
\

\begin{definition}[Even and Odd Permutations]
  \

  A permutation $\sigma \in S_n$ is \textbf{even} if it is the product of an even number of transpositions, it is \textbf{odd} otherwise.

  The \textbf{sign function} is given by $$\text{sgn}(\sigma) =
  \begin{cases}
    1 & \sigma \text{ is even}\\
    -1 & \sigma \text{ is odd}
  \end{cases} $$
\end{definition}

\begin{definition}[Inversions]
  \

  An \textbf{inversion} is an ordered pair $(\sigma(i), \sigma(j))$ with $i<j$, but $\sigma(i) >\sigma(j)$.
\end{definition}

\begin{proposition}
  \

  A permutation is even iff it has an even number of inversions.
\end{proposition}

\subsection{Multilinear Functions}
\

\begin{definition}[$k$-linear function]
  \

  If $V$ is a $\mathbb{R}$-vector space, a function $f: V^k \to \mathbb{R}$ is $k$-linear if $$f(v_1, v_2, \dots, v_{i-1}, ax + by, v_{i+1}, \dots, v_k) = af(v_1, \dots, v_{i-1}, x, v_{i+1}, \dots, v_k) + bf(v_1, \dots, v_{i-1}, y, v_{i+1}, \dots, v_k)$$for any term $v_i = ax + by$.
\end{definition}

\begin{definition}[Symmetric and Alternating Functions]
  \

  A $k$-linear function $f$ is \textbf{symmetric} if for any $\sigma \in S_k$, $$f(v_1, \dots, v_k) = f(v_{\sigma(1)}, \dots, v_{\sigma(k)})$$so reordering terms does not change the function.

  A $k$-linear function is \textbf{alternating} if for any $\sigma \in S_k$, $$f(v_1, \dots, v_k) = \text{sgn}(\sigma)f(v_{\sigma(1)}, \dots, v_{\sigma(k)})$$
\end{definition}

\begin{example}
  \

  \begin{enumerate}
    \item $f(u,v) = u \cdot v$ is symmetric
    \item $f(v_1, \dots, v_n) = \det[v_1, \dots, v_n]$, where $v_i \in \mathbb{R}^n$, is alternating
    \item $u \times v$ on $\mathbb{R}^3$ is alternating
  \end{enumerate}
\end{example}

\begin{definition}[Space of Functions]
  \

  The space of all $k$-linear functionals $f:V^k \to \mathbb{R}$ is denoted $L_k(V)$.

  The space of alternating $k$-linear functionals $f: V^k \to \mathbb{R}$ is denoted $\bigwedge_k(V)$. It's elements are called \textbf{alternating $k$-tensors}, \textbf{$k$-covectors}, or \textbf{multicovectors of degree $k$}

  When $k = 0$, $\bigwedge_0(V) \cong \mathbb{R}$ is the space of constants, and $A_1(V) \cong V^*$
\end{definition}

\subsection{The Permutation Action on Multilinear Functions}
\

\begin{definition}
  \

  Let $f:V^k \to \mathbb{R}$ be a $k$-linear function, $\sigma \in S_k$. Then, define $$(\sigma f)(v_1, \dots, v_k) = f(v_{\sigma(1)}, \dots, v_{\sigma(k)})$$
\end{definition}

\begin{proposition}
  \

  $f$ is symmetric iff $\sigma f = f$ for all $\sigma$, and $f$ is alternating iff $\sigma f = \text{sgn}(\sigma)f$ for all $\sigma$.
\end{proposition}

\begin{lemma}
  \

  If $\sigma, \tau \in S_k$, $f$ is a $k$-linear function on $V$, then $\tau(\sigma f) = (\tau\sigma)f$
\end{lemma}

\begin{customproof}
  \

  $$\tau(\sigma f) = \tau f(v_{\sigma(1)}, \dots, v_{\sigma(k)}) = f(v_{\tau(\sigma(1))}, \dots, v_{\tau(\sigma_k)}) = f(v_{\tau\sigma(1)}, \dots, v_{\tau\sigma(k)}) = (\tau\sigma)f$$
\end{customproof}

\subsection{The Symmetrizing and Alternizing Operators}
\

\begin{definition}[The Symmetrizing Operator]
  \

  Given some $k$-linear function $f: V^k \to \mathbb{R}$, define $S: L_k(V) \to L_k(V)$ by $$Sf(v_1, \dots, v_k) = \sum_{\sigma \in S_k} f(v_{\sigma(1)}, \dots, v_{\sigma(k)}) = \sum_{\sigma \in S_k} \sigma f$$and this turns $f$ into a symmetric functional.
\end{definition}

\begin{definition}[The Alternizing Operator]
  \

  Define $A: L_k(V) \to \bigwedge_k(V)$ by $$Af = \sum_{\sigma \in S_k} \text{sgn}(\sigma) \sigma f$$and this turns $f$ into an alternating functional.
\end{definition}

\begin{proposition}
  \

  If $f$ is a $k$-linear functional:

  \begin{enumerate}
    \item $Sf$ is symmetric
    \item $Af$ is alternating
  \end{enumerate}
\end{proposition}

\begin{proposition}
  \

  If $f \in A_k(V)$ then $$Af = k!f$$
\end{proposition}

\begin{customproof}
  \

  $$Af = \sum_{\sigma \in S_k} \text{sgn}(\sigma)\text{sgn}(\sigma) f = \sum_{\sigma \in S_k} 1 \cdot f = k!f$$since there are $k!$ permutations in $S_k$.
\end{customproof}

\subsection{The Tensor Product}
\

\begin{definition}[Tensor Product]
  \

  Let $f$ be a $k$-linear function, $g$ be an $l$-linear function. Then, their \textbf{tensor product} $f\otimes g$ is defined by $$(f \otimes g)(v_1, \dots, v_{k+l}) = f(v_1, \dots, v_k)g(v_{k+1},\dots,v_{k+l})$$By commutativity of multiplication of real numbers, the tensor product is commutative.
\end{definition}

\begin{example}
  \

  Consider the inner product of two vectors, $\langle v, w \rangle$. Then, $$\langle v, w\rangle = \sum v^iw^i\langle e_i, e_j\rangle$$but $v^i = \alpha^i(v)$ where $\alpha^i$ is a basis element for the dual space, since they are functions which pick out the $i$th coordinate. Then, we have $$\langle v,w \rangle = \sum \alpha^i(v) \alpha^j(w)\langle e_i, e_j \rangle = \sum \langle e_i, e_j\rangle (\alpha^i \otimes\alpha^j)(v,w)$$
\end{example}

\subsection{The Wedge Product}
\

\begin{definition}[Wedge Product]
  \

  Let $f \in \bigwedge_k(V)$, $g \in \bigwedge_l(V)$ alternating functions. We define the \textbf{wedge product} or the \textbf{exterior product} as $$f \wedge g = \frac{A(f \otimes g)}{k!l!}$$
\end{definition}

\begin{proposition}
  \

  The wedge product is anticommutative: $$f \wedge g = - g\wedge f$$and associative. The tensor product is commutative and associative.
\end{proposition}

\begin{proposition}
  \

  If $\alpha^1\dots \alpha^k$ are linear functions on a vector space $V$, and $v_1, \dots,v_k \in V$, then $$(a^1 \wedge \dots \wedge a^k)(v_1, \dots, v_k) = \text{det}[a^i(v_j)]$$
\end{proposition}

\begin{customproof}
  \

  $(a^1 \wedge \dots \wedge a^k)(v_1,\dots, v_j) = \sum_{\sigma \in S_k} \text{sgn}(\sigma) \alpha^1(v_{\sigma(1)}) \dots \alpha^k(v_{\sigma(k)}) = \det[\alpha^i(v_j)]$
\end{customproof}

\begin{definition}[Graded Algebras]
  \

  An algebra $A$ over a field $k$ is \textbf{graded} if $A = \bigoplus_{i=0}^\infty A^i$ for $k$-vector spaces (the superscripts are just labels, they don't mean anything else), and for $a \in A^i$, $b \in A^j$, $ab \in A^{i+j}$.

  A graded algebra is said to be \textbf{graded commutative} or \textbf{anticommutativt} or \textbf{supercommutative} if for $a \in A^i, b \in A^j$, $ab = (-1)^{i+j}ba$.
\end{definition}

\begin{remark}
  \

  Note that graded algebras can be regularly commutative, where $ab = ba$, or strictly graded commutative, where $ab = (-1)^{i+j}ba$ and $x^2 = 0$ for all $x \in A^i$, $i$ odd, which is implied if $k$ isn't a field of characteristic $2$.
\end{remark}

\begin{definition}[Exterior/Grassmann Algebra]
  \

  The \textbf{exterior algebra} or \textbf{Grassmann algebra} over a vector space $V$, where $\dim V= n$, is given by $$\bigwedge (V) = \bigoplus_{i=0}^\infty \bigwedge_i(V) = \bigoplus_{i=0}^n \bigwedge_i(V)$$
\end{definition}

\begin{remark}
  \

  We can do the sum up to $\dim V = n$ because $\bigwedge_k(V)$ is trivial when $k > n$. This is because any set of $k$ vectors is linearly dependent, and we can do $k$-linear magic and anticommutativiity to show that $f = -f = 0$.
\end{remark}

\subsection{A Basis for $k$-covectors}
\

\begin{remark}[Multi-Index Notation]
  \

  Let $I = (i_1, \dots, i_k)$, then $e_I = (e_{i_1},\dots, e_{i_k})$, and $\alpha^i = \alpha^{i_1} \wedge \dots \wedge \alpha^{i_k}$.
\end{remark}

\begin{lemma}
  \

  Let $e_1, \dots, e_n$ be a basis for $V$, $\alpha^1, \dots, \alpha^n$ be the dual basis in $V^*$. If $I = (1 \leq i_1 \leq \dots \leq i_k \leq n)$ and $J = (1 \leq j_1 \leq \dots \leq j_k \leq n)$, then $$\alpha^I(e_J) = \delta_J^I =
  \begin{cases}
    1 & \text{for $I = J$}\\
    0 & \text{for $I \neq J$}
  \end{cases}$$
\end{lemma}

\begin{customproof}
  \

  $\alpha^I(e_J) = \det[\alpha^i(e_j)]_{i \in I, j \in J}$, but if these are not exactly the same, there will be one row/column that is all $0$, hence the determinant will be all $0$.
\end{customproof}

\begin{lemma}
  \

  The alternating linear functions $\alpha^I$, $I = (i_1 < \dots < i_k)$ form a basis on $\bigwedge_k(V)$
\end{lemma}

\begin{corollary}
  \

  If a vector space $V$ has dimension $n$, $\bigwedge_k(V)$ has dimension $\binom{n}{k}$
\end{corollary}

\subsection{Exercises}
\

\begin{exercise}[3.2]
  \

  \begin{enumerate}
    \item Show that if $f: V \to \mathbb{R}$ is a linear functional, $\dim V = n$, then $\dim \ker f = n-1$. A $n-1$ dimensional subspace of a vector space is called a \textbf{hyperplane}.
    \item Show that a nonzero linear functional on a vector space $V$ is determined up to a multiplicative constant, by it's kernel, a hyperplane in $V$. In other words, if $f, g: V \to \mathbb{R}$, and $\ker f = \ker g$, then $g = cf$ for some $c \in \mathbb{R}$.
  \end{enumerate}
\end{exercise}

\begin{customproof}
  \

  \begin{enumerate}
    \item By the first isomorphism theorem of vector spaces, $\text{Im}(f) \cong V/\ker(f)$, so $\mathbb{R} \cong V/\ker(f)$ and $V/\ker(f)$ has dimension $1$. But $V = V/\ker(f) \oplus \ker(f)$, so $\dim \ker(f) = n-1$.
    \item Since $\ker f$ is a hyperplane, any vector $x$ can be written as $x = h = \alpha v_0$ with $h \in \ker f$, and $v_0$ a point where $f(v_0) \neq 0$ (this is because the kernel and one vector not in the kernel form a basis for the entire vector space.)

      We get $$f(x) = f(h) + \alpha f(v_0) = \alpha f(v_0)$$$$g(x) =  \alpha g(v_0)$$so $\alpha = \frac{f(x)}{f(v_0)}$ nonzero and well defined, and $g(x) = \frac{f(x)g(v_0)}{f(v_0)} = f(x)\frac{g(v_0)}{f(v_0)}$
  \end{enumerate}
\end{customproof}

\section{Differential Forms on $\mathbb{R}^n$}

\subsection{Differential $1$-forms and the Differential of a Function}
\

\begin{definition}[Cotangent Space]
  \

  The \textbf{cotangent space} denoted by $T^*_p \mathbb{R}^n$ is the dual space of the tangent space, $$T^*_p \mathbb{R}^n \cong (T_p \mathbb{R}^n)^* \cong \text{Hom}(T_p \mathbb{R}^n, \mathbb{R})$$
\end{definition}

\begin{definition}[Differential $1$-forms]
  \

  Just like how vector fields assign a tangent vector to each point of $\mathbb{R}^n$, a \textbf{differential $1$-form} is a function that assigns a cotangent vector to each point.
\end{definition}

\begin{example}
  \

  For any $C^\infty$ function $f: U \to \mathbb{R}$ we can construct a $1$-form $df$, defined as $$(df)_p(X_p) = X_pf$$
\end{example}

\begin{proposition}
  \

  If $x^1, \dots, x^n$ are the standard coordinates for $\mathbb{R}^n$, then at each point $p \in \mathbb{R}^n$, $\{(dx^1)_p, \dots, (dx^n)_p\}$ forms a basis $T^*_p \mathbb{R}^n$.
\end{proposition}

\begin{proposition}
  \

  If $f \to U \mathbb{R}$ is a $C^\infty$ function on an open set $U$, then $$df = \sum \frac{\partial f}{\partial x^i} dx^i$$
\end{proposition}

\subsection{Differential $k$-forms}
\

\begin{definition}[Differential $k$-forms]
  \

  In general, a \textbf{differential $k$-form} is a function that assigns to each $p \in U$ an alternating $k$-linear function, so if $\omega$ is a differential $k$-form, then $$\omega_p \in \bigwedge_k(T_p \mathbb{R}^n)$$and since bases for $\bigwedge_k(T_p \mathbb{R}^n)$ are $dx_p^I$ where $I = \{i_j\}$, $1 \leq i_1\leq \dots \leq i_k \leq n$, we get $$\omega_p = \sum \alpha_I(p) dx^I_p$$and so $$\omega = \sum \alpha_Idx^I$$(here $\alpha_I$ is a single function $\alpha_I: \mathbb{R}^n \to \mathbb{R}$, but $dx^I$ is the wedge $dx^{i_1} \wedge \dots \wedge dx^{i_k}$).

  If each $a^I$ is $C^\infty$, then $\omega$ is a $C^\infty$ form.

  $\Omega^k(U)$ is the vector space of $C^\infty$ $k$-forms on $U$.
\end{definition}

\begin{definition}[Wedge Product of Forms]
  \

  Let $\omega \in \Omega^k(U)$, $\tau \in \Omega^l(U)$. The \textbf{wedge product} is given pointwise by $$(\omega \wedge \tau)_p = \omega_p \wedge \tau_p$$In coordinates, if $\omega = \sum_{I} a_I dx^I$ and $\tau = \sum_J a_J dx^J$, then $$\omega \wedge \tau = \sum_{I,J} (a_Ib_J) dx^I \wedge dx^J$$

  The wedge product is a bilinear map $$\wedge: \Omega^k(U) \times \Omega^l(U) \to \Omega^{k+l}(U)$$
\end{definition}

\subsection{Differential Forms as Multilinear Functions on Vector Fields}
\

\begin{definition}
  \

  If $\omega$ is a $1$-form, $X$ is a $C^\infty$ vector field on an open set $U$ in $\mathbb{R}^n$, then define the function $$\omega(X)_p = \omega_p(X_p)$$or in coordinates, $$\omega(X) = (\sum_i a_idx^i)(\sum_j b^j \frac{\partial}{\partial x^j}) = \sum a_ib^i$$
\end{definition}

\begin{remark}
  \

  So any $1$-form $\omega$ defines a (linear) map $\mathcal{X}(U) \to C^\infty(U)$ taking vector fields to functions on $U$.
\end{remark}

\subsection{The Exterior Derivative}
\

\begin{definition}[Exterior Derivative]
  \

  For a $0$-form (a function) $f \in C^\infty(U)$, the \textbf{exterior derivative} is given by $$df = \sum \frac{\partial f}{\partial x^i} d x^i \in \Omega^1(U)$$

  For any $k$-form $\omega \in \Omega^k(U)$, with $k \geq 1$, the \textbf{exterior derivative} is given by $$d\omega = \sum_I da_I\wedge dx^I = \sum_I (\sum_j \frac{\partial a_I}{\partial x^j} dx^j) \wedge dx^I \in \Omega^{k+1}(U)$$
\end{definition}

\begin{example}
  \

  Let $\omega = fdx + gdy \in \Omega^1(\mathbb{R}^2)$, where $f,g \in C^\infty(\mathbb{R}^2)$, then $$d\omega = df \wedge dx + dg \wedge dy $$but in the definnition above, $df = \sum \frac{\partial f}{\partial x^i} dx^i$, so we have $$d\omega = (\frac{\partial f}{\partial x} dx + \frac{\partial f}{\partial y} dy) \wedge dx +(\frac{\partial g}{\partial x} dx + \frac{\partial g}{\partial y} dy) \wedge dy$$so cancelling, $$d\omega = (\frac{\partial g}{\partial x} - \frac{\partial f}{\partial y}) dx \wedge dy$$
\end{example}

\begin{definition}[Antiderivation]
  \

  Let $A = \bigoplus_{i=0}^\infty V^i$ be a graded algebra over a field $k$. An \textbf{antiderivation} on $A$ is a $k$-linear map $$D: A \to A$$such that for $a \in A^i$, $b \in A^l$, $$D(ab) = (Da)b + (-1)^iaDb$$If there is some integer $m$ such that an antiderivation sends $a \in A^i$ to $A^{i+m}$, we call it an antiderivation of \textbf{degree $m$}
\end{definition}

\begin{proposition}
  \

  \begin{enumerate}
    \item The exterior differential $d: \Omega^*(U) \to \Omega^*(U)$ is an antiderivation of degree $1$: $$d(\omega \wedge \tau) = (d\omega) \wedge \tau + (-1)^{\text{deg} \omega} \omega \wedge d\tau$$
    \item $d^2 = 0$
    \item If $f \in C^\infty(U)$, and $X \in \mathcal{X}(U)$ a vector field, then $(df)(X) = Xf$.
  \end{enumerate}
\end{proposition}

\begin{customproof}
  \

  \begin{enumerate}
    \item Since both sides are linear, we can reduce this to the case where $\omega = fdx^I$, $\tau = gdx^J$, and this will work for any sum by linearity.

      We have $$d(\omega \wedge \tau) = d(fg dx^I \wedge dx^J) = \sum \frac{\partial fg}{\partial x^i} dx^i \wedge dx^I \wedge dx^J = \sum g\frac{\partial f}{\partial x^i} dx^i \wedge dx^I \wedge dx^J + \sum f\frac{\partial g}{\partial x^i} dx^i \wedge dx^I \wedge dx^J$$and moving the $\frac{\partial g}{\partial x^i} dx^i$ across gives us a $(-1)^k$ term$$= \sum g\frac{\partial f}{\partial x^i} dx^i \wedge dx^I \wedge dx^J + (-1)^k \sum f dx^I \wedge \frac{\partial g}{\partial x^i} dx^i \wedge dx^J = d\omega \wedge \tau + (-1)^k \omega \wedge d\tau$$
    \item Again assume WLOG $\omega = fdx^I$. Then, $$d^2(fdx^I) = d(\sum \frac{\partial f}{\partial x^i} dx^i \wedge dx^I) = \sum \frac{\partial^2 f}{\partial x^j \partial x^i} dx^j \wedge dx^i \wedge d^I$$$\frac{\partial^2 f}{\partial x^j \partial x^i}$ is symmetric, but $dx^j \wedge dx^i$ is antisymmetric, so when $i \neq j$, the $ij$ term and the $ji$ term will add up and cancel. When $i = j$, $dx^i \wedge dx^j = 0$.
    \item For functions, the exterior derivative is just the differential, and we defined $(df)(X) = Xf$
  \end{enumerate}
\end{customproof}

\subsection{Closed and Exact Forms}
\

\begin{definition}[Closed Forms]
  \

  A form is \textbf{closed} if $d\omega= 0$.
\end{definition}

\begin{definition}[Exact Forms]
  \

  A form $\omega$ is \textbf{exact} if there is some $\tau$ such that $d\tau = \omega$
\end{definition}

\begin{proposition}
  \

  Every exact form is closed
\end{proposition}

\begin{customproof}
  \

  Let $\omega$ be an exact, form, $\omega = d\tau$ for some form $\tau$. Then, $d\omega= d^2\tau = 0$ so $\omega$ is closed.
\end{customproof}

\begin{definition}[Cochain Complex]
  \

  A collection of vector spaces $\{V^k\}_{k = 0}^\infty$ with linear maps $d_k: V^k \to V^{k+1}$ such that $d_{k+1} \circ d_k = 0$ is called a \textbf{differential complex} or a \textbf{cochain complex}.
\end{definition}

\begin{definition}[deRham Complex]
  \

  For any open subset $U \subseteq \mathbb{R}^n$, the exterior derivative $d$ gives the cochain complex \[
    \begin{tikzcd}
      0 & \Omega^0(U) & \Omega^1(U) & \Omega^2(U)  \dots
      \arrow[from=1-1, to=1-2]
      \arrow[from=1-2, to=1-3, "d"]
      \arrow[from=1-3, to=1-4, "d"]
  \end{tikzcd}\]called the \textbf{deRham complex} and the closed forms are the kernel of $d$, and exact forms are the image of $d$.
\end{definition}

\subsection{Exercises}
\

\begin{exercise}[4.1]
  \

  Let $\omega$ be the $1$-form $zdx - dz$ and let $X$ be the vector field $y \frac{\partial }{\partial x} + x\frac{\partial}{\partial y}$ on $\mathbb{R}^3$. Compute $\omega(X)$ and $d\omega$
\end{exercise}

\begin{solution}
  \

  $$d \omega = dz \wedge dx + d(-1) \wedge dz = - dx \wedge dz$$

  $$\omega(X) = \sum a_ib^i = yz$$
\end{solution}

\begin{exercise}[4.2]
  \

  At each point $p \in \mathbb{R}^3$, define a bilinear function $\omega_p$ on $T_p \mathbb{R}^3$ by $$\omega_p(a,b) = p^3 \text{det}
  \begin{bmatrix}
    a^1 & b^1\\
    a^2 & b^2
  \end{bmatrix}$$where $a^1, a^2, b^1, b^2$ are components of $a,b$ and $p^3$ is the third component of $p$.

  Since $\omega_p$ is an alternating bilinear function, it is a $2$-form. Write it in terms of the standard basis $dx^i \wedge dx^j$ at each point.
\end{exercise}

\begin{solution}
  \

  Since $(dx^1 \wedge dx^2)(a,b) = \frac{A(x^1 \otimes x^2)}{1!1!} = \sum_{\sigma \in S_2} x^1(a_{\sigma(1)})x^2(a_{\sigma(2)}) = a^1b^2 - a^2b^1$, we get:

  $\omega_p(a,b) = p^3(a^1b^2 - b^1a^2) = x^3(p) dx^1 \wedge dx^2$
\end{solution}

\section{Manifolds}

\subsection{Topological Manifolds}
\

\begin{definition}[Second Countable]
  \

  A topological space $X$ is \textbf{second countable} if it has a basis which is countable.
\end{definition}

\begin{definition}[Locally Euclidean]
  \

  A topological space $M$ is \textbf{locally Euclidean} with dimension $n$ if every point $p \in M$ has a neighborhood $U$ such that there is a homeomorphism $\varphi: U \to \mathbb{R}^n$, called a \textbf{chart}.
\end{definition}

\begin{definition}[Topological Manifolds]
  \

  A topological space which is Hausdorff, second countable, and locally Euclidean is called a \textbf{topological manifold}.
\end{definition}

\begin{proposition}[Invariance of Dimensino]
  \

  If a topological manifold is of dimension $m$ on some neighborhood, it must be of dimension $m$ on every neighborhood.
\end{proposition}

\begin{example}
  \

  A cross is not a topological manifold. Pick $p$ to be the intersection point, neighborhood $U$. It cannot be homeomorphic to any $\mathbb{R}^n$, since if we remove some point from a neighbhorhood of $\mathbb{R}^n$, it has either $2$ connected components (if $n=1$) or one connected component. But if we remove $p$, we see that the neighborhood $U$ has four connected components. Hence, a chart at $p$ cannot exist.
\end{example}

\subsection{Compatible Charts}
\

\begin{definition}[Compatible Charts]
  \

  Two charts $(U, \varphi), (V, \psi)$ of a topological manifold are \textbf{compatible} if $\varphi \circ \psi ^{-1}$ and $\psi \circ \varphi ^{-1}$ are smooth on $U \cap V$
\end{definition}

\begin{definition}[Atlas]
  \

  An \textbf{atlas} on a locally Euclidean space $M$ is a collection $\mathcal{A} = \{(U_\alpha, \varphi_\alpha)\}$ of charts that are pairwise compatible, and where $\bigcup_\alpha U_\alpha = M$
\end{definition}

\begin{proposition}
  \

  Let $\{(U_\alpha, \varphi_\alpha)\}$ be an atlas on a locally Euclidean space. If two charts $(V, \psi)$, $(W, \sigma)$ are compatible with the atlas, they are compatible with each other.
\end{proposition}

\begin{customproof}
  \

  If they are disjoint, then they are compatible. If not, pick some $(U_i, \varphi_i)$ such that $U_i$ forms an open cover of $V \cap W$, then if they are compatible with each $U_i$, they are compatible with each other by transitivity.
\end{customproof}

\subsection{Smooth Manifolds}
\

\begin{definition}[Maximal Atlases]
  \

  An atlas $\mathcal{A}$ is \textbf{maximal} if it is not contained in any larger atlas (meaning we can't insert another chart and still have it be compatible).
\end{definition}

\begin{definition}[Smooth Manifold]
  \

  A \textbf{smooth manifold} $M$ is a topological manifold equipped with a maximal atlas. The maximal atlas is also called the \textbf{differentiable structure} on $M$.
\end{definition}

\begin{proposition}
  \

  Any atlas $\mathcal{A}$ is contained in a unique atlas.
\end{proposition}

\begin{proposition}
  \

  If $\{(U_\alpha, \varphi_\alpha)\}$, and $\{(V_\beta, \psi_\beta)\}$ are two atlases for smooth manifolds $M, N$ respectively, then the collection $\{(U_\alpha \times V_\beta, \varphi_\alpha \times \psi_\beta)\}$ forms an atlas on $M \times N$
\end{proposition}

\subsection{Exercises}
\

\begin{exercise}[5.2]
  \

  Show that the sphere with a hair (some line coming out of it at a point $q$) is not locally Euclidean at $q$, and hence cannot be a manifold.
\end{exercise}

\begin{solution}
  \

  Suppose it was, and there was some neighborhood $q \in U$, and some chart $\varphi: U \to \mathbb{R}^n$. Then, on this chart, by removing $q$, we see that the preimage has two connected components, the sphere, and the hair, so $\varphi: U \to \mathbb{R}^1$, as it is the only Euclidean space where removing one point gets you two connected segments. However, the sphere at everywhere else is locally Euclidean to $\mathbb{R}^2$. By invariance of dimension of manifolds, this cannot be a topological manifold.
\end{solution}

\begin{exercise}[5.4]
  \

  Let $\{(U_\alpha, \varphi_\alpha)\}$ be the maximal atlas on a manifold $M$. For any open set $U$ in $M$ and a point $p \in U$, prove the existence of a coordinate open set $U_\alpha$ such that $p \in U_\alpha \subseteq U$.
\end{exercise}

\begin{solution}
  \

  Suppose not. Then, pick some chart $(U_\alpha, \varphi_\alpha)$ with $p \in U_\alpha$. Now, consider the chart $(U_\alpha \cap U, \varphi_\alpha \vert_{U_\alpha \cap U})$ is a compatible chart since it is just a subset of an existing chart, but it is not in the atlas since $U_\alpha \cap U \subseteq U$, a contradiction.
\end{solution}

\section{Smooth Maps on a Manifold}

\subsection{Smooth Functions on a Manifold}
\

\begin{definition}[Smooth Functions]
  \

  Let $M$ be a manifold, $f: M \to \mathbb{R}$ be a function. $f$ is \textbf{smooth} at $p$ if there is a chart $(U, \varphi)$ with $p \in U$ such that $f \circ \varphi ^{-1}$ is smooth at $p$. It is \textbf{smooth} if it is smooth at every $p \in M$.
\end{definition}

\begin{proposition}
  \

  The following are equivalent:

  \begin{enumerate}
    \item The function $f: M \to \mathbb{R}$ is $C^\infty$
    \item The manifold $M$ has an atlas such that for every (smooth) chart $(U, \varphi)$ in the atlas, $f \circ \varphi ^{-1}$ is $C^\infty$
    \item For every (smooth) chart $(V, \psi)$, the function $f \circ \psi ^{-1}$ is smooth.
  \end{enumerate}
\end{proposition}

\begin{customproof}
  \

  ($(2) \implies (1)$) Since the atlas forms an open cover of $M$, for every $p$, there is a chart $(U, \varphi)$ such that $f \circ \varphi ^{-1}$ is smooth, hence $f$ is smooth.

  $(3) \implies (2)$ If it is true for any chart, it is true for the charts on the atlas of $M$.

  $(1) \implies (3)$ Let $f$ be smooth, and let $(V,\psi)$ be some chart. Then, $f \circ \psi ^{-1} \vert_{U \cap V} = (f \circ \varphi ^{-1}) \circ (\varphi \circ\psi ^{-1})$ which is smooth.
\end{customproof}

\begin{definition}[Pullback]
  \

  Let $F: M \to N$ be a map, and $h: N \to \mathbb{R}$. The \textbf{pullback} of $h$ by $F$ is $$F^*h = h \circ F$$
\end{definition}

\subsection{Smooth Maps Between Manifolds}
\

\begin{definition}[Smooth Maps Between Manifolds]
  \

  A map $F: M \to N$ is \textbf{smooth} at $p$ if there is a chart on $p \in U \subseteq M$, $(U, \varphi)$ a chart $F(p) \in V \subseteq N$, $(V, \psi)$ such that $\psi \circ F \circ \varphi ^{-1}$ is smooth at $p$.

  $F$ is \textbf{smooth} if the above is true for all $p \in M$.
\end{definition}

\begin{proposition}
  \

  Let $F: M \to N$ be smooth. For any $p \in M$, any smooth charts $(U, \varphi)$ with $p \in U \subseteq M$, and $(V, \psi)$, with $F(p) \in V \subseteq N$, then $\psi \circ F \circ \varphi ^{-1}$ is a smooth function.
\end{proposition}

\begin{customproof}
  \

  Since $F$ smooth, let $(U', \varphi')$, $(V', \psi')$ be charts such that $\psi' \circ F \circ {\varphi'} ^{-1}$ is smooth. Then, on the intersection of the domains of $\psi', \psi, \varphi, \varphi'$, we have $\psi \circ {\psi'} ^{-1} \circ \psi' \circ F \circ {\varphi'} ^{-1} \circ \varphi' \circ \varphi ^{-1}$ is smooth, since it is the composition of $\psi \circ {\psi'} ^{-1}$, $\psi' \circ f \circ {\varphi'} ^{-1}$ and $\varphi' \circ \varphi ^{-1}$ are all smooth, and the composition of smooth functions is smooth.
\end{customproof}

\begin{proposition}
  \

  The following are equivalent:

  \begin{enumerate}
    \item The map $F: M \to N$ is $C^\infty$
    \item There exist atlases $\mathcal{A}$ in $M$ and $\mathcal{B}$ in $N$ such that for every chart $(U, \varphi)$ in $\mathcal{A}$ and $(V, \psi)$ in $\mathcal{B}$, $\psi \circ F \circ \varphi ^{-1}$ is smooth on their intersection
    \item For every smooth chart $(U, \varphi)$ on $M$ and $(V, \psi)$ on $N$, the map $\psi \circ F \circ \varphi ^{-1}$ is smooth.
  \end{enumerate}
\end{proposition}

\begin{proposition}
  \

  If $F: M \to N$ and $G: N \to P$ are both smooth, $G \circ F$ is smooth.
\end{proposition}

\begin{customproof}
  \

  $F$ being smooth means that for every chart $(U, \varphi)$ and $(V, \psi)$ on $M$ and $N$ respectively, $\psi \circ F \circ \varphi$ is smooth on their intersection.

  $G$ being smooth means that for every chart $(V,\psi)$ and $(W, \sigma)$ on $N$ and $P$ respectively, $\sigma \circ G \circ \psi ^{-1}$ is smooth.

  But then $$\sigma \circ G \circ \psi ^{-1} \circ \psi \circ F \circ \varphi ^{-1}$$is smooth since it is the composition of smooth functions, so $$\sigma \circ G \circ F \circ \varphi ^{-1}$$ is smooth, fulfilling the criteria for smooth functions on $M$
\end{customproof}

\subsection{Diffeomorphisms}
\

\begin{definition}[Diffeomorphisms]
  \

  If a function $F: M\to N$ is smooth and there exists an inverse function $F ^{-1}: N \to M$ that is also smooth, $F$ and $F ^{-1}$ are \textbf{diffeomorphisms}.
\end{definition}

\begin{proposition}
  \

  If $(U, \varphi)$ is a chart for a manifold $M$, then $\varphi$ is a diffomorphism between $U$ and $\mathbb{R}^m$.
\end{proposition}

\begin{customproof}
  \

  By definition, $\varphi$ is a homeomorphism, so it must be continuous with continuous inverse. Then, for any other chart $(V,\psi)$ on $M$, and the identity chart $(\mathbb{R}^m, \iota)$ on $\mathbb{R}^m$, we have $$\iota \circ \varphi \circ \psi ^{-1} = \varphi \circ \psi ^{-1}$$is smooth by chart compatibility, and $$\psi \circ \varphi ^{-1} \circ \iota = \psi \circ \varphi ^{-1}$$also smooth by chart compatibility, hence $\varphi$ is a diffeomorphism.
\end{customproof}

\begin{proposition}
  \

  If $U \subseteq M$ and $F: U \to \mathbb{R}^m$ is a diffeomorphism, then $(U, F)$ is a smooth chart on $M$
\end{proposition}

\subsection{Smoothness in Terms of Components}
\

\begin{proposition}
  \

  The following are equivalent

  \begin{enumerate}
    \item The map $F: M \to \mathbb{R}^m$ is smooth
    \item The manifold $M$ has an atlas such that for every chart $(U,\varphi)$ in the atlas, $F \circ \varphi ^{-1}$ is a smooth function
    \item For every chart $(U, \varphi)$ on $M$, $F \circ \varphi ^{-1}$ is smooth.
  \end{enumerate}
\end{proposition}

\begin{proposition}
  \

  Let $M$ be a manifold, $F: M \to \mathbb{R}^m$ a vector valued function. $F$ is smooth iff every component $F^k$ is smooth.
\end{proposition}

\begin{customproof}
  \

  The map $F: M \to \mathbb{R}^m$ is smooth iff for every chart $(U, \varphi)$ on $M$, $F \circ \varphi ^{-1}$ is smooth, but since this is a function between two Euclidean spaces, this is only smooth iff each component $F^i \circ \varphi ^{-1}$ is smooth, and this is true iff $F^i: M \to \mathbb{R}$ are smooth.
\end{customproof}

\begin{definition}[Lie Groups]
  \

  A \textbf{Lie group} is a smooth manifold $M$ such that the multiplication map $$\mu: M \times M \to M$$ and the inverse map $$\iota: M \to M$$are both smooth maps.
\end{definition}

\subsection{Partial Derivatives}
\

\begin{definition}[Partial Derivatives]
  \

  If $M^m$ is a manifold, $(U, \varphi)$ a chart, then $\varphi$ induces coordinates on $U$, by $x^i = r^i \circ \varphi$, if $r^i$ is the $i$th coordinate function on $\mathbb{R}^m$, then $x^i$ is a coordinate function on $M$.

  If $f: M \to \mathbb{R}^m$ a function, we then have $$\left.\frac{\partial}{\partial x^i} \right\vert_p f = \frac{\partial f}{\partial x^i}(p) = \frac{\partial (f \circ \varphi ^{-1})}{\partial r^i}(\varphi(p)) = \left.\frac{\partial }{\partial r^i}\right\vert_{\varphi(p)} (f \circ \varphi ^{-1})$$
\end{definition}

\begin{proposition}
  \

  Let $(U, x^1, \dots, x^n)$ be a chart on a manifold. Then, $$\frac{\partial x^i}{\partial x^j} = \delta^i_j$$
\end{proposition}

\begin{customproof}
  \

  $$\frac{\partial x^i}{\partial x^j} = \frac{\partial (x^i \circ \varphi ^{-1})}{\partial r^j} = \frac{\partial r^i}{\partial r^j} = \delta_j^i$$
\end{customproof}

\begin{definition}[Jacobian]
  \

  Let $F: M\to N$ be smooth, $(U, \varphi = x^1, \dots, x^m)$ be a chart on $M$, $(V, \psi = y^1, \dots, y^m)$ be a chart on $N$, and $F^i = y^i \circ F$. Then, the matrix $$[\frac{\partial F^i}{\partial x^j}]$$is called the \textbf{Jacobian matrix} and $$\det[\frac{\partial F^i}{\partial x^j}] = \frac{\partial (F^i, \dots, F^n)}{\partial(x^i, \dots, x^n)}$$is the \textbf{Jacobian determinant}
\end{definition}

\subsection{Inverse Function Theorem}
\

\begin{definition}[Locally Invertible]
  \

  A smooth map $F: M \to N$ is \textbf{locally diffeomorphic} or \textbf{locally invertible} at a point $p$ if for some neighbhorhood $U$, $p \in U$, $F \vert_U$ is a diffeomorphism.
\end{definition}

\begin{theorem}[Inverse Function Theorem]
  \

  Let $F: M \to N$ be smooth and $\dim M = \dim N$, pick some point $p \in M$, let there be charts $(U, x^1, \dots, x^n)$, $(V, y^1, \dots, y^n)$ with $p \in U$, $F(p) \in V$ respectively, let $F^i = y^i \circ F$. Then, $F$ is locally invertible at $p$ iff $\det[\frac{\partial F^i}{\partial x^j}(p)]$ is nonzero.
\end{theorem}

\begin{customproof}[Idea]
  \

  It's analagous to the inverse function theorem between $\mathbb{R}^n$ and $\mathbb{R}^n$, but we add coordinate maps $\psi \circ F \circ \varphi ^{-1}$ in the front and back.
\end{customproof}

\begin{corollary}
  \

  Let $M^m$ be a smooth manifold. The set of smooth functions $\{F^1, \dots, F^n\}$ on a neighborhood $(U, x^1, \dots, x^n)$ with $p \in U$ forms a coordinate system about $p$ iff $\det[\frac{\partial F^i}{\partial x^j}(p)]$ is nonzero.
\end{corollary}

\begin{customproof}
  \

  $\det[\frac{\partial F^i}{\partial x^j}(p)]$ is nonzero iff $F$ is diffeomorphicm iff $(U, F)$ forms a chart (aka a coordinate system).
\end{customproof}

\subsection{Exercises}
\

\begin{exercise}[6.1]
  \

  Consider the structure given by the maximal atlases of the following charts let $M$ be the real line with $(\mathbb{R}, \text{id}_\mathbb{R})$ and $N$ be the real line with $(\mathbb{R}, \varphi: x \mapsto x^3)$.
  \begin{enumerate}
    \item Show these are distinct structures
    \item Show there exists a diffeomorphsim $F: M \to N$
  \end{enumerate}
\end{exercise}

\begin{customproof}
  \

  \begin{enumerate}
    \item To show they are distinct structures, we want to show that the atlases are not compatible. See that $\text{id}_\mathbb{R} \circ \varphi ^{-1} = x^\frac{1}{3}$ which is indeed not a smooth map, hence the atlases are incompatible and they define distinct structures.
    \item The diffeomorphism is $F(x) = x^\frac{1}{3}$. We have $\varphi \circ F\circ \text{id}_ \mathbb{R} ^{-1} = (x^\frac{1}{3})^3= x$ is indeed smooth, and $\text{id}_ \mathbb{R} \circ F ^{-1} \circ \varphi ^{-1} = (x^3)^\frac{1}{3} = x$ also smooth.
  \end{enumerate}
\end{customproof}

\section{Quotients}

\subsection{The Quotient Topology}
\

\begin{definition}[The Quotient Topology]
  \

  Let $S$ be a topological space, $\sim$ be an equivalence relation of points on $S$. Let $S/\sim$ be the \textbf{quotient topology} comprised of equivalence classes of $S$ under $\sim$, with the map $\pi: S \to S/\sim$ by $\pi: x \mapsto [x]_\sim$. A set $U \subseteq S/\sim$ is open iff $\pi ^{-1}(U)$ is open.
\end{definition}

\subsection{Continuity of a Map on a Quotient}
\

\begin{proposition}
  \

  Let $S$ be a topological space, $\sim$ an equivalence relation on $S$, and $f: S \to Y$ a continuous map between topological spaces that is constant on each equivalence class. Then, there is a unique continuous map $\bar{f}: S/\sim \to Y$ such that the following diagram commutes: \[
    \begin{tikzcd}
      S & S/\sim\\
      & Y
      \arrow[from=1-1, to=1-2, "\pi"]
      \arrow[from=1-1, to=2-2, "f"']
      \arrow[from=1-2, to=2-2, "\bar{f}"]
  \end{tikzcd}\]
\end{proposition}

\begin{customproof}
  \

  Let $\bar{f}: [x] \mapsto f(x)$ for any representative $x$. This is well defined since $f$ is constant on each equivalence class.

  The diagram commutes since $f(x) = \bar{f}([x]) = \bar{f}(\pi(x))$

  $\bar{f}$ is a morphism of topological spaces since if $V$ is open in $Y$, then $f ^{-1}(V)$ is open in $S$ by continuity. Then, $\pi(S)$ is open by definition of the quotient topology, and $\bar{f}(\pi(S)) = V$ by commutativity, so $\bar{f}$ takes open images to open preimages, so it is continuous.
\end{customproof}

\subsection{Identification of a Subset to a Point}
\

\begin{definition}[Identifying with a Point]
  \

  If $A$ is a subspace of a topological space $S$, we can define an equivalence relation $\sim$ on $S$ by $x \sim x$ for all $x \in S$, and $x \sim y$ for all $x,y \in A$.

  Then, $S/\sim$ just "shrinks" $A$ to a single point.
\end{definition}

\begin{proposition}
  \

  Let $I = [0,1]$, $I/\sim$ be the quotient space obtained by identifying $0 \sim 1$. Show that $f(x) = e^{2\pi ix}: I/\sim \to S^1$ is a homeomorphism.
\end{proposition}

\begin{customproof}
  \

  It is true that $f(x)$ is continuous, and a bijection, and $I/\sim$ is compact, and $S^1$ Hausdorff (as a subspace of $\mathbb{C}$). Then, since $f$ is a continuous bijection from a compact space to a Hausdorff space, it is a homeomorphism.
\end{customproof}

\subsection{A Necessary Condition for the Hausdorff Quotient}
\

\begin{proposition}
  \

  If the quotient space $S/\sim$ is Hausdorff, the equivalence class $[p]$ of any point $p \in S$ must be closed.
\end{proposition}

\subsection{Open Equivalence Relations}
\

\begin{definition}[Open Equivalence Relations]
  \

  The equivalence relation $\sim$ on a topological space $S$ is \textbf{open} if the projective map $\pi: S \mapsto S/\sim$ is an open map (maps open sets in $S$ to open sets in $S/\sim$)
\end{definition}

\begin{definition}[Graph of an Equivalence Relation]
  \

  Given an equivalence relation $\sim$ on $S$, let $R$ be the subset of $S \times S$, where $$R = \{(x,y) \in S \times S \mid s \sim y\}$$Then $R$ is the \textbf{graph} of the equivalence relation.
\end{definition}

\begin{theorem}
  \

  If $\sim$ is an open equivalence relation, then $S/\sim$ is Hausdorff if and only if the graph of $\sim$ is closed (in the product topology).
\end{theorem}

\begin{corollary}
  \

  A topological space $S$ is closed if and only if the diagonal $\Delta$ is closed in $S \times S$ (in the product topology).
\end{corollary}

\begin{theorem}
  \

  Let $\sim$ be an open equivalence relation on a topological space $S$ with projection $\pi: S \to S/\sim$. If $\mathcal{B} = \{B_\alpha\}$ is a basis for $S$, then $\{\pi(B_\alpha)\}$ is a basis for $S/\sim$.
\end{theorem}

\begin{customproof}
  \

  Let $U$ be open in $S/\sim$. Then, $\pi ^{-1}(U)$ is open in $S$. But then, there exists some collection of basis elements $B_i$ such that $\bigcup B_i = \pi ^{-1}(U)$ but this implies that $\bigcup \pi(B_i) = U$ and each $\pi(B_i)$ is open, so $\pi(B_\alpha)$ forms a basis on $S/\sim$.
\end{customproof}

\begin{corollary}
  \

  If $\sim$ is an open equivalence relation, $S$ second countable, then $S/\sim$ is second countable.
\end{corollary}

\subsection{Real Projective Space}
\

\begin{definition}[Real Projective Space]
  \

  Consider a relation on $\mathbb{R}^{n+1} \setminus \{0\}$ by $$x \sim y \iff y = tx$$for some real number $t$. Then, the \textbf{real projective space} is $\mathbb{R} \mathbb{P}^n = (\mathbb{R}^{n+1}\setminus \{0\})/\sim$. The equivalence class of a point $(p^0, \dots, p^n) \in \mathbb{R}^{n+1} \setminus \{0\}$ is $[a^0, \dots, a^n]$ and called \textbf{homogeneous coordinates}.
\end{definition}

\begin{remark}
  \

  We can think of two points as being equivalent if they lie on the same line through the origin. If we then think of a unit sphere about the origin, this line passes through exactly two points on the sphere as well, so we can also think of $\mathbb{R} \mathbb{P}^n$ as half of a sphere.
\end{remark}

\begin{proposition}
  \

  The equivalence relation in the definition above is an open equivalence relation.
\end{proposition}

\begin{customproof}
  \

  The way we prove open equivalence relation is to show that $\pi ^{-1}(\pi(U))$ is open for any open set $U$. But $\pi ^{-1}(\pi(U)) = tU = \{tu \mid u \in U, t \in \mathbb{R}\}$ and since multiplication by a real number is a homeomorphism (true, look it up), then it is indeed open.
\end{customproof}

\begin{proposition}
  \

  $\mathbb{RP}^n$ is second countable.
\end{proposition}

\begin{proposition}
  \

  $\mathbb{RP}^n$ is Hausdorff.
\end{proposition}

\subsection{The Standard $C^\infty$ Atlas on $\mathbb{RP}^n$}
\

\begin{theorem}
  \

  There exists a smooth atlas for $\mathbb{RP}^n$
\end{theorem}

\begin{customproof}
  \

  Let $$U_i = \{[a^0, \dots, a^n] \mid a^i \neq 0\}$$and $$\varphi_i: U_i \to \mathbb{R}^n$$ by $$[a^0, \dots, a^n] \mapsto (\frac{a^1}{a^i}, \dots, \frac{a^n}{a^i})$$omitting the $i$th coordinate.

  Every $U_i$ is open, consider the intersection of any $U_i \cap U_j = \{[a^0, \dots, a^n] \mid a^i, a^j \neq 0\}$. On this intersection, we have two coordinate systems, set the one belonging to $U_i$ to be $x^k$, and the one belonging to $U_J$ to be $y^k$. Then, our transition function $\varphi_j \circ \varphi_i ^{-1}(x) = (\frac{x^1}{x^j}, \dots, \frac{x^{i-1}}{x^j}, \frac{1}{x^j}, \frac{x^{i+1}}{x^j}, \dots, \frac{x^n}{x^j})$ and this is a smooth function since $x^j \neq 0$.
\end{customproof}

\section{The Tangent Space}

\subsection{The Tangent Space at a Point}
\

\begin{definition}[Derivations on Manifolds]
  \

  Germs on a manifold $M$ at a point $p$, \textbf{germs} are equivalence classes of functions from any neighbhorhood around $p$ to $\mathbb{R}$, with the equivalence relation of agreeing on some neighborhood around $p$.

  A \textbf{derivation at $p$} is a vector space morphism (linear map) $D: C^\infty_p(M) \to \mathbb{R}$ (where $C^\infty_p(M)$ is the set of germs at $p$), that satisfies the Leibniz condition: $$D(fg) = fD(g) + D(f)g$$
\end{definition}

\begin{definition}[Tangent Vector]
  \

  A \textbf{tangent vector} at a point $p$ in a manifold is a derivation at $p$. The space of all tangent vectors at $p \in M$ is denoted $T_pM$
\end{definition}

\begin{definition}
  \

  Let $(U,\varphi)$ be a chart, $p \in U$. @e define $$\left.\frac{\partial}{\partial x^i} \right\vert_p f = \left.\frac{\partial}{\partial r^i} \right\vert_{\varphi(p)} (f \circ \varphi ^{-1})$$
\end{definition}

\subsection{The Differential of a Map}
\

\begin{definition}[Differential of a Map]
  \

  Given a smooth map $F: M \to N$, the \textbf{differential} of $F$ at $p$, denoted $F_{*,p}: T_pM \to T_{F(p)}N$, is given by $$F_{*,p}(v(f)) = v(f \circ F)$$for any $v \in T_pM$, $f: N \to \mathbb{R} \in C^\infty_{F(p)}(N)$.
\end{definition}

\subsection{The Chain Rule}
\

\begin{theorem}[The Chain Rule]
  \

  Let $M, N, P$ be smooth manifolds with smooth maps $F: M \to N$, $G: N \to P$, then for $p\in M$ $$(G \circ F)_{*,p} = G_{*, F(p)} \circ F_{*,p}$$
\end{theorem}

\begin{customproof}
  \

  Let $v \in T_pM$, $f: P \to \mathbb{R}$ smooth at $GF(p)$, then$$(G \circ F)_{*, p}(v)(f) = v(f \circ G \circ F)$$But we also can think about $F_{*,p}(v)$ as a vector in $T_{F(p)}N$, so we get $$(G_{*, F(p)} \circ F_{*, p}(v))(f) = G_{*,F(p)}(F_{*,p}(v))(f) = (F_{*,p}(v))(f \circ G) = v(f \circ G \circ F)$$hence the two are equal.
\end{customproof}

\begin{corollary}
  \

  If $F: M \to N$ is a diffeomorphism of manifolds, $p \in M$, then $F_*: T_pM \to T_{F(p)}N$ is a vector space isomorphism.
\end{corollary}

\begin{customproof}
  \

  Use chain rule with the fact that $F$ has a smooth inverse.
\end{customproof}

\begin{corollary}
  \

  If an open set $U \subseteq \mathbb{R}^n$ is diffeomorphic to another open subset $V \subseteq \mathbb{R}^m$, then $m = n$.
\end{corollary}

\subsection{Bases for the Tangent Space at a Point}
\

\begin{proposition}
  \

  Let $(U, \varphi = x^1, \dots, x^n)$, then $$\varphi_*(\left.\frac{\partial}{\partial x^i} \right\vert_p) = \left.\frac{\partial}{\partial r^i} \right\vert_{\varphi(p)}$$Note $r^i$ are the standard coordinate functions, the functions which pick out the $i$th component of a vector in $\mathbb{R}^n$, and $x^i = r^i \circ \varphi$
\end{proposition}

\begin{customproof}
  \

  Let $f: \mathbb{R}^n \to \mathbb{R}$ be smooth$$\varphi_*(\left.\frac{\partial}{\partial x^i}\right\vert_p)f = \left.\frac{\partial}{\partial x^i}f \circ \varphi\right\vert_{p} = \left.\frac{\partial}{\partial r^i} f \circ \varphi \circ \varphi ^{-1} \right\vert_{\varphi(p)} = \left.\frac{\partial}{\partial r^i}\right\vert_p f$$
\end{customproof}

\begin{proposition}
  \

  If $M$ is a manifold, $(U, \varphi = x^1 \dots x^n)$ is a chart containing $p$, then $$\left.\frac{\partial}{\partial x^1}\right\vert_p, \dots, \left.\frac{\partial}{\partial x^n}\right\vert_p$$forms a basis for $T_pU = T_pM$
\end{proposition}

\begin{customproof}
  \

  $\varphi$ is a diffeomorphism, so $\varphi_*$ is a vector space isomorphism, so since $\left.\frac{\partial}{\partial r^i}\right\vert_{\varphi(p)}$ forms a basis for $\mathbb{R}^n$, then so does $\left.\frac{\partial}{\partial x^i}\right\vert_p$
\end{customproof}

\begin{proposition}
  \

  Suppose $(U, x^i)$ and $(V, y^j)$ are two coordinate charts on a manifold $M$. Then, $$\frac{\partial}{\partial x^j} = \sum_{i} \frac{\partial y^i}{\partial x^j} \frac{\partial}{\partial y^i}$$
\end{proposition}

\begin{customproof}
  \

  Since both $\{\left.\frac{\partial}{\partial x^i}\right\vert_p\}$ and $\{\left.\frac{\partial}{\partial y^i}\right\vert_p\}$ form a basis for $T_pM$, there exists a transition mapping $$\left.\frac{\partial}{\partial x^j}\right\vert_p = \sum_k a^k_j \left.\frac{\partial}{\partial y^k}\right\vert_p$$But we can apply both sides of the equation to $y^i$ to get $$\left.\frac{\partial y^i}{\partial x^j}\right\vert_p = \sum_k \alpha^k_j \delta^k_j = \alpha^i_j$$and substituting, we get$$\frac{\partial}{\partial x^j} = \sum_{i} \frac{\partial y^i}{\partial x^j} \frac{\partial}{\partial y^i}$$
\end{customproof}

\subsection{A Local Expression for the Differential}
\

\begin{proposition}
  \

  Given a smooth map $F: M\to N$ and a point $p \in M$, let $(U, x^i)$ and $(V, y^j)$ be charts where $p \in U$, $F(p) \in V$. Relative to the bases $\{\left.\frac{\partial}{\partial x^i}\right\vert_p\}$ and $\{\left.\frac{\partial}{\partial y^j}\right\vert_{F(p)}\}$, the differential $F_{*, p}$ (which is a linear transformation) can be represented by the matrix $[\frac{\partial F^i}{\partial x^j}(p)]$ where $F^i = y^i \circ F$ the $i$th component of $F$.
\end{proposition}

\begin{customproof}
  \

  Again, since this is a linear transformation of vector spaces, we want to see where the basis vectors get sent. The map can be expressed as: $$F_*(\left.\frac{\partial}{\partial x^j}\right\vert_{p}) = \sum_k a^k_j \left.\frac{\partial}{\partial y^k}\right\vert_{F(p)}$$Applying $y^i$ to both sides, we get $$F_*(\left.\frac{\partial}{\partial x^j}\right\vert_{p})(y^i) = \left.\frac{\partial}{\partial x^j}\right\vert_p (y^i \circ F) = a^j_i$$
\end{customproof}

\subsection{Curves in a Manifold}
\

\begin{remark}
  \

We use $]a,b[ = (a,b)$ to denote an open interval to avoid confusion with ordered pairs.
\end{remark}

\begin{definition}[Smooth Curves]
  \

A \textbf{smooth curve} in a manifold $M$ is a map $\gamma: ]a,b[ \to M$. The \textbf{velocity vector} is given by $\gamma'(t_0) = \gamma_*(\left.\frac{d}{dt}\right\vert_{t_0})$.
\end{definition}

\begin{proposition}[Velocity of a Curve in Local Coordinates]
  \

Let $\gamma: ]a,b[ \to M$ be a smooth curve. Then, $$\gamma'(t) = \sum_{i=1}^n {\gamma^i}'(t)\left.\frac{\partial}{\partial x^i}\right\vert_{\gamma(t)}$$
\end{proposition}

\begin{customproof}
  \

  Once again, since we know the differential is a linear map, we can use the general formula: $$\gamma_*(\left.\frac{d}{dt}\right\vert_t) = \sum_k a_k \left.\frac{\partial}{\partial x^k}\right\vert_{\gamma(t)}$$Applying to coordinate functions $x^i$, we get $$\left.\frac{d}{dt}\right\vert_t x^i \circ \gamma = {\gamma^i}'(t)= a_i$$
\end{customproof}

\begin{proposition}
  \

For any $p \in M$, any $X_p \in T_pM$, there exists a $\varepsilon > 0$ such that there exists some $\gamma: ]-\varepsilon,\varepsilon[ \to M$ a smooth curve with $\gamma(0) = p$ and $\gamma'(0) = X_p$.
\end{proposition}

\begin{customproof}
  \

  Let $(U, \varphi = x^i)$ be a chart $p \in U$, with $\varphi(p) = 0$, then $x^i = r^i \circ \varphi$, and let $X_p = \sum a^i \left.\frac{\partial}{\partial x^i}\right\vert_p$.

Consider $\gamma(0) = \varphi ^{-1}((a^1t, \dots, a^nt))$ for $t \in ]-\varepsilon, \varepsilon[$ with $\varepsilon$ sufficiently small such that $\gamma$ lies entirely within $U$.

  Then, $\gamma(0) = \varphi ^{-1}((0, \dots, 0)) = \varphi ^{-1}(0) = p$, and if we let $\alpha(t) = (a^1t, \dots, a^nt)$, we have$$\gamma'(0) = (\varphi ^{-1})_* \circ \alpha_*(\left.\frac{\partial}{\partial t}\right\vert_0) = (\varphi ^{-1})_*(\sum_i a^i \left.\frac{\partial}{\partial r^i}\right\vert_0) = \sum_i a^i \left.\frac{\partial}{\partial ^i}\right\vert_p = X_p$$
\end{customproof}

\begin{proposition}
  \

Suppose $X_p$ is a tangent vector at a point $p$ of a manifold $M$ and $f \in C^\infty_p(M)$. If $\gamma: ]-\varepsilon, \varepsilon[ \to M$ is a smooth curve starting at $p$ with $c'(0) = X_p$ show that $$X_pf = \left.\frac{d}{dt}\right\vert_0(f \circ \gamma)$$
\end{proposition}

\begin{customproof}
  \

  $$X_pf = \gamma'(0)f= \gamma_* (\left.\frac{d}{dt}\right\vert_0)f = \left.\frac{d}{dt}\right\vert_0 f \circ \gamma$$
\end{customproof}

\subsection{Computing the Differential Using Curves}
\

\begin{proposition}
  \

  Let $F: M\to N$ be a smooth map, $p\in M$ and $X_p \in T_pM$. Let $\gamma$ be a smooth curve and $\gamma(0) = p$, $\gamma'(0) = X_p$. Then, $$F_{*,p}(X_p) = \left.\frac{d}{dt}\right\vert_0 F \circ \gamma$$
\end{proposition}

\begin{customproof}
  \

  $$F_{*,p}(X_p) = F_{*,p}(\gamma'(0)) = F_{*,p}(\left.\frac{d}{dt}\right\vert_0 \gamma) = \left.\frac{d}{dt}\right\vert_0 F \circ \gamma$$
\end{customproof}

\subsection{Immersions and Submersions}
\

\begin{definition}[Immersions and Submersions]
  \

  Let $F: M \to N$ be smooth. $F$ is an \textbf{immersion} at $p$ if $T_pM$ is injective, it is a \textbf{submersion} at $p$ if $T_pM$ is surjective.

  $F$ is an \textbf{immersion} (globally) if $T_pM$ is injective at all $p$, it s a \textbf{submersion} if $T_pM$ is surjective at all $p$.
\end{definition}

\subsection{Rank, and Critical and Regular Points}
\

\begin{definition}[Rank]
  \

  The \textbf{rank} of a linear transformation $L: V \to W$ is the dimension of the image $L(V)$ in $W$.

  The \textbf{rank} of a smooth map $F: M \to N$ at a point $p$, denoted $\text{rk} F(p)$ is the rank of $F_{*,p}$.
\end{definition}

\begin{definition}[Critical/Regular Points/Values]
  \

  A point $p \in M$ is \textbf{critical point} if $F_{*,p}$ is not surjective (not a submersion at $p$). It is a \textbf{regular point} if $F_{*,p}$ is a submersion at $p$.

  A point $q \in N$ is a \textbf{critical value} if any point in the preimage $F ^{-1}(p)$ is critical. It is a \textbf{regular value} if it is not critical
\end{definition}

\begin{remark}
  \

  Any point not in the image of $F$ is regular.
\end{remark}

\begin{proposition}
  \

  For a real-valued function $f: M \to \mathbb{R}$, a point $p \in M$ is a critical point iff relative to some chart $(U, x^i)$ with $p \in U$, all the partial derivatives satisfy $$\frac{\partial f}{\partial x^i}(p) = 0$$
\end{proposition}

\begin{customproof}
  \

  The matrix representing $f_{*,p}$ is represented by $$
  \begin{bmatrix}
    \frac{\partial f}{\partial x^1}(p) & \dots & \frac{\partial f}{\partial x^n}(p)
  \end{bmatrix}$$So this matrix can have either rank $0$ or $1$. It has rank $1$ (and hence is surjective) if any $\frac{\partial f}{\partial x^i}$ is nonzero, hence it is only not surjective (and hence $p$ is critical) if and only if all of them are zero.
\end{customproof}

\section{Submanifolds}

\subsection{Submanifolds}
\

\begin{definition}[Regular Submanifold]
  \

  A subset $S$ of a manifold $M^m$ is a \textbf{regular submanifold} of dimension $k$ if for every $p \in S$, there is some chart $(U, \varphi=x^i)$ with $p \in U$ such that $U \cap S$ is the zero set of $m-k$ coordinate functions, which we can assume are the coordinates $x^{k+1} \dots, x^m$.

  The chart $(U, \varphi = x^i)$ is called the \textbf{adapted chart}, and denote the restriction to $S$ by $\varphi_S: U \cap S \to \mathbb{R}^k$, making $(U \cap S, \varphi_S)$ a chart for $S$.
\end{definition}

\begin{definition}[Codimension]
  \

  If $S^k$ is a regular submanifold of $M^m$, then the \textbf{codimension} of $S$ in $M$ is $m - k$.
\end{definition}

\begin{proposition}
  \

  Let $S \subseteq M^m$ be a regular submanifold, $\{(U, \varphi)\}$ be a collection of compatible charts on $M$ that form an open cover of $S$. Then, $\{(U \cap S), \varphi_S\}$ is an atlas of $S$, and so a regular submanifold is in itself a manifold.

  If $M$ has dimension $m$ and $S$ is locally defined by the vanishing of $n-k$ coordinates, then $\dim S = k$.
\end{proposition}

\begin{customproof}
  \

  These charts are already compatible and form an open cover of $S$, so they are an atlas.
\end{customproof}

\subsection{Level Sets of a Function}
\

\begin{definition}[Level Sets]
  \

  Given a map $F: M \to N$, a \textbf{level set} of $F$ is the subset $F ^{-1}(\{c\}) = \{p \in M \mid F(p) = c\}$ or the preimage of some $c\in N$. $c \in N$ is called the \textbf{level}.

  If $F:M \to \mathbb{R}^m$, we denote $Z(F) = F ^{-1}(0)$ the \textbf{zero set} of $F$.

  If $c$ is a regular value, then $F ^{-1}(c)$ is called a \textbf{regular level set}, and if $Z(F)$ is regular, we call it a \textbf{regular zero set}.
\end{definition}

\begin{lemma}
  \

  Let $g: M\to \mathbb{R}$ be a smooth function. A regular level set $g ^{-1}(c)$ of level $x$ of the function $g$ is the regular zero set $f ^{-1}(0)$ of the function $f = g - c$
\end{lemma}

\begin{customproof}
  \

  For any $p \in M$, $$g(p) \iff f(p) = g(p) - c = 0$$so $g ^{-1}(c) = f ^{-1}(0)$. Because $f_{*,p} = g_{*,p}$ (since constants fall out of the differential), if $g ^{-1}(c)$ is a regular level set, then $f ^{-1}(0)$ is a regular zero set.
\end{customproof}

\begin{theorem}
  \

  Let $g: M \to \mathbb{R}$ be a smooth function of the manifold $M$. Then, a nonempty regular level set $S = g ^{-1}(c)$ is a regular submanifold of $M$ of codimension $1$.
\end{theorem}

\begin{customproof}
  \

  $S = g ^{-1}(c) = f ^{-1}(0)$ by the prior lemma, where $f = g -c$. Then, since $S$ is a regular zero set, we know that $0$ is a regular value, and all $p \in S$ is a regular value, meaning that $f_{*,p}$ is surjective for all $p \in S$, so for any $p$, chart $(U, x^i)$ with $p \in U$, for some $i$, we have $$f_{*, p}\left(\left.\frac{\partial}{\partial x^i}\right\vert_p \right) = \frac{\partial f}{\partial x^i}(p) \neq 0$$Assume WLOG that $f_{*,p}(\left.\frac{\partial}{\partial x^1}\right\vert_p) \neq 0$. Then, consider the map $\varphi: U \to \mathbb{R}^n$, with $\varphi = (f, x^2, \dots, x^n)$ replacing the $x^1$ component with $f$. Then, the Jacobian is $$
  \begin{bmatrix}
    \frac{\partial f}{\partial x^1} & \dots & \frac{\partial f}{\partial x^n}\\
    0 & \dots & 0 \\
    \vdots & \ddots & \dots \\
    0 & \dots & 1
  \end{bmatrix}$$with $1$s on the diagonal. Then, since we know no row or column is all $0$s, we know the determinant is nonzero, so by the inverse function theorem, there is a neighborhood $U_p$ of $p$ where $(f, x^2, \dots, x^n)$ forms a coordiante system.

  Then, the level set $U_p \cap S$ is defined by setting $f = 0$, making $(U_p, (f, x^2, \dots, x^n))$ an adapted chart, so $S$ is a regular submanifold of codimension $1$.
\end{customproof}

\subsection{Regular Level Set Theorem}
\

\begin{theorem}[Regular Level Set Theorem]
  \

  Let $F: M^m \to N^n$ be a smooth map. A nonempty regular level set $F ^{-1}(c) \subseteq M$ is a regular submanifold of dimension $m - n$.
\end{theorem}

\begin{customproof}
  \

  Let some $p \in M$, $(U, x^i)$ be a coordinate map containing $p$, and let $(V, y^j)$ be a coordinate map where $F(p)$ is centered, so it maps $F(p)$ to $0$. Let $F^i = y^i \circ F$. Then, the pushforward $F_{*,p}$ is defined by the matrix $$F_{*,p} =
  \begin{bmatrix}
    \frac{\partial F^1}{\partial x^1}(p)& \frac{\partial F^1}{\partial x^2}(p)& \dots& \frac{\partial F^1}{\partial x^m}(p)\\
    \vdots & \vdots & \ddots & \vdots\\
    \frac{\partial F^n}{\partial x^1}(p) & \frac{\partial F^n}{\partial x^2}(p) & \dots & \frac{\partial F^n}{\partial x^m}(p)
  \end{bmatrix}$$Note that if a regular level set $F ^{-1}(c)$ is nonempty, this implies that $F$ is a submersion at $c$, so $m \geq n$, so the Jacobian above at least has rank $n$, and assume WLOG the first $n \times n$ block has nonzero determinant. Then, by the inverse function theorem, $(U', \varphi = F^1, \dots, F^n, x^{n+1}, \dots x^m)$ forms a smooth chart, and at $p$, all $F^i(p) = 0$, so $F ^{-1}(c) = \varphi ^{-1}(0)$, so this is an adapted chart, and hence $F ^{-1}(p)$ is a regular submanifold.
\end{customproof}

\begin{lemma}
  \

  Let $F: M^m \to \mathbb{R}^n$ be a smooth map, let $S$ be the level set $F ^{-1}(0)$. Then, if relative to some coordinate chart about $p\in S$, the determinant of the Jacobian, denoted $\frac{\partial (F^1, \dots, F^m)}{\partial (x^1, \dots, x^m)}(p)$ is nonzero, then in some neighbhorhood of $p$, one may replace some $x^{j_i}, \dots, x^{j_m}$ coordinate maps in the chart with $F^1, \dots, F^m$ to obtain an adapted chart for $M$ relative to $S$.
\end{lemma}

\section{Categories and Functors}

\subsection{Categories}
\

\begin{example}
  \

  The collection of smooth manifolds with smooth maps as morphisms forms a category called the \textbf{smooth category}.
\end{example}

\begin{example}
  \

  A tuple $(M,p)$ where $M$ is a smooth manifold, and $p \in M$ is some identified point, is called a \textbf{pointed manifold}. The set of all such pointed manifolds, $(M,p)$, $(N,q)$, and morphisms being point-preserving smooth maps ($F(p) = q$), forms a category called \textbf{the category of pointed manifolds}
\end{example}

\subsection{Functors}
\

\begin{definition}[Functors]
  \

  A \textbf{covariant functor} $F: C \to D$ between categories $C$ and $D$ is a mapping which sends objects $c \in C$ to objects $F(c) \in D$, and morphisms $f \in C(x,y)$ to $Ff \in D(F(x), F(y))$, with the composition property $$F(f \circ g) = Ff \circ Fg$$and the property that $F(1_x) = 1_{F(x)}$

  A \textbf{contravariant functor} reverses the direction of morphisms, so $f \in C(x,y)$ gets sent to $Ff \in D(F(y), F(x))$, with the composition property $$F(f \circ g) = Fg \circ Ff$$and the same identity property.
\end{definition}

\begin{proposition}
  \

  Let $F: C \to D$ be a functor between categories. If $f \in C(x,y)$ is an ismorphism, then so is $Ff$.
\end{proposition}

\begin{customproof}
  \

  An isomorphism is a morphism $f: x \to y$ with an inverse $f ^{-1}$ such that $f \circ f ^{-1} = 1_y$ and $f ^{-1} \circ f = 1_x$. Then, by functoriality, $$1_{F(y)} = F(1_y) = F(f \circ f ^{-1}) = Ff \circ F ^{-1}$$and same for the other way, hence functors preserve isomorphisms.
\end{customproof}

\subsection{The Dual Functor and the Multicovector Functor}
\

\begin{definition}[Dual Linear Maps]
  \

  For every linear map $L: V \to W$ of vector spaces, there exists a \textbf{dual map} $L^*: W^* \to V^*$ by $$L^*(\alpha) = \alpha \circ L$$ for all $\alpha \in W^*$
\end{definition}

\begin{proposition}
  \

  Let $V, W, S$ be vector spaces.
  \begin{enumerate}
    \item If $1_V: V \to V$ is the identity map on $V$, show that $1_V^*$ is the identity map on $V^*$
    \item If $f: V \to W$ and $g: W \to S$ are linear maps, then $(g \circ f)^* = f^* \circ g^*$
  \end{enumerate}
\end{proposition}

\begin{customproof}
  \

  \begin{enumerate}
    \item The dual identity map takes the linear functional $\alpha$ and sends it to the functional$$1_V^*(\alpha) = \alpha \circ 1_V$$but these are both the same linear functional, we can see by how they act on vectors $v \in V$: $$(\alpha \circ 1_V)(v) = \alpha(1_V(v)) = \alpha(v)$$Hence $1_V^*$ takes linear functionals on $V$ to themselves, hence it is the identity.
    \item Let $\alpha \in S^*$, then $$(g \circ f)^*(\alpha) = \alpha \circ g \circ f = (\alpha \circ g) \circ f = f^*(\alpha \circ g) = f^*(g^*(\alpha)) = f^* \circ g^*(\alpha)$$
  \end{enumerate}
\end{customproof}

\begin{definition}[Pullbacks]
  \

  Fix a positive integer $k$. For any linear map $L: V \to W$ of vector spaces, the \textbf{pullback} $L^*: A_k(W) \to A_k(V)$ acting on multicovectors, is given by $$(L^*f)(v_1, \dots, v_k) = f(L(v_1), \dots, L(v_k))$$for $f \in A_k(W)$.

  The pullback is just the generalization of the dual map to the space of alternating functionals
\end{definition}

\begin{proposition}
  \

  \begin{enumerate}
    \item If $1_V: V \to V$ is the identity map on $V$, then $1_V^* = 1_{A_k(V)}$
    \item If $K: U \to V$ and $L: V \to W$ are linear maps, $$(L \circ K)^* = K^* \circ L^*$$
  \end{enumerate}
\end{proposition}

\begin{customproof}
  \

  \begin{enumerate}
    \item $$1_V^*(f) = f(1_V(v_1), \dots, 1_V(v_n)) = f(v_1, \dots, v_n) = f$$hence this is the identity map.
    \item $$(L \circ K)^*(f) = f(L \circ K(v_1), \dots, L \circ K(v_n)) = L^*(f(K(v_1), \dots, K(v_n))) = (K^* \circ L^*)(f)$$
  \end{enumerate}
\end{customproof}

\section{The Rank of a Smooth Map}

\subsection{Constant Rank Theorem}
\

\begin{remark}
  \

  The regular level set theorem gives us a sufficient but not necessary condition for verifying sets are submanifolds, by saying that the preimage of regular values (or alternatively, the pushforward of $F: M \to N$ for some $N$ at every point in the submanifold $S \subseteq M$ is surjective, and therefore, of maximal rank) are submanifolds. This is because we pass to IFT, which allows us to construct a diffeomorphism (which is basically an adapted chart) if we have an $n \times n$ matrix, which we do, its the submatrix of the Jacobian of $F$. But, we can actually drastically strengthen the condition with the following.
\end{remark}

\begin{theorem}[Constant Rank Theorem on Euclidean Spaces]
  \

  If $U \subset \mathbb{R}^n$ and $f: U \to \mathbb{R}^m$ has constant rank $k$ in a neighborhood of a point $p \in U$, then there are diffeomorphisms $G$ sending $p$ to the origin in $\mathbb{R}^n$ and $F$ sending $f(p)$ to the origin in $\mathbb{R}^m$ such that $$(F \circ f \circ G ^{-1})(x^1, \dots, x^n) = (x^1, \dots, x^k, 0, \dots, 0)$$
\end{theorem}

\begin{theorem}[Constant Rank Theorem]
  \

  Let $M^m, N^n$ be smooth manifolds. Let $f: M \to N$ have constant rank $k$ in a neighborhood of a point $p \in M$. Then, there are charts $(U, \varphi)$ centered at $p$ and $(V,\psi)$ centered at $f(p)$ such that for $(r^1, \dots, r^n)$ in $\varphi(U)$, $$(\psi \circ f \circ \varphi ^{-1})(r^1, \dots, r^n) = (r^1, \dots, r^k, 0, \dots, 0)$$
\end{theorem}

\begin{customproof}
  \

  Choose charts $(\bar{U}, \bar{\varphi})$ about $p$ in $N$ and $(\bar{V}, \bar{\psi})$ about $f(p)$ in $M$. Then, $\bar{\psi} \circ f \circ \bar{\varphi} ^{-1}$ is a map between open subsets of Euclidean spaces with constant rank $k$ (since $\bar{\psi}$, $\bar{\varphi}$ are diffeomorphisms, they are rank-preserving). Then, by the Constant Rank Theorem on Euclidean Spaces, there are diffeomorphisms $F, G$ such that $F \circ \bar{\psi} \circ f \circ \bar{\varphi} ^{-1} \circ G ^{-1}(r^1, \dots, r^n) = (r^1, \dots, r^k, 0,\dots, 0)$. Taking $\varphi = G \circ \bar{\varphi}$ and $\psi = F \circ \bar{\psi}$ yields our answer.
\end{customproof}

\begin{theorem}[Constant Rank Level Set Theorem]
  \

  Let $f: M \to N$ be a smooth map between manifolds, $c \in N$. If $f$ has constant rank $k$ in a neighborhood of the level set $f ^{-1}(c)$ in $M$, then $f ^{-1}(c)$ is a regular submanifold of $M$ with codimension $k$.
\end{theorem}

\begin{customproof}
  \

  By the prior theorem, there is some coordinate chart $(U,\varphi = x^1, \dots, x^m)$ centered on $p \in M$ and $(V, \psi = y^1,\dots, y^n)$ centered on $f(p) \in N$ such that $(\psi \circ f \circ \varphi ^{-1})(r^1, \dots, r^n) = (r^1, \dots, r^k, 0, \dots, 0)$ and so $(\psi \circ f \circ \varphi ^{-1}) ^{-1}(0)$ is a level set, since it is defined by the vanishing of the coordinate functions $r^1, \dots, r^k$ (note this is a level set in $\mathbb{R}^m$, and we want to "lift this up" to $M$)

  But then, since $\varphi(c) = 0$ (since $\varphi$ centered on $c$), and each $r^i$ is a diffeomorphism, we can define $x^i = r^i \circ \varphi$ a set of coordinate function, and we find that $$\varphi( f ^{-1}(c)) = \varphi( f ^{-1}(\psi ^{-1}(0))) = (\psi \circ f \circ \varphi ^{-1}) ^{-1}(0)$$is defined by the vanishing of each coordinate function $x^1, \dots, x^k$, so $f ^{-1}(c)$ is also a regular submanifold of $M$ of codimension $k$.
\end{customproof}

\subsection{The Immersion and Submersion Theorems}
\

\begin{proposition}
  \

  Let $M, N$ be manifolds of dimensions $m, n$ respectively. If a smooth map $f: M \to N$ is an immersion at $p$, it has constant rank $m$ in a neighborhood of $p$.

  If $f$ is a submersion, it has constant rank $n$ in a neighborhood of $p$.
\end{proposition}

\begin{customproof}
  \

  $f$ doesn't have maximal rank $k$ at $p$ if every $k \times k$ minor of the jacobian $[\frac{\partial f^i}{\partial x^j}(p)]$ vanishes. But this is the union of the zero sets of some continuous functions, and these are closed sets, hence the points where $k$ is not maximal form a closed set, so if $k$ is maximal at $p$, it is maximal at a neighborhood in $p$.

  If a map is a submersion at a point, we are saying it is of maximal rank $n$, and same for immersions, hence the proposition is true.
\end{customproof}

\begin{theorem}
  \

  \begin{enumerate}
    \item (Immersion Theorem) Suppose $f: M \to N$ is an immersion at $p \in M$. Then, there are charts $(U, \varphi)$ centered at $p \in M$ and $(V,\psi)$ centered at $f(p) \in N$ such that in a neighborhood of $\varphi(p)$, $$(\psi \circ f \circ \varphi ^{-1})(r^1, \dots, r^n) = (r^1, \dots, r^n, 0, \dots, 0)$$
    \item (Submersion Theorem) Suppose $f: M \to N$ is a submersion at $p \in M$. Then, there are charts $(U, \varphi)$ centered at $p \in M$ and $(V,\psi)$ centered at $f(p) \in N$ such that in a neighborhood of $\varphi(p)$, $$(\psi \circ f \circ \varphi ^{-1})(r^1, \dots, r^m, r^{m+1}, \dots, r^n) = (r^1, \dots, r^m)$$
  \end{enumerate}
\end{theorem}

\begin{customproof}
  \

  Just combine the prior two proofs.
\end{customproof}

\begin{corollary}
  \

  A submersion $f: M \to N$ is an open map.
\end{corollary}

\begin{customproof}
  \

  Projections are open maps (by a result in topology), and locally, a submersion is a projection.
\end{customproof}

\subsection{Images of Smooth Maps}
\

\begin{definition}[Embeddings]
  \

  A smooth map $f: M \to N$ is an \textbf{embedding} if both the follwing are true
  \begin{enumerate}
    \item it is a one to one immersion
    \item the image $f(M)$ with the subspace topology is homeomorphic to $M$ under $f$ (meaning topologically, $f(M)$ with the subspace topology from $N$ is the same as $f(M)$ with the topology "open sets $V \subseteq f(M)$ are open if $f(U) = V$ for some open set $U \subseteq M$)
  \end{enumerate}
\end{definition}

\begin{theorem}
  \

  If $f: M \to N$ is an embedding, then $f(M)$ is a regular submanifold of $M$
\end{theorem}

\begin{customproof}
  \

  By the immersion theorem, about ever $p \in M$, there are local coordinates $(U, x^i)$ near $p$ and $(V, y^j)$ near $f(p)$ such that $f(x^1, \dots, x^n) = (x^1, \dots, x^n, 0,\dots, 0)$, so $f(U)$ is defined as the vanishing of $y^{m+1}, \dots, y^n$. But, since $V \cap f(M)$ may be larger than $f(U)$, we aren't done yet.

  Since the topology on $f(M)$ is exactly the subspace topology, there is an open set $V' \subseteq N$ such that $V' \cap f(N) = f(U)$, and so $V \cap V' \cap f(N) = f(U)$, and so $(V \cap V', y^j)$ is now our new adapted chart for $f(N)$. By repeating this for all $f(p)$, we can assemble an open cover of adapted charts. (I actually have no clue why we have to do this last step but okay ig)
\end{customproof}

\begin{theorem}
  \

  If $N$ is a regular submanifold of $M$, the inclusion map $\iota: N \to M$ is an embedding.
\end{theorem}

\begin{customproof}
  \

  Since a regular submanifold has the subspace topology, and so does $\iota(N)$ by definition, $\iota: N \to \iota(N)$ is a homeomorphism. In addition, let $(V, y^1, \dots, y^n, y^{n+1}, \dots, y^m)$ be an adapted chart, so $V \cap N$ is the zero set of $y^{n+1}, \dots, y^m$. Then, this extends to a chart on $M$ through $\iota$ by $$\iota: (y^1,\dots, y^n) \mapsto (y^1, \dots, y^n, 0, \dots, 0)$$and so $\iota$ is an immersion.
\end{customproof}

\subsection{Smooth Maps into a Submanifold}
\

\begin{theorem}
  \

  Suppose $f: M\to N$ is smooth, and $f(M) \subseteq S \subseteq N$. If $S$ is a regular submanifold of $N$, then $\bar{f}: M \to S$ is smooth.
\end{theorem}

\begin{customproof}
  \

  Let $(V, \varphi= y^i)$ be an adapted chart for $S$. Then, since $f$ smooth, for all $q \in M$, $(\varphi \circ f)(q) = (y^1 \circ f(q), \dots, y^s \circ f(q), 0, \dots, 0)$ and each $y^i \circ f(q)$ is smooth (by definition of $f$ being a smooth map). But then each component function of the restriction $(\varphi \circ f) \vert_S(p) = (y^1 \circ f(p), \dots, y^s \circ f(p))$ has components which are all smooth as well, so $(\varphi \circ f) \vert_S = \bar{f}$ is smooth as well.
\end{customproof}

\section{The Tangent Bundle}
\

\subsection{The Topology of the Tangent Bundle}
\

\begin{definition}[Tangent Bundles]
  \

  The \textbf{tangent bundle} of a smooth manifold $M$ is the union of all tangent spaces of $M$: $$TM = \bigcup_{p \in M} T_pM = \bigsqcup_{p \in M} T_pM$$We commonly denote a member of the tangent bundle as $(p,v)$ where $p \in M$ and $v \in T_pM$. There is then a natural (doesn't depend on choice of atlas or local coordinates) map $$\pi: TM \to M$$$$\pi: (p, v) \mapsto p$$

  Now we want to give a topology to $TM$. For any $v \in T_pM$, $v = \sum_{i =1}^m c^i \left.\frac{\partial}{\partial x^i}\right\vert_p$, and $c^i(v)$, the function which picks out the $i$th coefficient of $v$, is a smooth function. Then, for charts $(U, \varphi = x^1, \dots, x^m)$ containing $p$, let $$\bar{\varphi}: TM \to \mathbb{R}^{2m}$$$$\bar{\varphi}: (p, v) \mapsto (x^1(p), \dots, x^m(p), c^1(v), \dots, c^m(p))$$And a set $V \subseteq TM$ is open iff $\bar{\varphi}(V)$ is open in $\mathbb{R}^{2m}$.
\end{definition}

\begin{remark}
  \

  Coordinate open sets are open sets which live entirely inside a chart.
\end{remark}

\begin{lemma}
  \

  \begin{enumerate}
    \item For any manifold $M$, the set $TM$ is the union of all $A$ open in $TM$
    \item Let $U, V$ be coordinate open sets in $M$. If $A$ open in $TU$ and $B$ open in $TV$, then $A \cap B$ is open in $T(U \cap V)$
  \end{enumerate}
\end{lemma}

\begin{customproof}
  \

  \begin{enumerate}
    \item Let $\{U_\alpha, \varphi_\alpha\}$ be an atlas on $M$, consider just the sets $A_\alpha = \{(p,v) \mid p \in U_\alpha, v \in T_pM\}$ are open, since $T_pM$ is an open set and $U_\alpha$ is open, and these form an open cover of $TM$.
    \item If $A$ is open in $TU$, then $A \cap (U \cap V)$ is open in $T(U \cap V)$ (by the subspace topology). If $B$ is open in $TV$, then $B \cap (U \cap V)$ is open in $T(U \cap V)$. Then, the intersection of two open sets $A \cap B$ is open in $T(U \cap V)$.
  \end{enumerate}

  It follows that $\mathcal{B} = \{A \subseteq TM \mid \pi(A) \text{ is a coordiante open set on $M$}, A \text{ open in TM}\}$ (any open set of $TM$ whose projection to $M$ is a coordinate open set) forms an adequate basis for $TM$
\end{customproof}

\begin{lemma}
  \

  A manifold $M$ has a countable basis consisting of coordinate open sets.
\end{lemma}

\begin{customproof}
  \

  We know that a manifold $M$ has a countable basis, and it has an open cover of coordinate open sets, but we want to merge the two. Let $U_\alpha$ be a coordinate open set, $p \in U_\alpha$, and choose some $B_{p, \alpha} \in \mathcal{B}$ such that $p \in B_{p,\alpha} \subset U_\alpha$. The collection of all such $\{B_{p, \alpha}\}$ is a subset of $\mathcal{B}$ and so is countable, and for any open set $U \in M$, $p \in U$, there is a coordinate open set $U_\alpha$ such that $p \in U_\alpha \subset U$, so $p \in B_{p,\alpha} \subset U_\alpha \subset U$, so $\{B_{p,\alpha}\}$ is a basis of coordinate open sets.
\end{customproof}

\begin{proposition}
  \

  The tangent bundle $TM$ of a manifold $M$ is second countable.
\end{proposition}

\begin{customproof}
  \

  If $\{B_{p,\alpha}\}$ is a coordinate open set basis for $M$, $\mathcal{A}$ is a countable basis for $\mathbb{R}^m$, and $\mathcal{B}$, the basis defined above for $TM$, we have: $$\mathcal{B} = \{B_{p,\alpha}\} \times \mathcal{A}$$and the product of two countable sets is countable, hence it is a countable basis.
\end{customproof}

\begin{proposition}
  \

  The tangent bundle $TM$ of a manifold $M$ is Hausdorff.
\end{proposition}

\subsection{The Manifold Structure on the Tangent Bundle}
\

\begin{theorem}
  \

  For a manifold $M$, $TM$ is a smooth manifold with atlas $\{(TU_\alpha, \bar{\varphi}_\alpha)\}$.
\end{theorem}

\begin{customproof}
  \

  To note, $TU_\alpha = \bigsqcup_{p \in U_\alpha} T_pM$ and $\bar{\varphi}_\alpha: (p, v) \mapsto (x^1(p), \dots, x^m(p), c^1(v), \dots, c^m(p))$

  Since $\{U_\alpha\}$ forms an open cover of $M$, $\{TU_\alpha\}$ forms an open cover of $TM$.

  We wnat to show that for $TU_\alpha \cap TU_\beta$, $\bar{\varphi}_\alpha$ and $\bar{\varphi}_\beta$ are compatible. We just show first that $\bar{\varphi}_\beta \circ \bar{\varphi}_\alpha ^{-1}$ is smooth. Since they are coordinate maps, the $x^i$ terms are smooth. For any $v \in TU_\alpha \cap TU_\beta$, $$v = \sum_j a^j \left.\frac{\partial}{\partial x^j}\right\vert_p = \sum_i b^i \left.\frac{\partial}{\partial y^i}\right\vert_p$$Applying both sides to $x^k$, we find $$a^k = \sum_i b^i \frac{\partial x^k}{\partial y^i}$$and since each $$b^i = \sum_j a^j \frac{\partial(\varphi_\beta \circ \varphi_\alpha ^{-1})^i}{\partial r^j}(\varphi_\alpha(p))$$ is smooth, then $\bar{\varphi}_\beta \circ \bar{\varphi}_\alpha = (\dots, b^1, \dots, b^m)$ is smooth and so these are compatible charts.
\end{customproof}

\subsection{Vector Bundles}
\

\begin{definition}[Fibers]
  \

  Given any $\pi: E \to M$, the inverse image $\pi ^{-1}(p)$ is the \textbf{fiber at $p$}, denoted $E_p$. For any maps $\pi: E \to M$ and $\pi': E' \to M$, a map $\varphi: E \to E'$ is \textbf{fiber preserving} if $\varphi(E_p) \subseteq E'_p$ for all $p \in M$.
\end{definition}

\begin{definition}[Locally Trivial Maps]
  \

  A surjective smooth map $\pi: E \to M$ is \textbf{locally trivial of rank $r$} if
  \begin{enumerate}
    \item each fiber $E_p$ has an $r$-dimensional vector space structure
    \item for each $p \in M$, there is an open neighborhood $U$ of $p$, and a fiber preserving diffeomorphism $\varphi: \pi ^{-1}(U) \to U \times \mathbb{R}^r$ such that for every $q \in U$, the restriction $$\varphi \vert_{\pi ^{-1}(q)}: \pi ^{-1}(q) \to \{q\} \times \mathbb{R}^r$$is a vector space isomorphism.
  \end{enumerate}

  Such an open set $U$ as defined above is called a \textbf{trivializing open set} for $E$, and $\varphi$ is a trivialization of $E$ over $U$, the collection $\{(U, \varphi)\}$ where $\{U\}$ forms an open cover of $M$ is called a \textbf{local trivialization} for $E$, and $\{U\}$ is called a \textbf{trivializing open cover} of $M$ for $E$.
\end{definition}

\begin{definition}[Vector Bundles]
  \

  A \textbf{smooth vector bundle of rank $r$} is a triple $(E, M, \pi)$ consisting of manifolds $E, M$, and $\pi: E \to M$ a surjective locally trivial map of rank $r$. The manifold $E$ is the \textbf{total space} of the vector bundle, and $M$ is the \textbf{base space}, and we say that $E$ is a \textbf{vector bundle} over $M$.
\end{definition}

\begin{example}
  \

  The tangent bundle of a manifold $M$, $(TM, M, \pi)$ is a vector bundle of rank $m$.
\end{example}

\begin{example}[Product Bundle]
  \

  Given a manifold $M$, let $\pi: M \times \mathbb{R}^r \to M$ be a projection onto $M$. Then, $M \times \mathbb{R}^r \to M$ is a vector bundle of rank $r$, called the \textbf{product bundle} of rank $r$ over $M$.

  It has vector space structure by $b(p, u) + (p,v) = (p, bu+v)$. We can also construct a trivialization by taking the open set of the entire manifold itself, $M$, and doing $\varphi: M \times \mathbb{R}^r \to M \times \mathbb{R}^r$ as the identity map. Hence, this is locally trivial of rank $r$.
\end{example}

\begin{definition}[Bundle Maps]
  \

  Let $\pi_E: E \to M$ and $\pi_F: F \to N$ be two vector bundles, possibly of different ranks. A \textbf{bundle map} from $E$ to $F$ is a pair of maps $(f, \bar{f})$, with $f: M \to N$, $\bar{f}: E \to F$, such that \[
    \begin{tikzcd}
      E & F\\
      M & N
      \arrow[from=1-1,to=1-2, "\bar{f}"]
      \arrow[from=1-1, to=2-1, "\pi_E"]
      \arrow[from=1-2, to=2-2, "\pi_F"]
      \arrow[from=2-1, to=2-2, "f"]
  \end{tikzcd}\]commutes, and $\bar{f}$ is linear on fibers, so $\bar{f}_p: E_p \to F_p$ is a linear transformation.
\end{definition}

\subsection{Smooth Sections}
\

\begin{definition}[Sections]
  \

  A \textbf{section} of a vector bundle $\pi: E \to M$ is a map $s: M \to E$ such that $\pi \circ s = 1_M$, the identity map on $M$. Basically, $s$ maps a point in $p$ to a point in the fiber $E_p$.

  If $s$ is also a smooth map between manifolds, we say it is a \textbf{smooth section}.
\end{definition}

\begin{definition}[Vector Fields]
  \

  A \textbf{vector field} $X$ on a manifold $M$ is a function that assigns a tangent vector $X_p \in T_pM$ to each point $p \in M$. A vector field is simply a section of the tangent bundle $TM$.

  A vector field is smooth if it is a smooth section.
\end{definition}

\begin{proposition}
  \

  Let $s,t$ be smooth sections of a smooth vector bundle $\pi: E\to M$, and let $f$ be a smooth real-valued function on $M$. Then:

  \begin{enumerate}
    \item The sum $s + t: M \to E$ defined as $$(s+t)(p) = s(p) + t(p)$$is a smooth section.
    \item The product $fs: M \to E$ defined by $(fs)(p) = f(p)s(p)$ is a smooth section
  \end{enumerate}
\end{proposition}

\begin{customproof}
  \

  \begin{enumerate}
    \item We want to show that as a map between manifolds, $s+t$ is smooth. Since the sum of two points in a fiber still lies in the fiber, $s+t$ is still a section. Since $\pi$ is a smooth vector bundle, at every $p\in M$, for some $V$ neighborhood of $p$, it has a trivialization $\varphi: \pi ^{-1}(V) \to V \times \mathbb{R}^r$. Then, suppose $$(\varphi \circ s)(q) = (q, a^1(q), \dots, q^r(q))$$$$(\varphi \circ t)(q) = (q, b^1(q), \dots, b^r(q))$$where each $a^i(q), b^i(q)$ takes the vector that $s, t$ assigns to $q$, and gets the $i$th component. Since $s,t$ are smooth, so are $a^i, b^i$. Since $\varphi$ is linear on fibers, $$(\varphi \circ (s+t))(q) = (q, a^1(q) + b^1(q), \dots, a^r(q) + b^r(q))$$Which proves that $s+t$ is a smooth map on $V$, and since this works for any $p \in M$, it is smooth.
    \item Same argument.
  \end{enumerate}
\end{customproof}

\begin{corollary}
  \

  The set of all smooth sections/vector fields over $E$, $\Gamma(E)$, is a module over the ring of smooth functions on $M$.
\end{corollary}

\subsection{Smooth Frames}
\

\begin{definition}[Frames]
  \

  A \textbf{frame} for a vector bundle $\pi: E \to M$ is a collection of sections $s_1, \dots, s_r$ such that at each point $p \in U$, $s_1(p), \dots, s_r(p)$ form a basis for the fiber $E_p$. A frame is smooth if each $s_i$ is a smooth section.

  A frame for the tangent bundle $TM \to M$ is called a \textbf{frame on $M$}
\end{definition}

\begin{example}
  \

  The vector fields $\frac{\partial}{\partial x}$, $\frac{\partial}{\partial y}$, and $\frac{\partial}{\partial z}$ form a smooth frame on $\mathbb{R}^3$. This is because for fixed $p \in \mathbb{R}^3$, they give $\left. \frac{\partial}{\partial x}\right\vert_p$, $\left. \frac{\partial}{\partial y}\right\vert_p$, $\left. \frac{\partial}{\partial z}\right\vert_p$, which form a basis for $T_p \mathbb{R}^3$.
\end{example}

\begin{lemma}
  \

  Let $\varphi: E\vert_U \to U \times \mathbb{R}^r$ be a trivialization over an open set $U$ of a smooth vector bundle $E \to M$, and $t_1, \dots, t_r$ be a frame over $U$ of the trivialization, meaning $t_i(p) = \varphi ^{-1}(p, e_i)$, where $e_i$ is the $i$th unit vector in $\mathbb{R}^r$. Then, a section $s = \sum b^it_i$ of $E$ over $U$ is smooth iff its coefficients $b^i$ are smooth functions.
\end{lemma}

\begin{customproof}
  \

  $(\impliedby)$ By the prior proposition, if each $b^i$ is a smooth real valued function, the linear combination $\sum b^it_i$ defines a smooth section.

  $(\implies)$ Let $s = \sum b^it_i$ be a smooth section. Then, $\varphi \circ s$ is smooth, so $$(\varphi \circ s)(p) = \sum b^i(p)\varphi(t_i(p)) = \sum b^i(p)(p, e_i) = (p, \sum b^i(p)e_i)$$so each $b^i(p)$ are just the $i$th coefficients of the vector that $s$ maps to in $E$, where the coordinates are determined by the trivialization $\varphi$. Since $\varphi \circ s$ is smooth, so must all the $b^i$ be.
\end{customproof}

\begin{proposition}
  \

  Let $\pi: E \to M$ be a smooth vector bundle, $U$ an open subset of $M$. Suppose $s_1, \dots, s_r$ is a smooth frame for $E$ over $U$. Then, a section $s = \sum c^js_j$ is smooth iff the coefficients $c^j$ are smooth real values functions on $U$.
\end{proposition}

\begin{customproof}
  \

  Note that if $s_1, \dots, s_r$ are the frame over some trivialization, then we win, by the prior lemma.

  $(\implies)$ If the $c^j$ are smooth functions on $U$, then $s = \sum c^js_j$ is a smooth function on $U$.

  $(\impliedby)$ Suppose $s = \sum c^js_j$ is a smooth section of $E$ over $U$. Fix a point $p \in U$, choose a trivializing open set $V \subset U$ for $E$ containing $p$, with frame $t_1, \dots, t_r$. Then, $s = \sum b^i t_j$, and $s_j = \sum a^i_j t_i$, and all the coefficients are smooth functions by the prior lemma. Then, $s = \sum_{i, j} c^ja^i_jt_i = \sum b^i t_i$, for some more coefficients $c^j$, and $b^i = c^ja^i_j$ and by some linear algebra bullshit, we get that they are smooth.
\end{customproof}

\section{Bump Functions and Partitions of Unity}

\subsection{$C^\infty$ Bump Functions}
\

\begin{definition}[Support]
  \

  The \textbf{support} of a function $f: M \to \mathbb{R}$ is the closure of the set of all points $p \in M$ where $f(p) \neq 0$, denoted $$\text{supp} f = \overline{\{q \in M \mid f(q) \neq 0\}}$$
\end{definition}

\begin{definition}[Bump Functions Supported in $U$]
  \

  A \textbf{bump function at $q$ supported in $U$} is any continuous nonnegative function $\rho$ on $M$, which takes on the value $1$ in a neighborhood of $q$, and where $\text{supp}(\rho) \subset U$. These functions exist (result in analysis).
\end{definition}

\begin{proposition}
  \

  Suppose $f$ is a smooth function defined on a neighborhood $U$ of a point $p$ in a manifold $M$. There is a function $\bar{f}$ on $M$ that agrees with $f$ on some (possibly smaller) neighborhood of $p$.
\end{proposition}

\begin{customproof}
  \

  Let $\bar{f} =
  \begin{cases}
    \rho(q)f(q) & q \in U\\
    0 & q \notin U
  \end{cases}$Where $\rho$ is a bump function that is supported in $U$, and identically $1$ on a neighborhood $V$ of $p$. Then for $p \in V$, $\bar{f} = \rho(p)f(p) = f(p)$.
\end{customproof}

\subsection{Partitions of Unity}
\

\begin{definition}[Partitions of Unity]
  \

  A \textbf{smooth partition of unity} on a manifold is a collection of nonnegative smooth functions $\{\rho_\alpha: M \to \mathbb{R}\}_{\alpha \in A}$ such that
  \begin{enumerate}
    \item The collection of supports $\{\text{supp} (\rho_\alpha)\}_{\alpha \in A}$ is locally finite (meaning that for all $p \in M$, there exists a neighborhood of $p$ that intersects with only finitely many $\{\text{supp} (\rho_\alpha)\}_{\alpha \in A}$)
    \item $\sum \rho_\alpha = 1$
  \end{enumerate}If every $\text{supp}(\rho_\alpha) \subset U_\alpha$ for open sets $U_\alpha$, we say the partition of unity $\{\rho_\alpha\}_{\alpha \in A}$ is \textbf{subordinate to the open cover} $\{U_\alpha\}$.
\end{definition}

\subsection{Existence of a Partition of Unity}
\

\begin{lemma}
  \

  If $\rho_1, \dots, \rho_m$ are real-valued functions on a manifold $M$, then $\text{supp}(\sum \rho_i) \subset \bigcup \text{supp}( \rho_i)$
\end{lemma}

\begin{customproof}
  \

  Since each $\rho_i$ is nonnegative, a point $p \in M$ is supported by $\sum \rho_i$ if and only if it is supported by at least one $\rho_i$, therefore, $\text{supp}(\sum \rho_i) \subset \bigcup \text{supp}( \rho_i$). We add the subset because the union of a closure of sets is sometimes greater than the closure of the union of those sets.
\end{customproof}

\begin{proposition}
  \

  Let $M$ be a compact manifold, and $\{U_\alpha\}_{\alpha \in A}$ be an open cover of $M$. There exists a smooth partition of unity $\{\rho_\alpha\}_{\alpha \in A}$ subordinate to $\{U_\alpha\}_{\alpha \in A}$.
\end{proposition}

\begin{customproof}
  \

  For each $q \in M$, find an open set $q \in U_\alpha$, and construct a bump function $\psi_q$ at $q$ supported in $U_\alpha$. Since it is a bump function at $q$, there is a neighborhood $q \in W_q$ where $\psi_q > 0$. Consider all such $\{W_q\}$, which forms an open cover of $M$.

  By compactness of $M$, there is a finite subcover $\{W_{q_i}\}$, and $\psi = \sum \psi_{q_i}$ is nonzero at each $p$, since every $p$ is a member of some $W_{q_i}$, and there is at least one $\psi_{q_i}$ nonzero on $W_{q_i}$. Then, call $\varphi_i = \frac{\psi_{q_i}}{\psi}$, and so $\sum \varphi_i = 1$

  Since $\psi > 0$, $\text{supp}(\varphi_i) = \text{supp}(\psi_{q_i}) \subset W_{q_i} \subset U_\alpha$ for some $\alpha \in A$.

  Then, we want to reindex so that the index on the partitions matches that of the $U_\alpha$ they are subordinate to. For each $i = 1, \dots, m$, choose $\tau(i) \in A$ to be an index such that $\text{supp}(\varphi_i) \subseteq U_{\tau(i)}$, then for each $\alpha \in A$, define $\rho_\alpha = \sum_{\tau(i) = \alpha} \varphi_i$.
\end{customproof}

\begin{theorem}
  \

  Let $\{U_\alpha\}_{\alpha \in A}$ be an open cover of a (not necessarily compact) manifold $M$.
  \begin{enumerate}
    \item There is a smooth partition of unity $\{\varphi_k\}_{k = 1}^\infty$ with every $\varphi_k$ having compact support, and such that for each $k$, $\text{supp}(\varphi_k) \subset U_\alpha$ for some $\alpha \in A$
    \item If we do not require compact support, there is a partition of unity $\{\rho_\alpha\}$ subordinate to $\{U_\alpha\}$
  \end{enumerate}
\end{theorem}

\section{Vector Fields}

\subsection{Smoothness of a Vector Field}
\

\begin{lemma}
  \

  Let $(U, \phi = x^1, \dots, x^n)$ be a chart on a manifold $M$. A vector field $X = \sum a^i \frac{\partial}{\partial x^i}$ on $U$ is smooth iff the coefficient functions $a^i$ are all smooth on $U$
\end{lemma}

\begin{customproof}
  \

  Basically, we want to show that a vector field is smooth as a section of the tangent bundle iff each component function is smooth.

  ($\implies$) Let $X: M \to TM$ be the vector field as defined above. At a chart $(U, \phi)$, we have $$X_p = \sum a^i(p)\left.\frac{\partial}{\partial x^i} \right\vert_p$$and so each $a^i$ can be thought of as a real valued functions. We want to show these functions must be smooth if $X$ is smooth.

  We can extend the chart $(U, x^i)$ in $M$ to a chart $(\tilde{U}, \tilde{\phi} = x^1 \circ \pi, \dots, x^m \circ \pi, c^1, \dots, c^m)$ for $c^i$ picking out the $i$th coordinate of some vector $v \in T_pM$. If $X$ is smooth, by definition, it is a smooth map $X: M \to TM$, and so $\tilde{\phi} \circ X$ must be smooth. But at any $p \in M$, $$X_p = \sum a^i(p) \left.\frac{\partial}{\partial x^i}\right\vert_p = \sum c^i(X_p) \left.\frac{\partial}{\partial x^i}\right\vert_p $$and since $\tilde{\phi} \circ X$ is smooth, and $c^i(X_p) = (\tilde{\phi} \circ X)^i(p)$, $c^i$ varies smoothly across $p \in M$. Then, matching coefficients, we see $a^i(p) = c^i(X_p)$ so it is smooth as well.

  ($\impliedby$) If $(U, \phi)$ is a chart about $p\in M$, $(\tilde{U}, \tilde{\phi} = (x^1 \circ \pi, \dots, x^m \circ \pi, c^1, \dots, c^m))$ is a chart about $(p, X_p) \in TM$, then $a^i = c^i \circ X$, hence if all $a^i$ is smooth, $X$ defines a smooth map from $X: M \to TM$, and hence is a smooth section.
\end{customproof}

\begin{proposition}
  \

  Let $X$ be a vector field on a manifold $M$. TFAE:
  \begin{enumerate}
    \item The vector field $X$ is smooth on $M$
    \item The manifold $M$ has an atlas such that on any chart $(U, x^i)$ of the atlas, the coefficients $a^i$ of $X = \sum a^i \frac{\partial}{\partial x^i}$ relative to the frame $\frac{\partial}{\partial x^i}$ are all smooth.
    \item On any chart $(U, x^i)$ on the manifold $M$, the coefficients $a^i$ of $X = \sum a^i \frac{\partial}{\partial x^i}$ relative to the frame $\frac{\partial}{\partial x^i}$ are all smooth.
  \end{enumerate}
\end{proposition}

\begin{customproof}
  \

  $(2) \implies (1)$ By the prior lemma, $X$ is smooth on every chart $(U, \phi)$ of an atlas of $M$, so $X$ is smooth on $M$

  $(1) \implies (3)$ A smooth vector field on $M$ is smooth on every chart $(U,\phi)$ of $M$, then the prior lemma implies $(3)$.

  $(3) \implies (2)$ If something is true on every chart of a manifold, and those charts assemble an atlas, then it is true for some atlas.
\end{customproof}

\begin{proposition}
  \

  A vector field $X$ on $M$ is smooth iff for every smooth function $f$ on $M$, the function $Xf$ is smooth on $M$
\end{proposition}

\begin{customproof}
  \

  As a reminder, $Xf(p) = X_pf = \sum a^i(p) \frac{\partial f}{\partial x^i}(p)$.

  ($\implies$) Suppose $X$ is smooth and $f$ is smooth. Then, $\frac{\partial f}{\partial x^i}$ is smooth, $a^i$ is smooth, and $Xf = \sum a^i \frac{\partial f}{\partial x^i}$ is smooth.

  ($\impliedby$) Choose $f = x^i$ on coordinates, then $X x^i = a^i$.
\end{customproof}

\begin{proposition}
  \

  Suppose $X$ is a smooth vector field defined on a neighborhood $U$ of a point $p$ in a manifold $M$. Then, there is a smooth vector field $\tilde{X}$ on $M$ that agrees with $X$ on some possibly smaller neighborhood of $p$.
\end{proposition}

\begin{customproof}
  \

  Choose a bump function $\rho: M \to \mathbb{R}$ supported in $U$ that is identically equal to $1$ in a neighborhood $V$ of $p$, define $\tilde{X}(q) =
  \begin{cases}
    \rho(q)X_q & q \in U\\
    0 & q \notin U
  \end{cases}$
\end{customproof}

\subsection{Integral Curves}
\

\begin{definition}[Integral Curves]
  \

Let $X$ be a smooth vector field on a manifold $M$, $p \in M$. An \textbf{integral curve} of $X$ is a smooth curve $c: ]a,b[ \to M$ such that $c'(t) = X_{c(t)}$ for all $t \in ]a,b[$

  An integral curve is \textbf{maximal} if it's domain cannot be extended to a larger interval
\end{definition}

\subsection{Local Flows}
\

\begin{theorem}
  \

Let $V$ be an open subset of $\mathbb{R}^n$, $p_0 \in V$, and $f: V \to \mathbb{R}^n$ a smooth function. Then, the differential equation $$\frac{dy}{dt} = f(y)$$with the initial condition $y(0) = p_0$ has a unique smooth solution $y: ]a(p_0), b(p_0)[ \to V$, where $]a(p_0), b(p_0)[$ is the maximal open interval containing $0$ on which $y$ is defined.
\end{theorem}

\begin{theorem}
  \

Let $V$ be an open subset of $\mathbb{R}^n$, and $f: V\to \mathbb{R}^n$ a smooth function on $V$. For each point $p_0 \in V$, there are a neighborhood $W$ of $p_0$ in $V$, a number $\varepsilon > 0$, and a smooth function $$y: ]-\varepsilon, \varepsilon[ \times W \to V$$such that $\frac{\partial y}{\partial t}(t,q) = f(y(t,q))$, and $y(0,q) = q$ for all $(t,q) \in ]-\varepsilon, \varepsilon[ \times W$
\end{theorem}

\begin{definition}[Local Flows]
  \

A \textbf{local flow} about a point $p$ in an open set $U$ of a manifold is a smooth function $F: ]-\varepsilon, \varepsilon[ \times W\ to U$, where $\varepsilon > 0$ and $W$ is a neighborhood of $p \in U$, such that
  \begin{enumerate}
    \item $F(0,q) = q$for all $q \in W$
    \item $F(t, F(s,q)) = F_{t+s}(q)$ whenever both sides are defined.
  \end{enumerate}

  Given a local flow, $F(0,q) = q$, and $\frac{\partial F}{\partial t}(0,q) = X_{F(0,q)} = X_q$ so we can recover the vector field from it's flow.
\end{definition}

\subsection{The Lie Bracket}
\

\begin{remark}[Motivation]
  \

  For two vector fields $X,Y$, they individually satisfy the derivation property, $Y(fg) = (Yf)g + f(Yg)$, but their composition does not:,$$XY(fg) = X(Y(fg)) = X((Yf)g + f(Yg)) = (XYf)g + (Yf)(Xg) + (Xf)(Yg) + f(XYG)$$we get two extra terms $(Yf)(Xg) + (Xf)(Yg)$. Since these terms are symmetric, $XY(fg) - YX(fg)$ will make those terms disappear. The Lie Bracket seeks to measure how much the composition of two vector fields fails to satisfy the Leibniz property by doing this subtraction.
\end{remark}

\begin{definition}[Lie Bracket]
  \

  Given two vector fields $X,Y$ on $U$, some $p \in U$, their \textbf{Lie bracket} is given by $[X,Y]_p(f) = (X_pY - Y_pX)f$
\end{definition}

\begin{proposition}
  \

  If $X,Y$ are smooth on $M$, then $[X,Y]$ is smooth on $M$.
\end{proposition}

\begin{customproof}
  \

  $[X,Y]$ is smooth on $M$ iff for all $f$ smooth, $[X,Y]f$ is smooth. But $$[X,Y]f = (XY-YX)f = X(Yf) - Y(Xf)$$and since $Xf, Yf$ smooth, so are $X(Yf)$ and $Y(Xf)$, hence this is indeed a smooth vector field.
\end{customproof}

\begin{definition}[Lie Algebras]
  \

  Let $K$ be a field. A \textbf{Lie algebra} over $K$ is a vector space $V$ over $K$ with a product $[,]: V \times V\to V$ called the \textbf{bracket}, where the following are true:

  \begin{enumerate}
    \item $[aX + bY, Z] = a[X,Z] + b[Y,Z]$ \\ $[Z, aX + bY] = a[Z,X] + b[Z,Y]$
    \item $[Y,X] = -[X,Y]$
    \item $\sum_\text{cyclic}[X,[Y,Z]] = [X,[Y,Z]] + [Y, [Z,X]] + [Z, [X,Y]] = 0$
  \end{enumerate}The last criterion is called the \textbf{Jacobi identity}
\end{definition}

\begin{definition}[Derivations of Lie Algebras]
  \

  A \textbf{derivation} of a Lie algebra $V$ over a field $K$ is a $K$-linear map $D: V \to V$ satisfying the product rule: $D[Y,Z] = [DY,Z] + [Y,DZ]$
\end{definition}

\subsection{The Pushforward of Vector Fields}
\

\begin{remark}
  \

  If $F: M \to N$, then $F_*: T_pM \to T_pN$ is the pushforward. In general, the pushforward does not extend to vector fields, because $F(p) = F(q)$ and so $X_p$ and $X_q$ will both be pushed forward to tangent vectors at $F(p)$ but there is no reason that $F_*(X_p) = F_*(X_q)$

  If $F$ is a diffeomorphism, then the pushforward makes sense.

  In absence of a well defined notion of a pushforward, we introduce vector field relations.
\end{remark}

\subsection{Related Vector Fields}
\

\begin{definition}[$F$-related]
  \

  Let $F: M \to N$ be a smooth map. A vector field $X$ on $M$ is $F$-related to a vector field $\tilde{X}$ on $N$ if for all $p \in M$, $F_{*, p}(X_p) = \tilde{X}_{F(p)}$
\end{definition}

\begin{proposition}
  \

  Let $F: M \to N$ be a smooth map of manifolds. A vector field $X$ on $M$ and a vector field $\tilde{X}$ on $N$ are $F$-related iff for all $g \in C^\infty(M)$, $$X(g \circ F) = (\tilde{X}g) \circ F$$
\end{proposition}

\begin{customproof}
  \

  We have that $$F_{*,p}(X_p)g = \tilde{X}_{F(p)}g$$$$X_p(g \circ F) = (\tilde{X}g)(F(p))$$$$(X(g \circ F))(p) = (\tilde{X}g)(F(p))$$since this is true for all $p \in M$, $X(g \circ F) = (\tilde{X}g) \circ F$. Reversing the above gives the converse.
\end{customproof}

\begin{proposition}
  \

  Let $F: M \to N$ be smooth map between manifolds. If the smooth vector fields $X$ and $Y$ are $F$-related to the smooth vector fields $\tilde{X}$ and $\tilde{Y}$ respectively, then $[X,Y]$ is $F$-related to $[\tilde{X}, \tilde{Y}]$
\end{proposition}

\begin{customproof}
  \

  For any $g \in C^\infty(M)$, $$[X,Y](g \circ F) = XY(g \circ F) - YX(g \circ F)$$$$=X((\tilde{Y}g)\circ F) - Y((\tilde{X}g) \circ F)$$$$= (\tilde{X}\tilde{Y}g) \circ F - (\tilde{Y}\tilde{X}g) \circ F$$$$ - ((\tilde{X}\tilde{Y} - \tilde{Y}\tilde{X})g) \circ F$$$$= ([\tilde{X},\tilde{Y}]g) \circ F$$
\end{customproof}

\section{Lie Groups}

\subsection{Examples of Lie Groups}
\

\begin{definition}[Lie Groups]
  \

  A \textbf{Lie group} is a smooth manifold $G$ that has the group operations multiplication and inversion: $$\mu: G \times G \to G$$$$\iota: G \to G$$which are also smooth maps between manifolds.
\end{definition}

\begin{exercise}
  \

  We denote $$\ell_a: G \to G$$$$\ell_a: x \mapsto \mu(a,x) = ax$$the left multiplication operation. Show that for any $a \in G$, $\ell_a$ is a diffeomorphism.
\end{exercise}

\begin{customproof}
  \

  As a group operation, $\ell_a$ is a bijection from $G$ to $G$. Then, let $(U, \varphi = x^i)$ be a chart on $G$. For $x,y \in G$, $\varphi \circ \mu(x,y)$ is a smooth map, so for fixed $x = a$, $\varphi \circ \mu(a,y)$ is also a smooth map. But then, $\mu(a,y) = \ell_a(y)$, hence $\ell_a(y)$ is a bijective smooth map and a diffeomorphism.
\end{customproof}

\begin{definition}[Lie Group Homomorphisms]
  \

  A map $F: G \to H$ between two Lie groups $G, H$ is a \textbf{Lie group homomorphism} if for all $a,b \in G$, $$F(ab) = F(a)F(b)$$Equivlanently, we may write $$F \circ \ell_a = \ell_{F(a)} \circ F$$
\end{definition}

\begin{proposition}
  \

  Lie group homomorphisms map the identity to the identity.
\end{proposition}

\begin{customproof}
  \

  $$F(e_G) = F(e_Ge_G) = F(e_G)F(e_G)$$so $F(e_G)$ is the identity in $H$.
\end{customproof}

\begin{example}
  \

  $\text{GL}(n, \mathbb{R})$ is a Lie group under matrix multiplication and inversion.
\end{example}

\subsection{Lie Subgroups}
\

\begin{definition}[Lie Subgroups]
  \

  A subset $H$ is a \textbf{Lie subgroup} of a Lie group $G$ fulfills all of the following:

  \begin{enumerate}
    \item A subgroup $H \leq G$ (just considering group structure)
    \item An immersed submanifold via some inclusion map
    \item The group operations on $H$ are smooth.
  \end{enumerate}
\end{definition}

\begin{remark}
  \

  By immersed submanifold, we mean that the inclusion map $\iota: H \to G$ must be injective, and an immersion.
\end{remark}

\begin{example}
  \

  Let $G$ be the torus $\mathbb{R}^2/ \mathbb{Z}^2$ and $L$ a line through the origin in $\mathbb{R}^2$ (we can think of this as a unit square with opposite edges identified). Let $\pi: \mathbb{R}^2 \to \mathbb{R}^2/ \mathbb{Z}^2$ be a projection map, and call $H = \pi(L)$. If $L$ has rational slope, then $H$ will form a closed curve, beacuse once both coefficients $(m,n)$ are integers, the curve will return to $(m,n) \sim (0,0) \in \mathbb{R}^2/ \mathbb{Z}^2$.

  If the slope is irrational, then $H$ will never loop, and $f = \pi \vert_L$ is an injective immersion, and forms a Lie subgroup of $G$.
\end{example}

\begin{proposition}
  \

  If $H$ is an abstract subgroup of a Lie group $G$ (when viewing $H$ and $G$ purely as groups), and $H$ is also a regular submanifold of $G$, then $H$ is a Lie subgroup of $G$.
\end{proposition}

\begin{customproof}
  \

  There exists an injective immersion $\iota$ from any submanifold into its parent manifold.

  Now we want to show that the inherited operations are smooth. Since $\iota$ is smooth, so is $\iota \times \iota: H \times H \to G \times G$, and since $\mu: G \times G\to G$ is smooth, and so their composition, $\mu \circ (\iota \times \iota): H \times H \to G$ is smooth, but since $H$ is a subgroup, we can ensure that it is closed under multiplication, so it is actually $\mu \circ (\iota \times \iota): H \times H \to H$, exactly the multiplication operation on $H$, hence the multiplication operation on $H$ is smooth.

  A similar argument works for inversion.
\end{customproof}

\begin{theorem}
  \

  Let $G$ be a Lie group, $H$ be an abstract subgroup of $G$, and also $H$ is a closed set in $G$ (as a topological space). Then, $H$ is an embdded Lie subgroup of $G$.
\end{theorem}

\subsection{The Matrix Exponential}
\

\begin{definition}[Norm]
  \

  The \textbf{norm} of a vector space $V$ is a real valued function $\vert\vert \cdot\vert\vert : V \to \mathbb{R}$, where for all $r\in \mathbb{R}$, $v,w \in V$, the following are true:

  \begin{enumerate}
    \item $\vert\vert v\vert\vert \geq 0$, and $\vert\vert v\vert\vert = 0$ iff $v = 0$
    \item $\vert\vert rv\vert\vert = \vert r\vert \vert\vert v\vert\vert$
    \item$ \vert\vert v+w \vert\vert \leq \vert\vert v\vert\vert + \vert\vert w\vert\vert$
  \end{enumerate}

  A vector space with a norm is called a \textbf{normed vector space} or a \textbf{normed linear space}.
\end{definition}

\begin{example}
  \

  The vector space $\mathbb{R}^{2^n} \cong \mathbb{R}^{n \times n}$ of $n \times n$ real matrices can be given the Euclidean norm. For matrix $X = [x_{ij}]$ $$\vert\vert X\vert\vert = (\sum x_{ij}^2)^\frac{1}{2} $$
\end{example}

\begin{definition}[Matrix Exponential]
  \

  For a matrix $X \in \mathbb{R}^{n \times n}$, $$e^X = I + X + \frac{X^2}{2!} + \frac{X^3}{3!} + \dots$$This formula doesn't make sense yet because we aren't sure if the sum converges, but we will use the rest of the chapter to prove this.
\end{definition}

\begin{definition}[Normed Algebra]
  \

  A \textbf{normed linear space} $V$ that is also an $\mathbb{R}$-algebra is called a \textbf{normed algebra} if for all $v, w \in V$, $$\vert\vert vw\vert\vert \leq \vert\vert v\vert\vert \vert\vert w\vert\vert  $$
\end{definition}

\begin{proposition}
  \

  For $X,Y$ in $\mathbb{R}^{n \times n}$, $$\vert\vert XY\vert\vert\leq \vert\vert X\vert\vert \vert\vert Y\vert\vert$$
\end{proposition}

\begin{customproof}
  \

  $$\vert\vert XY\vert\vert ^2 \leq \sum_{i,j} (\sum_k x^2_{ik} \sum_{k} y^2_{kj}) = (\sum_{i,k} x^2_{ik})(\sum_{j,k} y^2_{kj}) = \vert\vert X\vert\vert^2 \vert\vert Y\vert\vert^2  $$
\end{customproof}

\begin{proposition}
  \

  Let $V$ be a normed algebra. If $a \in V$, $s_m \to s$ is a sequence in $V$ that converges to $s$, then $as_m \to as$

  If $a \in V$, $\sum_{k=0}^\infty b_k$ is convergent, then $a\sum_{k=0}^\infty b_k = \sum_{k=0}^\infty ab_k$.
\end{proposition}

\begin{proposition}
  \

  $$e^X = I + X + \frac{X^2}{2!} + \frac{X^3}{3!} + \dots$$converges
\end{proposition}

\begin{customproof}
  \

  It is sufficient to show $$\sum_{k=0}^\infty \vert\vert \frac{X^k}{k!}\vert\vert$$converges.

  Note that since $\mathbb{R}^{n \times n}$ is a normed algebra, $\vert\vert X^k\vert\vert \leq \vert\vert X\vert\vert^k$. Then,
  $$\sum_{k=0}^\infty \vert\vert \frac{X^k}{k!}\vert\vert \leq \sqrt{n} + \vert\vert X\vert\vert + \frac{1}{2} \vert\vert X\vert\vert^2 + \frac{1}{3!} \vert\vert X\vert\vert^3 + \dots = (\sqrt{n}-1)+e^{\vert\vert X\vert\vert }$$and this is obviously nonzero for any matrix $X$, since $\vert\vert X\vert\vert$ is just a number.
\end{customproof}

\begin{proposition}
  \

  $$\frac{d}{dt}e^{tX} = Xe^{tX} = e^{tX}X$$
\end{proposition}

\begin{customproof}
  \

  $$\frac{d}{dt}e^{tX} = \frac{d}{dt}[I + tX + \frac{t^2X^2}{2!} + \frac{t^3X^3}{3!} + \dots ] = X(I + tX + \frac{t^2X^2}{2!} + \dots) = Xe^{tX}$$$$ = (I + tX + \frac{t^2X^2}{2!} + \dots)X = e^{tX}X$$
\end{customproof}

\begin{remark}
  \

  Everything above works for complex valued matrices with the Hermitian norm.
\end{remark}

\subsection{Trace of a Matrix}
\

\begin{definition}[Trace]
  \

  For $[x_{ij}] = X \in \mathbb{R}^{n \times n}$, the \textbf{trace} is given by $$\text{tr}(X) = \sum x_{ii}$$
\end{definition}

\begin{proposition}
  \

  \begin{enumerate}

    \item For $X,Y \in \mathbb{R}^{n \times n}$, $\text{tr}(XY) = \text{tr}(YX)$
    \item For $X \in \mathbb{R}^{n \times n}$ and $A \in \text{GL}(n, \mathbb{R})$, $\text{tr}(AXA ^{-1}) = \text{tr}(X)$.
  \end{enumerate}
\end{proposition}

\begin{customproof}
  \

  \begin{enumerate}
    \item $$\text{tr}(XY) = \sum_i (XY)_{ii} = \sum_i x_{ik}y_{ki} = \sum_i y_{ki}x_{ik} = \text{tr}(YX)$$
    \item $\text{tr}(AXA ^{-1}) = \text{tr}(XAA ^{-1}) = \text{tr}(X)$
  \end{enumerate}
\end{customproof}

\begin{definition}[Eigenvalues]
  \

  The \textbf{eigenvalues} of a matrix $X$ is the solution to the equation $$\det(\lambda I - X) = 0$$
\end{definition}

\begin{proposition}
  \

  The trace of a matrix is equal to the sum of it's complex eigenvalues (remember real numbers are complex numbers too).
\end{proposition}

\begin{customproof}
  \

  Given any matrix $X \in \mathbb{\mathbb{C}}^{n \times n}$, there exists a nonsingular matrix $A \in \text{GL}(n, \mathbb{C})$ such that $AXA ^{-1}$ is a upper triangular matrix with it's eigenvalues on the diagonal.

  Then, $$\text{tr}(X) = \text{tr}(AXA ^{-1}) = \sum_i \lambda_i$$
\end{customproof}

\begin{proposition}
  \

  For any $X \in \mathbb{R}^{n \times n}$, $\det(e^X) = e^{\text{tr}(X)}$
\end{proposition}

\begin{customproof}
  \

  If $X$ is upper triangular, then $$e^X = \sum \frac{1}{k!} X^k = \sum \frac{1}{k!}
  \begin{bmatrix}
    \lambda_1^k & & *\\
    & \ddots & \\
    0 & & \lambda_n^k
  \end{bmatrix} =
  \begin{bmatrix}
    e^{\lambda_1} & & *\\
    & \ddots & \\
    0 & & e^{\lambda_n}
  \end{bmatrix}$$And so $\det(e^X) = \Pi e^{\lambda_i} = e^{\sum \lambda_i} = e^{\text{tr} X}$

  We can find a nonsingular matrix $A$ such that $AXA ^{-1}$ is upper triangular matrix. Then, $$e^{AXA ^{-1}} = I + AXA ^{-1} + \frac{(AXA ^{-1})^2}{2!} + \dots = I + AXA ^{-1} + \frac{AX^2A ^{-1}}{2!} + \dots = Ae^XA ^{-1}$$And $\det(e^X) = \det(Ae^XA ^{-1}) = \det(e^{AXA ^{-1}}) = e^{\text{tr}(X)}$
\end{customproof}

\begin{corollary}
  \

  $e^X$ is never a singular matrix, since it always has nonzero determinant.
\end{corollary}

\subsection{The Differential of $\text{det}$ at the Identity}
\

\begin{proposition}
  \

  For any $X \in \mathbb{R}^{n \times n}$, $\det_{*, I}(X) = \text{tr}(X)$
\end{proposition}

\begin{customproof}
  \

  We use a curve $c(t) = e^{tX}$ to compute, because $c(0) = I$, and $c'(0) = X$.

  Then, $$\det_{*, I} = \left.\frac{d}{dt}\right\vert_{t=0} \det(e^{tX}) = \left.\frac{d}{dt}\right\vert_{t=0} e^{t\text{tr}(X)} = \left. \text{tr}(X)e^{t \text{tr}(X)} \right\vert_{t=0} = \text{tr}(X)$$
\end{customproof}

\section{Lie Algebras}

\subsection{Tangent Space at the Identity of a Lie Group}
\

\begin{remark}
  \

  Given a Lie group $G$, $g \in G$, the left multiplication action $\ell_g: G \to G$ is a diffeomorphism with inverse $\ell_{g ^{-1}}$. Looking at just the identity, the pushforward is an isomorphism of tangent spaces $$\ell_{g*} = (\ell_g)_{*,e}: T_eG \to T_gG$$So by describing $T_eG$, we have a description of every $T_gG$.
\end{remark}

\begin{example}
  \

  In $\text{GL}(n, \mathbb{R})$, the tangent space at any $g$ is just $\mathbb{R}^{n \times n}$, the space of $n \times n$ real matrices.

We want to find the pushforward. Let $\gamma: ]-\varepsilon, \varepsilon[ \to \text{GL}(n, \mathbb{R})$ be a curve such that $\gamma(0) = I$ and $\gamma'(0) = X$. Then, $$(\ell_g)_{*, I}(X) = \left.\frac{d}{dt}\right\vert_{t=0} \ell_g \circ \gamma(t) = \left.\frac{d}{dt}\right\vert_{t=0}g\gamma(t) = g\gamma'(0) = gX$$
\end{example}

\begin{proposition}
  \

  The tangent space $T_I \text{SL}(n, \mathbb{R})$ is the subspace of all matrices with zero trace.
\end{proposition}

\begin{customproof}
  \

There is a curve $\gamma: ]-\varepsilon, \varepsilon[ \to \text{SL}(n, \mathbb{R})$ with $\gamma(0) = I$ and $\gamma'(0) = X$ for any $x \in T_I\text{SL}(n,\mathbb{R})$, and $\gamma$ lies entirely within $\text{SL}(n, \mathbb{R})$, so $\det \gamma(t) = 1$ for any $t$.

  If we differentiate, we get that
  \begin{align*}
    \left.\frac{d}{dt} \det(\gamma(t)) \right\vert_{t=0} &= (\det \circ \gamma)_*(\left.\frac{d}{dt}\right\vert_0) && \text{definition of pushforward acting on a tangent vector}\\
    &= \text{det}_{*,\gamma(0)} (\gamma_* \left.\frac{d}{dt}\right\vert_0) && \text{by chain rule}\\
    &= \text{det}_{*, I}(\gamma'(0)) && \text{the pushforward of $\left.\frac{d}{dt}\right\vert_0$ by $\gamma$ is $\gamma'(0)$}\\
    &= \text{det}_{*,I}(X)&&\text{we defined this}\\
    &= \text{tr}(X) = 0 && \text{proved in a prior proposition}
  \end{align*}So $T_I\text{SL}(n, \mathbb{R})$ is just the space of all matrices with zero trace (and this does indeed form a vector space).
\end{customproof}

\begin{proposition}
  \

  The tangent space $T_IO(n)$ (which is the space of $n \times n$ matrices $A$ where $AA^T = I$), is the space of all skew-symmetric matrices (where $B = -B^T$).
\end{proposition}

\begin{customproof}
  \

Just trust that $O(n)$ forms a manifold. Let $\gamma: ]-\varepsilon, \varepsilon[ \to O(n)$ be a curve and $\gamma(0) = I$, $\gamma'(0) = X$ for some $X \in T_I O(n)$. Then, for all $t$, we have $\gamma(t)\gamma(t)^T = I$, differentiating we get $\gamma'(t)\gamma(t)^T + \gamma(t)\gamma'(t)^T = 0$ and evaluating at $T=0$, we get $X^T + X = 0$, which only happens if $X$ is skew-symmetric.
\end{customproof}

\subsection{Left-Invariant Vector Fields on a Lie Group}
\

\begin{definition}[Left Invariance]
  \

  Let $X$ be a vector field on $G$ a Lie group, $X$ need not be smooth. $X$ is \textbf{left invariant} if for all $g \in G$, $$\ell_{g, *}(X) = X$$or for every $g \in G$, every $h \in G$, $$\ell_{g, *}(X_h) = X_{gh}$$or $X$ is $\ell_g$-related to itself for all $g \in G$ (all three definitions above are equivalent).

  As a consequence, every left invariant vector field $X$ is completely defined by the value $X_e$ at the identity, since $X_g = \ell_{g, *}(X_e)$ for any other $g \in G$.

  Also, we can construct a left invariant vector field $A$ out of any vector $A_e \in T_eG$, by setting $A_g = \ell_{g, *}(A_e)$. We call $A$ the \textbf{left invariant vector field generated by $A_e$}
\end{definition}

\begin{proposition}
  \

  Any left invariant vector field on a Lie group $G$ is smooth.
\end{proposition}

\begin{customproof}
  \

  It suffices to show for any smooth function $f: G \to \mathbb{R}$, $Xf$ is also smooth.

  Let $\gamma: I \to G$ be a smooth curve defined on some interval $0 \in U \subseteq \mathbb{R}$. Let $\gamma(0) = e$ and $\gamma'(0) = X_e$. For all $g \in G$, $g\gamma(0) = g$ and $(g\gamma)'(0) = \ell_{g*}\gamma'(0) = X_g$.

  Then, $$(Xf)(g) = X_gf = \left.\frac{d}{dt}\right\vert_{t=0} f(gc(t))$$But since $f(gc(t))$ is a composition of smooth functions, it is smooth. Then, $ \left.\frac{d}{dt}\right\vert_{t=0} f(gc(t))$ is also smooth.
\end{customproof}

\begin{proposition}
  \

  If $X$ and $Y$ are left invariant, $[X,Y]$ is left invariant
\end{proposition}

\begin{customproof}
  \

  By a prior proposition, if $X$ is $\varphi$-related to $\bar{X}$, and same with $Y$, then $[X,Y]$ is $\varphi$-related to $[\bar{X}, \bar{Y}]$.

  For two left invariant vector fields on a Lie group $G$, $X$ and $Y$ $\ell_g$ related to themselves, so $[X,Y]$ is $\ell_g$ related to itself, so it is also left invariant.
\end{customproof}

\subsection{The Lie Algebra of a Lie Group}
\

\begin{remark}
  \

  Recall that the a Lie algebra is a vector space with a bracket operation that is bilinear, alternating, and satisfies the Jacobi identity.

  The space of vector fields on a smooth manifold $M$ is a Lie algebra, denoted $\mathfrak{X}(M)$
\end{remark}

\begin{definition}[Lie Subalgebra]
  \

  If $\mathfrak{g}$ is a Lie algebra, $\mathfrak{h} \subset \mathfrak{g}$ is a Lie subalgebra if it is a vector subspace of $\mathfrak{g}$, and if the bracket operation inherited from $\mathfrak{g}$ is algebraically closed.
\end{definition}

\begin{example}
  \

  The space of left invariant vector fields of a Lie group $G$, denoted $L(G)$, forms a Lie subalgebra of $\mathfrak{X}(G)$, since if $X,Y$ are left invariant, so is $[X,Y]$.
\end{example}

\begin{remark}
  \

  Recall that since every left-invariant vector field is generated by some vector in $T_eG$, there is a vector space isomorphism $\varphi: T_eG \to L(G)$.

  This isomorphism is very helpful, because $L(G)$ has Lie bracket structure we can bring to $T_eG$, and $T_eG$ has pushforwards (which if we recall, aren't always well defined for vector fields in general)
\end{remark}

\begin{definition}
  \

  For $A,B \in T_eG$, if they generate left invariant vector fields $\tilde{A}, \tilde{B}$, the Lie bracket operation on this space is given by $$[A,B] = [\tilde{A}, \tilde{B}]_e$$Basically, we bracket the two vector fields generated by those vectors, then take the vector at the identity in that vector field.
\end{definition}

\begin{proposition}
  \

  If $A,B \in T_eG$ are vectors, and they generate left invariant vector fields $\tilde{A}, \tilde{B}$, then $$[\tilde{A},\tilde{B}] = \squigglyover{[A,B]}$$
\end{proposition}

\begin{customproof}
  \

  We have $$[A,B] = [\tilde{A}, \tilde{B}]_e$$looking at the vector fields they generate, we get $$\squigglyover{[A,B]} = [\tilde{A}, \tilde{B}]$$Since the vector field generated by $[\tilde{A}, \tilde{B}]_e$ is just $[\tilde{A}, \tilde{B}]$.
\end{customproof}

\begin{definition}[Lie Algebras of a Lie Group]
  \

  Given a Lie group $G$, $T_eG$ with the Lie bracket from $L(G)$ gives the \textbf{Lie algebra} of the Lie group $G$.
\end{definition}

\subsection{The Lie Bracket on $\mathfrak{gl}(n, \mathbb{R})$}
\

\begin{definition}
  \

  Denote the Lie algebra of $\text{GL}(n, \mathbb{R})$ as $\mathfrak{gl}(n, \mathbb{R})$
\end{definition}

\begin{remark}
  \

  The tangent space $T_I\text{GL}(n, \mathbb{R})$ is isomorphic to $\mathbb{R}^{n \times n}$, and we identify $\sum a_{ij} \left.\frac{\partial}{\partial x_{ij}}T_I\right\vert_I \in \text{GL}(n, \mathbb{R})$ with $[a_{ij}] \in \mathbb{R}^{n \times n}$
\end{remark}

\begin{proposition}
  \

  Let $A = \sum a_{ij} \left.\frac{\partial}{\partial x_{ij}}\right\vert_I$, $B = \sum b_{ij} \left.\frac{\partial}{\partial x_{ij}}\right\vert_I$ be in $T_I\text{GL}(n, \mathbb{R})$. Then, $$[A,B] = [\tilde{A},\tilde{B}]_I = \sum_{i,j} (\sum_ka_{ik} b_{kj} - b_{ik}a_{kj})\left.\frac{\partial}{\partial x_{ij}}\right\vert_I$$and abusing a bit of notation, as matrices, $$[A,B] = AB - BA$$
\end{proposition}

\begin{customproof}
  \

  Just some matrix multipliation type shit.
\end{customproof}

\subsection{The Pushforward of Left-Invariant Vector Fields}
\

\begin{definition}[Pushforward of a Left-Invariant Vector Field]
  \

  Let $F: G \to H$ be a Lie group homomorphism, $A \in T_eG$, $\tilde{A}$ the vector field $A$ generates. Then, define $F_*: L(G) \to L(H)$ by $$F_*(\tilde{A}) = \squigglyover{F_{*,e}A}$$
\end{definition}

\begin{proposition}
  \

  If $F: G \to H$ is a Lie group homomorphism and $X$ is a left invariant vector field on $G$, then $F_*X$ is $F$-related to $X$.
\end{proposition}

\begin{customproof}
  \

  We want to show for all $g \in G$, $F_{*,g}X_g = (F_*X)_{F(g)}$.

  \begin{align*}
    F_{*,g}X_g &= F_{*,g}\ell_g(X_e) && \text{by left invariant vector field}\\
    &= (F \circ \ell_g)_{*, e}(X_e) && \text{by chain rule}\\
    &= (\ell_{F(g)} \circ F)_{*, e}(X_e) && \text{by rules of Lie group morphisms}\\
    &= \ell_{F(g), *} \circ F_{*,e}(X_e) && \text{by chain rule}\\
    &= \squigglyover{(F_{*,e}X_e)}_{F(g)} && \text{by left invariance, $\ell_{g,*}(X_h) = X_{gh}$}\\
    &= (F_*X)_{F(g)}
  \end{align*}
\end{customproof}

\begin{definition}
  \

  $F_*X$ is called the \textbf{pushforward of $X$ under $F$}
\end{definition}

\subsection{The Differential as a Lie Algebra Homomorphism}
\

\begin{proposition}
  \

  If $F: G \to H$ is a Lie group homomorphism, then its differential at the identity, $F_* = F_{*,e}: T_eG \to T_eH$ is a Lie algebra homomorphism, in other words, it is a linear map where $F_*[A,B] = [F_*A, F_*B]$
\end{proposition}

\begin{customproof}
  \

  $F_*\tilde{A}$ is $F$-related to $\tilde{A}$, and same for $B$, so $[F_*\tilde{A}, F_*\tilde{B}]$ is $F$-related to $[\tilde{A}, \tilde{B}]$, so $F_*([\tilde{A},\tilde{B}]_e) = [F_*\tilde{A}, F_*\tilde{B}]_e$, and taking everything back to the tangent space, $F_*[A,B] = [F_*A, F_*B]$.
\end{customproof}

%21000 words and counting!

\section{Differential $1$-Forms}

\subsection{The Differential of a Function}
\

\begin{definition}[Cotangent Space]
  \

  Let $M$ be a smooth manifold, $p \in M$. The \textbf{cotangent space} of $p$ is $$T_p^*M = (T_pM)^* = \text{Hom}(T_pM, \mathbb{R})$$the set of linear transformations from $T_pM$ to $\mathbb{R}$.

  An element of the cotangent space is called a \textbf{covector}, so it is a linear function $$\omega_p: T_pM \to \mathbb{R}$$

  A \textbf{covector field} or a \textbf{differential $1$-form} is a function that assigns each $p \in M$ a covector at $p$.
\end{definition}

\begin{remark}
  \

  Vector fields cannot be pushed forward under a map in general, but covector fields can be pulled back. This is becase a smooth map can assign two different points in the preimage to the same point in the image, but cannot assign the same point in the preimage to two different points in the image.
\end{remark}

\begin{definition}[Differential of a Function]
  \

  If $f: M \to \mathbb{R}$ is a smooth function, its \textbf{differential} is defined to be the $1$-form $df$ on $M$, such that for any $p \in M$, $X_p \in T_pM$, $$(df)_p(X_p) = X_pf$$
\end{definition}

\begin{proposition}
  \

  If $f: M \to \mathbb{R}$ is a smooth function, then for $p \in M$, $X_p \in T_pM$, $$f_*(X_p)= (df)_p(X_p)\left.\frac{d}{dt}\right\vert_{f(p)}$$
\end{proposition}

\begin{customproof}
  \

  Since the tangent space of any point in $\mathbb{R}$ has only one vector, we have$$f_*(X_p) = a \left.\frac{d}{dt}\right\vert_{f(p)}$$for some coefficient $a$. Applying both sides to $x$, we get $$a = f_*(X_p)(t) = X_p(t \circ f) = X_pf = (df)_p(X_p)$$
\end{customproof}

\subsection{Local Expression for a Differential $1$-Form}
\

\begin{proposition}
  \

  At each point $p \in U$, the covectors $(dx^1)_p, \dots, (dx^n)_p$ form a basis for the cotangent space $T_p^*M$ dual to the basis $\left.\frac{\partial}{\partial x^1}\right\vert_p, \dots, \left.\frac{\partial}{\partial x^n}\right\vert_p$ for the tangent space $T_pM$
\end{proposition}

\begin{customproof}
  \

  For $i,j$, we have $(dx^i)_p \left.\frac{\partial}{\partial x^j}\right\vert_p = \left.\frac{\partial x^i}{\partial x^j}\right\vert_p = \delta_i^j$. Since we have a set of functions $f_j$ where $f_j$ is dual to $e^j \in T_pM$, we know $f_j$ forms a basis for the dual space $T_p^*M$.

  Since every vector can be broken down into components $X_p = \sum a^i \left.\frac{\partial}{\partial x^i}\right\vert_p$, the way every $1$-form acts on vectors can also be broken down into components $$\omega = \sum a_i dx^i$$where $a_i: M \to \mathbb{R}$
\end{customproof}

\begin{proposition}
  \

  If $(U, x^i)$ is a chart on $M$, $f: U \to \mathbb{R}$, then$$df = \sum \frac{\partial f}{\partial x^i} dx^i$$
\end{proposition}

\begin{customproof}
  \

  By the prior proposition, $$df = \sum a_idx^i$$Applying both sides to $\frac{\partial}{\partial x^j}$, we have $$\frac{\partial f}{\partial x^j} = \sum_i a_i \delta_j^i = \alpha_j$$
\end{customproof}

\subsection{The Cotangent Bundle}
\

\begin{definition}[Cotangent Bundle]
  \

  The \textbf{cotangent bundle} is given by $$T^*M = \bigsqcup_{p \in M} T^*_pM$$Much like the tangent bundle, given $(U, \varphi=x^i)$ a chart, there is a map $$\bar{\varphi}: T^*U \to \varphi(U) \times \mathbb{R}^n$$$$\bar{\varphi}: \alpha_p\mapsto (x^1(p), \dots, x^m(p), c^1(\alpha), \dots, c^m(\alpha))$$For $\alpha = \sum c^i(\alpha)dx^i\vert_p \in T_p^*M$, and becomes a smooth manifold.

  More prescisely, the triple $(T^*M, M, \pi)$ is the cotangent bundle, with $T^*M$ being the \textbf{total space} and $M$ being the \textbf{base space}.

  Also, by virtue of $M$ being a (co)tangent bundle, it is also a vector bundle.
\end{definition}

\begin{example}
  \

  There is a $1$-form on $T^*M$ that can be defined independently of coordinates. A point in $T^*M$ is just a covector $\omega_p \in T_p^*M$. Let $X_{\omega_p}$ be a tangent vector to $T^*M$ at $\omega_p$, and consider $\pi_*(X_{\omega_p})$ is a tangent vector in $T_pM$. Then, $\omega_p(\pi_*(X_{\omega_p}))$ is a real number, and we can set the \textbf{Liouville form} $\lambda$ to be $$\lambda_{\omega_p}(X_{\omega_p}) = \omega_p(\pi_*(X_{\omega_p}))$$
\end{example}

\subsection{Characterization of Smooth $1$-Forms}
\

\begin{definition}[Smooth $1$-Forms]
  \

  A $1$-form $\omega$ is \textbf{smooth} if $\omega: M \to T^*M$ is smooth as a section of the cotangent bundle, or more prescisely, if $\omega: M \to T^*M$ is smooth.

  The space of smooth $1$-forms is denoted $\Omega^1(M)$
\end{definition}

\begin{lemma}
  \

  Let $(U,\varphi = x^i)$ be a chart on a manifold $M$. A $1$-form $\omega$ on $U$ is smooth iff the coefficient functions $a_i$ are all smooth.
\end{lemma}

\begin{customproof}
  \

  Let $\bar{\varphi}: T^*U \to \varphi(U) \times \mathbb{R}^n$ take $\omega_p$ to $(\varphi(p), c_1(\omega_p), \dots, c_m(\omega_p))$, a diffeomorphism, so $\omega: U \to T^*M$ is smooth iff $\bar{\varphi} \circ \omega$ is smooth. But for $p \in U$, $$(\bar{\varphi} \circ \omega)(p) = (x^1(p), \dots, x^m(p), c_1(\omega_p), \dots, c_m(\omega_p)) = (x^1(p),\dots, x^m(p), a_1(p),\dots, a_m(p))$$And each $x^i$ is smooth, so $\omega$ is smooth iff all $a_i$ are smooth.
\end{customproof}

\begin{proposition}
  \

  Let $\omega$ be a $1$-form on $M$. Then, TFAE:

  \begin{enumerate}
    \item The $1$-form $\omega$ is smooth on $M$
    \item The manifold $M$ has an atlas such that on any chart $(U, x^i)$, the coefficients $a_i$ of $\omega = \sum a_idx^i$ relative to the frame $dx^i$ are all smooth
    \item On any chart $(U, x^i)$, the coefficients $a_i$ relative to the frame $dx^i$ are all smooth.
  \end{enumerate}
\end{proposition}

\begin{customproof}
  \

  $(2) \implies (1)$ If something is smooth on every chart in an atlas, it is smooth on the entire manifold

  $(3) \implies (2)$ If something is true for every chart, it is true for an atlas (which is made up of charts)

  $(1) \implies (3)$ If something is smooth on $M$, it is smooth on every chart (this is something we proved far back).
\end{customproof}

\begin{corollary}
  \

  If $f$ is a smooth function on a manifold $M$, then $df$ is a smooth $1$-form on $M$.
\end{corollary}

\begin{customproof}
  \

  $df = \sum \frac{\partial f}{\partial x^i} dx^i$ and each $\frac{\partial f}{\partial x^i}$ is smooth.
\end{customproof}

\begin{proposition}
  \

  Let $\omega$ be a $1$-form on a manifold $M$. If $f$ is a function, and $X$ is a vector field on $M$, then $\omega(fX) = f\omega(X)$
\end{proposition}

\begin{remark}
  \

  Note that $f\omega$ is given by $f\omega(X_p) = f(p)\omega_p(X_p)$.
\end{remark}

\begin{customproof}
  \

  Pointwise, $\omega(fX)_p = \omega_p(f(p)X_p) = f(p)\omega_p(X_p) = (f\omega(X))_p$
\end{customproof}

\begin{proposition}
  \

  A $1$-form $\omega$ on a manifold $M$ is smooth iff for every smooth vector field $X$ on $M$, $\omega(X)$ is smooth on $M$.
\end{proposition}

\begin{customproof}
  \

  ($\implies$) $$\omega(X) = (\sum a_id^i)(\sum b^j\frac{\partial}{\partial x^j}) = \sum a_ib^i$$which is smooth

  ($\impliedby$) Let $\omega = \sum a_idx^i$, choose some chart $(U, x^i)$. By a prior theorem, we can extend the vector field $X = \frac{\partial}{\partial x^j}$, which is only defined on $U$, to $\bar{X}$ on $M$, and have $\bar{X} = X$ on some neighborhood $p \in V_p^j \subseteq U$. Then, $$\omega(\bar{X}) = a_j$$and $a_j$ smooth on $V_p^j$. But since this is true for all $a_j$ and all $p\in M$, $\omega$ is smooth.
\end{customproof}

\subsection{Pullback of $1$-Forms}
\

\begin{definition}[Pullback]
  \

  If $F: M \to N$ a smooth map, the \textbf{pullback} of a $1$-form by $F$ is defined as the dual of the pushforward: $$F^* = (F_{*,p})^*: T^*_{F(p)}M \to T^*_pM$$More concretely, if $\omega_{F(p)} \in T_{F(p)}^*M$, $X_p \in T_pM$, then $$F^*(\omega_{F(p)})(X_p) = \omega_{F(p)}(F_{*,p}X_p)$$
\end{definition}

\begin{proposition}
  \

  Let $F: M \to N$ be a smooth map. For any $h: M \to \mathbb{R}$, $F^*(dh) = d(F^*h)$
\end{proposition}

\begin{customproof}
  \

  For any $p \in M$, $X_p \in T_pM$,
  \begin{align*}
    (F^*dh)_p(X_p) &= (dh)_{F(p)}(F_{*,p}X_p) && \text{definition of the pullback}\\
    &= (F_*(X_p))h && \text{definition of $dh$}\\
    &= X_p(h \circ F) && \text{definition of $F_*$}\\
    &= X_p(F^*h) && \text{definition of $F^*$ (precomposing) of a function}\\
    &= (dF^*h)_p(X_p) && \text{definition of $d$}
  \end{align*}
\end{customproof}

\begin{proposition}
  \

  Let $F: M \to N$ be a smooth map of manifolds, $\omega, \tau \in \Omega^1(N)$ and $g \in C^\infty(N)$ Then
  \begin{enumerate}
    \item $F^*(\omega + \tau) = F^*\omega + F^*\tau$
    \item $F^*(g\omega) = (F^*g)(F^*\omega)$
  \end{enumerate}
\end{proposition}

\begin{customproof}
  \

  \begin{enumerate}
    \item For any $p \in M$, any $X_p \in T_pM$, we have
      \begin{align*}
        (F^*(\omega+\tau))_p(X_p) &= (\omega_{F(p)} + \tau_{F(p)})(F_{*,p}X_p)\\
        &= \omega_{F(p)}(F_{*,p}X_p) + \tau_{F(p)}(F_{*,p}X_p)\\ &= (F^*\omega)_p(X_p) + (F^*\tau)_p(X_p)
      \end{align*}
    \item For any $p \in M$, any $X_p \in T_pM$, we have
      \begin{align*}
        F^*(g\omega)_p(X_p) &= g(F(p))\omega_{F(p)}(F_{*,p}X_p)\\ &= (F^*g)(p)(F^*\omega)_p(X_p)
      \end{align*}
  \end{enumerate}
\end{customproof}

\begin{proposition}
  \

  The pullback $F^*\omega$ of a smooth $1$-form $\omega$ on $M$ under a smooth map $F: M \to N$ is a smooth form on $M$
\end{proposition}

\begin{customproof}
  \

  Fix $p \in M$, choose a chart $(V,\psi = y^i)$ in $N$ about $F(p)$. Since $F$ is continuous, we can find a chart $(U, \varphi = x^i)$ about $p \in M$ such that $F(U) \subset V$. Then,
  \begin{align*}
    F^* \omega &= \sum (F^*a_i)F^*(dy^i)\\
    &= \sum (a_i \circ F)dF^*y^i\\
    &= \sum (a_i \circ F)d(y^i \circ F)\\
    &= \sum (a_i \circ F) \frac{\partial F^i}{\partial x^j} dx^j
  \end{align*}and since $(a_i \circ F)$ smooth, $\frac{\partial F^i}{\partial x^j}$ smooth, the whole thing is smooth in $M$.
\end{customproof}

\subsection{Restriction of $1$-Forms to an Immersed Submanifold}
\

\begin{definition}[Restriction of $1$-Forms]
  \

  Let $S\subset M$ be an immersed submanifold with $\iota: S \to M$ an inclusion map. Since $\iota_*$ is injective, $T_pS$ is a vector subspace of $T_pM$, and so the \textbf{restriction} of $\omega \in \Omega^1(M)$ to $S$ is given by $$(\omega \vert_S)_p(v) = \omega_p(v)$$for all $p \in S$ and $v \in T_pS$ (it's the same)
\end{definition}

\begin{proposition}
  \

  If $\iota: S \hookrightarrow M$ is the inclusion map of an immersed submanifold, then $\iota^* \omega = \omega \vert_S$
\end{proposition}

\begin{customproof}
  \

  Let $p \in S$, $v \in T_pS$. Then, $$(\iota^* \omega)_p(v) = \omega_{p}(\iota_{*,p}v) = \omega_p(v)$$
\end{customproof}

\subsection{Exercises}
\

\begin{exercise}[1]
  \

  Denote the standard coordinates on $\mathbb{R}^2$ by $x,y$ and let $X = -y \frac{\partial}{\partial x} + x \frac{\partial }{\partial y}$ and $Y = x\frac{\partial}{\partial x} + y\frac{\partial}{\partial y}$ be two vector fields on $\mathbb{R}^2$. Find a $1$-form $\omega$ on $\mathbb{R}^2 \setminus \{(0,0)\}$ such that $\omega(X) = 1$ and $\omega(Y) = 0$
\end{exercise}

\begin{solution}
  \

  For every $(x,y) = p \in \mathbb{R}^2\setminus \{(0,0)\}$, let $\omega_p = a_x(p)dx + a_y(p)dy$. Then, $$\omega_p(X_p) = (a_x(p)dx + a_y(p)dy)(-y\frac{\partial}{\partial x} + x \frac{\partial}{\partial y}) = -ya_x(p) + xa_y(p) = 1$$and $$\omega_p(Y_p) = xa_x(p) + ya_y(p) = 0$$

  So we have the system $$xa_x(p) = -ya_y(p)$$$$xa_y(p) = 1 + ya_x(p)$$Solving, we get $$a_x((x,y)) = \frac{-y}{x^2+y^2}$$$$a_y((x,y)) = \frac{x}{x^2+y^2}$$
\end{solution}

\begin{exercise}[2]
  \

  Suppose $(U, x^i)$ and $(V, y^i)$ are two charts on $M$, $U \cap V \neq \varnothing$. A smooth $1$-form on $U \cap V$ has two different coordinate expressions: $$\omega = \sum a_j dx^j = \sum b_idy^i$$Find $a_j$ in terms of $b_i$
\end{exercise}

\begin{solution}
  \

  We have both forms act on $\frac{\partial}{\partial x^j}$. We get $$a_j = \sum b_i \frac{\partial y^i}{\partial x^j}$$
\end{solution}

\section{Differential $k$-Forms}

\subsection{Differential Forms}
\

\begin{remark}
  \

  Recall that a \textbf{$k$-tensor} or a \textbf{$k$-covector} on a vector space $V$ is a $k$-linear function $$f:  V^k \to \mathbb{R}$$and it is \textbf{alternating} if for some $\sigma \in S_k$, $$f(v_{\sigma(1)},\dots, v_{\sigma(k)}) = \text{sgn}(\sigma)f(v_1, \dots, v_k)$$Denote the space of alternating $k$-tensors as $$\bigwedge^k(V^*)$$
\end{remark}

\begin{definition}[$k$-forms]
  \

  A \textbf{$k$-covector field} or \textbf{differential $k$-form} on a manifold $M^m$ is a function which assigns every $p \in M$ to a $k$-covector $\omega_p \in \bigwedge^k(T^*_pM)$. If $k = m$, we call it a \textbf{top form}.

  If $\omega$ is a $k$-form, $X_i$ vector fields, then the function $\omega(X_1, \dots, X_k )$ is given by $$(\omega(X_1,\dots, X_k))(p) = \omega_p((X_1)_p, \dots, (X_k)_p)$$and this function assigns a real number to every point.
\end{definition}

\begin{proposition}
  \

  Let $\omega$ be a $k$-form on a manifold $M$. For any vector fields $X_1, \dots, X_k$, and any function $h$ on $M$, $$\omega(X_1, \dots, hX_i, \dots, X_k) = h\omega(X_1, \dots, X_i, \dots, X_k)$$
\end{proposition}

\begin{customproof}
  \

  At every point, $$\omega(X_1, \dots, hX_i, \dots, X_k)(p) = \omega_p((X_1)_p, \dots, h(p)(X_i)_p, \dots, (X_k)_p)$$and since $\omega_p$ is linear, we can pull out the $h(p)$ to get $$h(p)\omega_p((X_1)_p, \dots, (X_i)_p, \dots, (X_k)_p)$$and since this is true for all $p \in M$, we get our result above.
\end{customproof}

\subsection{Local Expression for a $k$-Form}
\

\begin{remark}
  \

  In chapter 3, we found the following: if $(U,x^i)$ is a coordinate chart, a $k$-form on $U$ is a linear combination $\omega = \sum a_I dx^I$, where $I = i_1 < \dots, < i_k$, and $1 \leq i_j \leq n$. Write as shorthand $\partial_i = \frac{\partial}{\partial x^i}$, and we have $$\frac{\partial (x^{i_1}, \dots, x^{i_k})}{\partial (x^{j_1}, \dots, x^{j_k})} = dx^I(\partial_{j_1}, \dots, \partial_{j_k}) = \det
  \begin{bmatrix}
    x^{i_1}(\partial_{j_1}) & \dots & x^{i_k}(\partial_{j_k})\\
    \vdots & \ddots & \vdots\\
    x^{i_1}(\partial_{j_k}) & \dots& x^{i_k}(\partial_{j_k})
  \end{bmatrix} = \delta^I_J =
  \begin{cases}
    1 & I=J\\
    0 & I \neq J
  \end{cases}$$
\end{remark}

\begin{proposition}
  \

  Let $(U, x^i)$ be a chart on a manifold, $f^1, \dots, f^k$ smooth functiosn on $U$. Then $$df^1 \wedge \dots \wedge df^k = \sum_{I = \{i_1 < \dots < i_k\}} \frac{\partial(f^1,\dots, f^k)}{\partial(x^{i_1}, \dots, x^{i_k})} dx^{i_1} \wedge \dots \wedge dx^{i_k}$$
\end{proposition}

\begin{customproof}
  \

  We just apply both sides to a list of coordinate vectors $\{\partial_{i_1}, \dots, \partial_{i_k}\}$ and prove equality.

  On the left, $$(df^1 \wedge\dots\wedge df^k)(\partial_{i_1}, \dots, \partial_{i_k}) = \det[\frac{\partial f^i}{\partial x^{i_j}}] = \frac{\partial(f^1,\dots, f^k)}{\partial(x^{i_1}, \dots, x^{i_k})}$$That last equality is because by definition, $\det[\frac{\partial f^i}{\partial x^{i_j}}] = \frac{\partial(f^1,\dots, f^k)}{\partial(x^{i_1}, \dots, x^{i_k})}$

  On the right we have $$\sum_J c_J dx^(\partial_{i_1}, \dots, \partial_{i_k}) = \sum_J c_J\delta_I^J = c_I$$but the $I$ coefficient is just $$c_I = \frac{\partial(f^1,\dots, f^k)}{\partial(x^{i_1}, \dots, x^{i_k})}$$
\end{customproof}

\begin{corollary}
  \

  Let $(U, x^i)$ be a chart on a manifold, let $f, f^1, \dots, f^n$ be smooth functions on $U$. Then,
  \begin{enumerate}
    \item $$df = \sum_i \frac{\partial f}{\partial x^i} dx^i$$
    \item $$df^1 \wedge \dots \wedge df^n = \det[\frac{\partial f^j}{\partial x^i}] dx^1 \wedge \dots \wedge dx^n$$
  \end{enumerate}
\end{corollary}

\subsection{The Bundle Point of View}
\

\begin{definition}[Exterior Powers]
  \

  The set $$\bigwedge^k(T^*M) = \bigsqcup_{p \in M} \bigwedge^k(T_p^*M)$$is called the \textbf{$k$th exterior power} of the cotangent bundle.

  If $(U, \varphi)$ is a chart, there is a bijection $$\bigwedge^k(T^*U) \cong \varphi(U) \times \mathbb{R}^\binom{n}{k}$$

  Much like the tangent bundle and cotangent bundle, we can bestow a topology on this structure, and a $k$-form is a section of this bundle.
\end{definition}

\subsection{Smooth $k$-Forms}
\

\begin{remark}
  \

  All the proofs for this section are basically analogues of the $1$-form versions.
\end{remark}

\begin{lemma}
  \

  Let $(U, x^i)$ be a chart on a manifold $M$. A $k$-form $\omega = \sum a_I dx^I$ is smooth iff each $a_I: M \to \mathbb{R}$ is smooth.
\end{lemma}

\begin{proposition}
  \

  Let $\omega$ be a $k$-form on a manifold $M$. TFAE:

  \begin{enumerate}
    \item The $k$-form $\omega$ is smooth on $M$
    \item The manifold $M$ has an atlas such that on every chart $(U, x^i)$ in the atlas, the coefficients $a_I$ of $\omega = \sum a_Idx^I$ relative to the coordinate frame $\{dx^I\}$ are all smooth.
    \item On every chart $(U, x^i)$, on $M$, the coefficients $a_I$ of $\omega = \sum a_Idx^I$ relative to the frame $\{dx^I\}$ are all smooth.
    \item For any $k$ smooth vector fields $X_1, \dots, X_k$ on $M$, $\omega(X_1, \dots, X_k): M \to \mathbb{R}$ is smooth on $M$.
  \end{enumerate}
\end{proposition}

\begin{proposition}
  \

  Suppose $\tau$ is a smooth differential form defined on a neighborhood $U$ of a point $p$ in a manifold $M$. Then, there is a smooth form $\bar{\tau}$ on $M$ that agrees with $\tau$ on a (possibly smaller) neighborhood of $p$.
\end{proposition}

\subsection{Pullback of $k$-Forms}
\

\begin{definition}[Pullback of $k$-Forms]
  \

  A linear map $L: V \to W$ of vector spaces induces a pullback map $L^*: \bigwedge_k(W) \to \bigwedge_k(V)$ byt $(L^*\alpha)(v_1, \dots, v_k) = \alpha(L(v_1),\dots, L(v_k))$ for $\alpha \in A_k(W)$ and $v_1, \dots, v_k\in V$

  Now, let $F: M \to N$, and the pushforward $F_{*,p}: T_pM \to T_pN$ induces a \textbf{pullback} of $k$-covectors $$F^* = (F_{*,p}): \bigwedge_k(T_{F(p)}N) \to \bigwedge_k(T_pM)$$If $\omega_{F(p)}$ is a $k$-covector at $F(p)$ in $N$, then $$F^*(\omega_{F(p)})(v_1, \dots, v_k) = \omega_{F(p)}(F_{*,p}(v_1), \dots, F_{*,p}(v_k))$$And on entire forms, $F^*\omega$ is given pointwise by $$(F^*\omega)_p = F^*(\omega_{F(p)})$$
\end{definition}

\begin{proposition}
  \

  Let $F: M \to N$ be a smooth map, $\omega, \tau$ are $k$-forms on $N$, and $a$ a real number. Then,

  \begin{enumerate}
    \item $F^*(\omega + \tau) = F^*\omega + F^*\tau$
    \item $F^*(a\omega) = aF^*\omega$
  \end{enumerate}
\end{proposition}

\begin{customproof}
  \

  \begin{enumerate}
    \item Pointwise, $$(F^*(\omega + \tau))_p(v_1, \dots, v_k) = (\omega + \tau)_{F(p)}(F_{*,p}(v_1), \dots, F_{*,p}(v_k)) =$$$$\omega_{F(p)}(F_{*,p}(v_1), \dots, F_{*,p}(v_k)) + \tau_{F(p)}(F_{*,p}(v_1), \dots, F_{*,p}(v_k)) = (F^*\omega)_p(v_1, \dots, v_k) + (F^*\tau)_p(v_1, \dots, v_k)$$
    \item Analagous
  \end{enumerate}
\end{customproof}

\subsection{The Wedge Product}
\

\begin{remark}
  \

  For any vector space $V$, $\alpha$ a $k$-tensor, $\beta$ an $l$-tensor, the wedge product $\alpha \wedge \beta$ is a $k+l$ tensor given by $$(\alpha \wedge\beta)(v_1, \dots, v_{k+k}) = \sum_{\sigma \in S_{k+l}} (\text{sgn}\sigma)\alpha(v_{\sigma(1)}, \dots, v_{\sigma(k)})\beta(v_{\sigma(k+1)}, \dots, v_{\sigma(k+l)})$$
\end{remark}

\begin{definition}[Wedge Product of Forms]
  \

  We define the \textbf{wedge product of forms} pointwise: $$(\omega \wedge \tau)_p = \omega_p \wedge \tau_p$$
\end{definition}

\begin{proposition}
  \

  If $\omega, \tau$ are smooth forms on $M$, then $\omega \wedge \tau$ is also smooth.
\end{proposition}

\begin{customproof}
  \

  Let $\omega = \sum a_I dx^I$, $\tau = \sum b_J dx^J$, then $a_I, b_J$ smooth. The wedge product is $$\omega \wedge \tau = \sum a_Ib_J dx^I\wedge dx^J$$but the coefficients are all smooth.
\end{customproof}

\begin{proposition}
  \

  For forms $\omega, \tau$ on $N$, $F: M \to N$ smooth, $$F^*(\omega \wedge\tau) = F^*\omega \wedge F^*\tau$$
\end{proposition}

\begin{definition}[Space of Forms]
  \

  For a manifold $M^m$, write $$\Omega^*(M) = \bigoplus_{i=0}^m \Omega^i(M)$$
\end{definition}

\subsection{Differential Forms on a Circle}
\

\begin{proposition}
  \

  Set $h: \mathbb{R} \to S^1$ by $h(t) = (\cos t, \sin t)$. For $k = 0,1$, under the pullback map $h^*: \Omega^*(S^1) \to \Omega^*(\mathbb{R})$, smooth $k$-forms on $S^1$ are identified with smooth periodic $k$-form of period $2\pi$ on $\mathbb{R}$.
\end{proposition}

\begin{customproof}
  \

  We first show that being a smooth $0$ form on $S^1$ makes you a $2\pi$ periodic form on $\mathbb{R}$.

  If $f \in \Omega^0(S^1)$, then $h^*f = f \circ h$, but since $h$ is $2\pi$ periodic, so is $f \circ h$.

  Next, show that every $2\pi$ periodic form on $\mathbb{R}$ is a form on $S^1$.

  Let $\bar{f} \in \Omega^0(\mathbb{R})$ be $2\pi$ periodic. Let $s_1$ and $s_2$ be two inverses of $h$, so $s_1 = s_2 + 2\pi n$ for $n \in \mathbb{Z}$. Since $\bar{f}$ is $2\pi$ periodic, $\bar{f} \circ s_1 = \bar{f} \circ s_2$, and set $f = \bar{f} \circ s_1$. Then, $\bar{f} = f \circ s ^{-1} = f \circ h = h^*f$, so the space of $2\pi$ periodic forms is prescisely the pullback of all smooth forms on $S^1$.

  The arguments generalize to $1$-forms, since $1$-forms are just $0$ forms with a $dt$ appended.
\end{customproof}

\subsection{Invariant Forms on a Lie Group}
\

\begin{definition}[Left Invariance for $k$-forms]
  \

  If $G$ is a Lie group, $\omega$ is a $k$-form on $G$, $\omega$ is \textbf{left invariant} if for every $g \in G$, $$\ell_g^*(\omega_{gx}) = \omega_x$$
\end{definition}

\begin{proposition}
  \

  Every left invarnant invariant $k$-form on a Lie group is smooth.
\end{proposition}

\begin{customproof}
  \

  Let $(Y_1)_e, \dots, (Y_n)_e$ be a basis for the tangent space $T_eG$, and $Y_1, \dots, Y_n$ be the vector fields they generate. Then, since every vector in $T_eG$ is a linear combination of $(Y_i)_e$, every left invariant vector field $X$ is a linear combination of $Y_i$. Hence, by linearity stuff, it is sufficient to show that $\omega \in \bigwedge^k(G)$ is smooth on $Y_i$.

  \begin{align*}
    \omega(Y_{i_1}, \dots, Y_{i_k})(g) &= \omega_g((Y_{i_1})_g, \dots (Y_{i_k})_g)\\
    &= (\ell^*_{g ^{-1}}(\omega_e))(\ell_{g*}(Y_{i_1})_e, \dots, \ell_{g*}(Y_{i_k})_e)\\
    &= (\ell_g \circ \ell_{g ^{-1}} )^*(\omega_e((Y_{i_1})_e, \dots, (Y_{i_k})_e))\\
    &= \omega_e((Y_{i_1})_e, \dots, (Y_{i_k})_e)
  \end{align*}But this shows that $\omega_e((Y_{i_1})_e, \dots, (Y_{i_k})_e)$ is a constant function, so it is smooth.

\end{customproof}

\section{The Exterior Derivative}

\subsection{The Exterior Derivative on a Coordinate Chart}
\

\begin{remark}
  \

  Recall that an \textbf{antiderivation} on a graded algebra $A = \oplus_{k=0}^\infty A^k$ is a linear map $D: A \to A$ such that $D(\omega \cdot \tau) = (D\omega) \cdot \tau + (-1)^k\omega \cdot D\tau$ for $\omega \in A^k$, $\tau \in A^l$. An element of $A^k$ is called a \textbf{homogeneous element of degree $k$}, and the antiderivation is \textbf{of degree} $m$ if $\text{deg}D\omega = \deg \omega = m$.

  $\Omega^*(M)$ is a graded algebra of smooth differential forms, and has a natural antiderivation called the \textbf{exterior derivative}.
\end{remark}

\begin{definition}[Exterior Derivative]
  \

  An \textbf{exterior derivative} on a manifold $M$ is an $\mathbb{R}$-linear map $$D: \Omega^*(M) \to \Omega^*(M)$$such that
  \begin{enumerate}
    \item $D$ is an antiderivation of degree $1$
    \item $D \circ D = 0$
    \item if $f$ is a smooth function and $X$ is a smooth vector field, $(Df)(X) = Xf$ (the exterior derivative of a $0$ form is a differential)
  \end{enumerate}

\end{definition}

\begin{proposition}
  \

  On any manifold $M$ with chart $(U,x^i)$ coordinates, $\omega = \sum a_Idx^I$ a $k$-form, the exterior derivative exists in the chart, is unique, and is given by $$D\omega = \sum_I\sum_j \frac{\partial a_I}{\partial x^j} dx^j \wedge dx^I$$
\end{proposition}

\begin{customproof}
  \

  \begin{align*}
    D\omega &= \sum (Da_I) \wedge dx^I + \sum a_I Ddx^I &&\text{by linearity and the first property}\\
    &= \sum(Da_I) \wedge dx^I && \text{by $Dd(x^I) = D^2x^I = 0$}\\
    &= \sum_I \sum_j \frac{\partial a_I}{\partial x^j}dx^j \wedge dx^I && \text{the derivation of a function}
  \end{align*}And as we see, the exterior derivative is uniquely defined, and exists for any manifold $M$.
\end{customproof}

\subsection{Local Operators}
\

\begin{definition}[Operators]
  \

  An endomorphism of a vector space $W$ is called an \textbf{operator} on $W$. An operator $D: \Omega^*(M) \to \Omega^*(M)$ is \textbf{local} if for all $k \geq 0$, whenever a $k$-form $\omega \in \Omega^k(M)$ is such that there is an open set $U$ where for all $p \in U$, $\omega_p = 0$, then $D\omega = 0$ on $U$.
\end{definition}

\begin{proposition}
  \

  Any antiderivation $D$ on $\Omega^*(M)$ is a local operator.
\end{proposition}

\begin{customproof}
  \

  Let $\omega \in \Omega^k(M)$ and $\omega = 0$ on an open subset $U$, choose some $p\in U$. Then, choose a smooth bump function $f$ at $p$ supported in $U$, so $f\omega = 0$ on $M$. Then, $$0 = D(0) = (Df\omega) = (Df) \wedge \omega + f\wedge D\omega$$and at $p$, $f(p) = 1$, so $(D\omega)_p = 0$
\end{customproof}

\subsection{Existence of an Exterior Derivative on a Manifold}
\

\begin{proposition}
  \

  An exterior derivative exists and is unique on manifolds.
\end{proposition}

\begin{customproof}
  \

  We know that for a manifold $M$, a chart $(U,x^i)$, form $\omega = \sum a_Idx^I$, there exists a unique exterior derivative $$d_U\omega = \sum da_I \wedge dx^I$$If $(Y, y^i)$ is another chart, then on $U \cap V$, let $\omega = \sum a_Idx^I = \sum b_Jdy^J$. Since $U \cap V$ is another chart, there is a unique exterior derivative $d_{U \cap V}$, and $$d_{U \cap V}(\sum a_I dx^I) = d_{U \cap V}(\sum b_J dy^J)$$so $$(\sum da_I\wedge dx^I)_p = (\sum db_J \wedge dy^J)$$so they do agree on intersections, and so there is indeed at least an exterior derivative $d: \Omega^*(M) \to \Omega^*(M)$ on a manifold.
\end{customproof}

\subsection{Uniqueness of the Exterior Derivative}
\

\begin{theorem}
  \

  On any manifold $M$ there exists a unique exterior derivative $d: \Omega^*(M) \to \Omega^*(M)$.
\end{theorem}

\begin{customproof}
  \

  We have already shown existence, so we want to show that if $D: \Omega^*(M) \to \Omega^*(M)$ is a derivative, it must be the exterior derivative defined above.

  By definition, $$(Df)(X) = Xf = (df)(X)$$so acting on $0$-forms, they are identical.

  Next, consider a $k$-form $\omega \in \Omega^k(M)$. Fix $p \in M$, choose a chart $(U, x^i)$ about $p$, and let $\omega = \sum a_Idx^I$. Extend each $a_I$, each $x^i$ to smooth functions $\tilde{a_I}, \tilde{x^i}$ via a bump functions, so that they agree on some $p \in V \subset U$. Then, let $\tilde{\omega} = \sum \tilde{a_I}d\tilde{x}^I \in \Omega^k(M)$, and since $D$ is a local operator, $D\omega = D\tilde{\omega}$. Thus: $$(D\omega)_p = (D\tilde{\omega})_p = (\sum D\tilde{a}_I \wedge d\tilde{x}^I + \sum \tilde{a}_I \wedge Dd\tilde{x}^I) = (\sum d\tilde{a}_I \wedge d\tilde{x}^I) = (d\omega)_p$$
\end{customproof}

\subsection{Exterior Differentiation Under a Pullback}
\

\begin{proposition}
  \

  Let $F: M \to N$ be a smooth map of manifolds. If $\omega \in \Omega^k(M)$, then $dF^*\omega = F^*d\omega$.
\end{proposition}

\begin{customproof}
  \

  When $k = 0$, we have already proved this. Consider $k\geq1$, fix some $p \in M$. Then, $$F^* \omega = \sum (F^*a_I)dF^{i_1} \wedge \dots \wedge dF^{i_k}$$so $$dF^*\omega = \sum d(a_I \circ F) \wedge dF^{i_1} \wedge \dots \wedge dF^{i_k}$$

  From the other side, we have $$
  F^* d\omega = F^*(\sum da_I \wedge dy^I) = \sum F^*da_I \wedge dF^{i_1} \wedge \dots \wedge dF^{i_k} = \sum d(a_I \circ F)\wedge dF^{i_1} \wedge \dots \wedge dF^{i_k}$$And we have equality.
\end{customproof}

\begin{corollary}
  \

  If $U$ is an open subset of $M$, then $$(d\omega) \vert_U = d(\omega \vert_U)$$
\end{corollary}

\begin{customproof}
  \

  Choose $F = \iota$, the inclusion map, in the prior proof
\end{customproof}

\begin{proposition}
  \

  If $F: M\to N$ is a smooth map between smooth manifolds, $\omega$ is a smooth $k$-form on $N$, then $F^*\omega$ is a smooth $k$-form on $M$
\end{proposition}

\begin{customproof}
  \

  Let $(V, y^i)$ be a chart in $N$, $\omega = \sum_I a_I dy^I$, $(U, x^i)$ be a chart in $M$ such that $F(U) \subset V$. Then,

  $$F^*\omega = \sum (a_I \circ F) dF^{i_1} \wedge \dots \wedge dF^{i_k} = \sum_{I,J}(a_I \circ F) \frac{\partial F^{i_1}, \dots, F^{i_k}}{\partial x^{j_1}, \dots x^{j_k}} dx^J$$and for $a_I$ smooth, $F$ smooth, each otf these is smooth as well.
\end{customproof}

\subsection{Restriction of $k$-Forms to a Submanifold}
\

\begin{definition}[Restriction of $k$-Forms]
  \

  The \textbf{restriction} of $\omega \in \Omega^k(M)$ to $S \subseteq M$ is given by $$(\omega \vert_S)_p(v_1, \dots, v_k) = \omega_p(v_1, \dots, v_k)$$for $v_1, \dots, v_k \in T_pS$

  A nonzero form can restrict to a zero form on a submanifold $S$, if the form happens to be zero on that submanifold.

  A \textbf{nowhere zero} or \textbf{nowhere vanishing} form is nowhere zero.
\end{definition}

\subsection{A Nowhere-Vanishing $1$-form on the Circle}
\

\begin{example}
  \

  Consider $S^1\subset \mathbb{R}^2$. The $1$-form $dx$ restricts from $\mathbb{R}^2$ to $S^1$, but at $p = (1,0)$, the basis for $T_p(S^1)$ is $\frac{\partial}{\partial y}$, and we have that $$(dx)_p(\frac{\partial}{\partial y}) = (x dx)(\frac{\partial}{\partial y}) = 0$$Even though $dx$ is nowhere vanishing on $\mathbb{R}^2$, it vanishes at $(1,0)$ when restricted to $S^1$.

  To find a nowhere vanishing form on $S^1$, take the exterior derivative of both sides of the equation $$x^2 + y^2 = 1$$to get $$2xdx + 2ydy = 0$$When both $x,y$ nonzero, we get $$\frac{dy}{x} = -\frac{dx}{y}$$and we can define a $1$-form as $$\omega =
  \begin{cases}
    \frac{dy}{x} & x \neq 0\\
    -\frac{dx}{y} & y \neq 0
  \end{cases}$$when both nonzero, they agree, and this is defined on all $S^1$, and is nowhere vanishing.
\end{example}

\section{The Lie Derivative and Interior Multiplication}

\subsection{Families of Vector Fields and Differential Forms}
\

\begin{definition}[$1$-Parameter Families]
  \

  A collection $\{X_t\}$ or $\{\omega_t\}$ of vector fields or differential forms is a \textbf{$1$-parameter family} if $t \in I \subseteq \mathbb{R}$.

  On a coordinate neighborhood $(U, x^i)$, we denote the \textbf{limit} of $\{X_t\}$ by $$\lim_{t \to t_0} X_t \vert_p = \sum_{i=1}^n \lim_{t \to t_0}a^i(t,p) \left.\frac{\partial}{\partial x^i}\right\vert_p$$

  A $1$-parameter family \textbf{depends smoothly} on $t$ if for all $p \in M$, there is a coordinate neighborhood $(U, x^i)$ where $$(X_t)_p = \sum a^i(t,p) \left.\frac{\partial}{\partial x^i}\right\vert_p$$for smooth functions $a^i$ on $I \times U$, and we call this $\{X_t\}$ a \textbf{smooth family of vector fields}.

  Given a smooth family of vector fields, we can define the derivative $$(\left.\frac{d}{dt}\right\vert_{t=t_0} X_t) = \sum \frac{\partial a^i}{\partial t}(t_0,p) \left.\frac{\partial}{\partial x^i}\right\vert_p$$And much of the same goes for $k$-forms.
\end{definition}

\begin{proposition}
  \

  If $\{\omega_t\}$ and $\{\tau_t\}$ are smooth families of $k$-forms and $l$-forms respectively on a manifold $M$, then $$\frac{d}{dt}(\omega_t \wedge \tau_t) = (\frac{d}{dt}\omega_t) \wedge \tau_t + \omega_t \wedge \frac{d}{dt} \tau_t$$
\end{proposition}

\begin{customproof}
  \

  For any $t_0 \in I$, $p \in M$, we have $$(\left.\frac{d}{dt}\right\vert_{t=t_0} \omega_t \wedge \tau_t)_p = \sum \frac{\partial}{\partial t}(a_Ib_J(t_0,p)) dx^I \wedge dx^J \vert_p =$$$$ \sum \frac{\partial a_I}{\partial t}b_J dx^I \wedge dx^J \vert_p + \sum \frac{\partial b_J}{\partial t}a_I dx^I\wedge dx^J \vert_p = ((\left.\frac{d}{dt}\right\vert_{t=t_0} \omega_t)\wedge \tau_t + \omega_t \wedge (\left.\frac{d}{dt}\right\vert_{t=t_0}\tau_t))_p$$
\end{customproof}

\begin{proposition}
  \

  If $\{\omega_t\}_{t \in I}$ is a smooth family of differential forms on a manifold $M$ then $$\left.\frac{d}{dt}\right\vert_{t=t_0} d\omega_t = d(\left.\frac{d}{dt}\right\vert_{t=t_0} \omega_t)$$
\end{proposition}

\begin{customproof}
  \

  Independent of evaluating at $t = t_0$, we just show that exterior derivative and the time derivative commute: $$\frac{d}{dt}(d\omega_t) = \frac{d}{dt} \sum_{J,i} \frac{\partial b_J}{\partial x^i} dx^i \wedge dx^J = \sum_{i,J} \frac{\partial}{\partial x^i} (\frac{\partial b_J}{\partial t}) dx^i \wedge dx^J = d(\sum_J \frac{\partial b_J}{\partial t}dx^J) = d(\frac{d}{dt}\omega_t)$$

  And evaluation commutes as well, because $d$ doesn't involve any time derivative.
\end{customproof}

\subsection{The Lie Derivative of a Vector Field}
\

\begin{remark}
  \

Recall that given a smooth vector field $X$, $p \in M$, there is a neighborhood $U$ of $p$ where there is a \textbf{local flow}, a real number $\varepsilon > 0$ and a map $\varphi: ]-\varepsilon, \varepsilon[ \times U \to M$ such that $\varphi_t(q) = \varphi(t,q)$ and $\frac{\partial}{\partial t}\varphi(t,q) = \frac{\partial}{\partial t}\varphi_t(q) = X_{\varphi_t(q)}$, and $\varphi_0(q) = 0$. Basically, $\varphi_t(q)$ for fixed $q$, variable $t$, is an integral curve with initial point $q$.

  Local flows have the properties: $\varphi_s \circ \varphi_t = \varphi_{s+t}$ and they are also diffeomorphisms.

  If $Y$ is a smooth vector field, we can use the pushforward of $\varphi_{-t}: \varphi_t(U) \to U$ to bring $Y_{\varphi_t(p)} \in T_{\varphi_t(p)}M$ into $\varphi_{-t*}(Y_{\varphi_t(p)}) \in T_pM$
\end{remark}

\begin{definition}[Lie Derivative]
  \

  For $X,Y$ smooth vector fields on $M$, define the \textbf{Lie derivative} $\mathcal{L}_XY$ to be $$(\mathcal{L}_XY)_p = \lim_{t \to 0} \frac{\varphi_{-t*}(Y_{\varphi_t(p)}) - Y_p}{t} =  \lim_{t\to0} \frac{(\varphi_{-t*}Y)_p - Y_p}{t} = \left.\frac{d}{dt}\right\vert_{t=0} (\varphi_{-t*}Y)_p$$On local coordinates, $(U, x^i)$, $$(\mathcal{L}_XY)_p = \left.\frac{d}{dt}\right\vert_{t=0} \varphi_{-t*}(Y_{\varphi_t(p)}) = \sum_{i,j} \left.\frac{\partial}{\partial t}\right\vert_{t=0} (b^j(\varphi(t,p))\frac{\partial \varphi^i}{\partial x^j}(-t,p))\left.\frac{\partial}{\partial x^i}\right\vert_p $$
\end{definition}

\begin{remark}
  \

  $\varphi_{-t*}$ moves $Y_{\varphi_t(p)}$ back along $X$ (or an integral curve of $X$) into $T_pM$, and the Lie derivative compares that with $Y_p$ which may or may not be different.
\end{remark}

\begin{theorem}
  \

  The Lie derivative (as defined above), and the Lie bracket are the same.
\end{theorem}

\begin{customproof}
  \

Let $X,Y$ be smooth vector fields, let a local flow for $X$ be $\varphi: ]-\varepsilon, \varepsilon[ \times U \to M$, where $(U, x^i)$ is a coordinate chart. Let $X = \sum a^i \frac{\partial}{\partial x^i}$, $Y = b^j \frac{\partial}{\partial x^j}$ on $U$.

  Since $\varphi_t(p)$ is an integral curve, $\frac{\partial \varphi^i(t,p)}{\partial t}(t,p) = \frac{\partial x^i \circ \varphi}{\partial t}(t,p) = a^i(\varphi(t,p))$, and at $t=0$, $\frac{\partial \varphi^i}{\partial t}(0,p) = a^i(p)$

  We know the coordinates for the Lie bracket are given by $$[X,Y] = \sum_{i,k}(a^k \frac{\partial b^i}{\partial x^k} - b^k \frac{\partial a^i}{\partial x^k})\frac{\partial}{\partial x^i}$$

  The Lie derivative gies us $$(\mathcal{L}_XY)_p = \sum_{i,j,k}(\frac{\partial b^j}{\partial x^k}(p)a^k(p)\frac{\partial \varphi^i}{\partial x^j}(0,p))\frac{\partial}{\partial x^i} - \sum_{i,j} (b^j(p)\frac{\partial a^i}{\partial x^j}(p))\frac{\partial}{\partial x^i} = \sum_{i,k}(a^k \frac{\partial b^i}{\partial x^k} - b^k \frac{\partial a^i}{\partial x^k})\frac{\partial}{\partial x^i}$$Hence the two are equal.
\end{customproof}

\subsection{The Lie Derivative on a Differential Form}
\

\begin{definition}[Lie Derivative]
  \

  For $X$ a smooth vector field, $\omega$ a smooth $k$-form on a manifold $M$, the \textbf{Lie derivative} $\mathcal{L}_X\omega$ at $p \in M$ is $$(\mathcal{L}_X\omega)_p = \lim_{t \to 0} \frac{\varphi_t^*(\omega_{\varphi_t(p)}) - \omega_p}{t} = \lim_{t \to 0} \frac{(\varphi_t^*)_p - \omega_p}{t} = \left.\frac{d}{dt}\right\vert_{t=0}(\varphi_t^*\omega)_p$$
\end{definition}

\begin{proposition}
  \

  If $f$ is a smooth function and $X$ a smooth vector field, then $$\mathcal{L}_Xf = Xf$$
\end{proposition}

\begin{customproof}
  \

  \begin{align*}
    \mathcal{L}_Xf &= \left.\frac{d}{dt}\right\vert_p (\varphi_t^* f)_p \\
    &= \left.\frac{d}{dt}\right\vert_p (f \circ \varphi_t)(p)\\
    &= X_pf && \text{by a proposition in chapter 8}
  \end{align*}
\end{customproof}

\subsection{Interior Multiplication}
\

\begin{definition}[Interior Multiplication]
  \

  If $\beta$ is a $k$-covector on a vector space $V$, $v \in V$, for $k \geq 2$, the \textbf{interior multiplication} or \textbf{contraction} of $\beta$ with $v$ is the $(k-1)$-covector $\iota_v\beta$, given by $$(\iota_v\beta)(v_2, \dots, v_k) = \beta(v, v_2, \dots, v_k)$$
\end{definition}

\begin{proposition}
  \

  For $1$-covectors $\alpha^1, \dots, \alpha^k$ on a vector space $V$, and $v \in V$, $$\iota_v(\alpha^1 \wedge \dots \wedge \alpha^k) = \sum_{i=1}^k (-1)^{i-1} \alpha^i(v)\alpha^1\wedge \dots \wedge \hat{\alpha^i} \wedge \dots \wedge \alpha^k$$where the hat means that the $\alpha^i$ is omitted.
\end{proposition}

\begin{customproof}
  \

  Let $v_i \in V$, then
  \begin{align*}
    \iota_v(\alpha^1 \wedge \dots \wedge \alpha^k)(v_2, \dots, v_k) &= (\alpha^1 \wedge \dots \wedge \alpha^k)(v, v_2, \dots, v_k) \\
    &= \det
    \begin{bmatrix}
      a^1(v) &  \dots& a^1(v_k)\\
      \vdots & \ddots & \vdots\\
      a^k(v) & \dots & a^k(v_k)
    \end{bmatrix} \\
    &= \sum_{i=1}^k (-1)^{i-1} \alpha^i(v)(\alpha^1\wedge \dots \wedge \hat{\alpha^i} \wedge \dots \wedge \alpha^k)(v_2, \dots, v_k)
  \end{align*}
\end{customproof}

\begin{proposition}
  \

  For $v$ in a vector space $V$, let $\iota_v: \bigwedge^*(V^*) \to \bigwedge^{*-1} (V^*)$ be interior multiplication by $v$. Then,
  \begin{enumerate}
    \item $\iota_v \circ \iota_v = 0$
    \item for $\beta \in \bigwedge^k(V^*)$ and $\gamma \in \bigwedge^l(V^*)$, $$\iota_v(\beta \wedge \lambda) = (\iota_v \beta) \wedge \lambda + (-1)^k \beta \wedge \iota_v\lambda$$
  \end{enumerate}

  So $\iota_v$ is an antiderivation of degree $-1$, and applying it twice yields zero.
\end{proposition}

\begin{customproof}
  \

  \begin{enumerate}
    \item For any $\beta \in \bigwedge^k(V^*)$, we have $$\iota_v \circ \iota(v)(\beta)(v_3, \dots, v_k) = \beta(v,v,v_3, \dots, v_k) = 0$$Since $\beta$ is alternating and there is a repeated variable in its arguments
    \item Since both sides are linear, assume WLOG $\beta = \alpha^1 \wedge \dots \wedge \alpha^k$ and $\gamma = \alpha^{k+1} \wedge \dots \wedge \alpha^{k+l}$. Then,
      \begin{align*}
        \iota_v(\beta \wedge \gamma) &= \iota_v(\alpha^1\wedge\dots\wedge\alpha^{k+l})\\
        &= (\sum_{i=1}^k (-1)^{i-1} \alpha^i(v)\alpha^1 \wedge \dots \wedge \hat{\alpha^i}\wedge\dots\wedge \alpha^k) \wedge \alpha^{k+1} \wedge \dots \wedge \alpha^{k+l} \\ &+(-1)^k \alpha^1 \wedge \dots \wedge \alpha^k \wedge \sum_{i=1}^k (-1)^{i+1} \alpha^{k+i}(v)\alpha^{k+1}\wedge \dots \wedge \hat{\alpha^{k+i}}\wedge \dots \wedge \alpha^{k+l} \\ &=(\iota_v\beta)\wedge \gamma + (-1)^k\beta\wedge\iota_v\gamma
      \end{align*}
  \end{enumerate}
\end{customproof}

\begin{definition}[Interior Multiplication]
  \

  On manifolds, interior multiplication is defined pointwise: $$(\iota_X \omega)_p = \iota_{X_p}\omega_p$$
\end{definition}

\subsection{Properties of the Lie Derivative}
\

\begin{theorem}
  \

  Assume $X$ is a smooth vector field on a manifold $M$.
  \begin{enumerate}
    \item $\mathcal{L}_X: \Omega^*(M) \to \Omega^*(M)$ is a derivation
    \item The Lie derivative commutes with the external derivative
    \item $\mathcal{L}_X = d\iota_X + \iota_Xd$
    \item $$\mathcal{L}_X(\omega(Y_1,\dots, Y_k)) = (\mathcal{L}_X\omega)(Y_1,\dots, Y_k) + \sum_{i=1}^k \omega(Y_1, \dots, \mathcal{L}_XY_i, \dots, Y_k)$$
  \end{enumerate}
\end{theorem}

\subsection{Global Formulas for the Lie and Exterior Derivatives}
\

\begin{theorem}
  \

  For a smooth $k$-form $\omega$, and smooth vector fields $X, Y_1, \dots, Y_k$ on a manifold $M$, $$(\mathcal{L}_X\omega)(Y_1,\dots,Y_k) = X(\omega(Y_1,\dots,Y_k)) - \sum_{i=1}^k\omega(Y_1, \dots, [X, Y_i], \dots, Y_k)$$
\end{theorem}

\begin{proposition}
  \

  If $\omega$ is a smooth $1$-form and $X$, $Y$ are smooth vector fields on a manifold $M$, then $$d\omega(X,Y) = X\omega(Y) - Y\omega(X) - \omega([X,Y])$$
\end{proposition}

\begin{theorem}
  \

  If $k\geq 1$, $\omega$ a smooth $k$-form, and $Y_0,\dots, Y_k$ smooth vector fields on a manifold $M$, $$(d\omega)(Y_0, \dots, Y_k) = \sum_{i=0}^k(-1)^iY_i\omega(Y_0, \dots, \hat{Y_i}, \dots, Y_k) + \sum_{0 \leq i < j \leq k} (-1)^{i+j} \omega([Y_i, Y_j], Y_0, \dots, \hat{Y_i}, \dots, \hat{Y_j}, \dots, Y_k)$$
\end{theorem}

\section{Orientations}

\subsection{Orientations of a Vector Space}
\

\begin{definition}[Orientations]
  \

  Let $V$ be a finite-dimensional real vector space. An \textbf{orientation} on $V$ is an equivalence class of ordered bases, where two bases, $(v_1, \dots, v_n)$ and $(u_1, \dots, u_n)$ are related if $$(v_1, \dots, v_n) = (u_1,\dots, u_n)A$$and $\det A > 0$, where $A$ is the unique transition matrix between the bases.
\end{definition}

\begin{example}
  \

  $\mathbb{R}^3$ admits two orientations: $(e_1,e_2,e_3)$ and $(e_2,e_1,e_3)$.
\end{example}

\subsection{Orientations and $n$-covectors}
\

\begin{lemma}
  \

  Let $u_1, \dots, u_n$ and $v_1, \dots, v_n$ be vectors in a vector space $V$, let $u_j = \sum_{i=1}^n v_i a^i_j$ so that a matrix $A = [a^i_j]$ transitions betwen them. If $\beta$ is an $n$-covector on $V$, then $$\beta(u_1, \dots, u_n) = (\det A)\beta(v_1, \dots, v_n)$$
\end{lemma}

\begin{customproof}
  \

  $$\beta(u_1, \dots, u_n) = \beta(\sum_I v_{i_1} a_1^{i_1}, \dots, \sum_I v_{i_n}a^{i_n}_n) = \sum a_1^{i_1} \dots a_n^{i_n}\beta(v_{i_1} \dots, v_{i_n}) = \sum a_1^{i_1} \dots a_n^{i_n}\text{sgn} \sigma_I\beta(v_{1} \dots, v_{n})$$$$= (\det A)\beta(v_1, \dots, v_n)$$
\end{customproof}

\begin{corollary}
  \

  The ordered bases $(u_1, \dots, u_n)$ and $(v_1, \dots, v_n)$ are equivalent iff for any $n$-form $\beta$, $\beta(u_1, \dots, u_n)$ and $\beta(v_1, \dots, v_n)$ have the same sign.
\end{corollary}

\begin{definition}[Equivalence of Covectors]
  \

  Given an ordered basis $(v_1, \dots, v_n)$, we say that an $n$-covector $\beta$ \textbf{determines} the orientation if $\beta(v_1, \dots, v_n) > 0$. We can then construct an equivalence class of covectors, by saying two $n$-covectors $\beta, \beta'$ are related iff $\beta = a \beta'$ for some positive real number $a$.
\end{definition}

\subsection{Orientations on a Manifold}
\

\begin{definition}[Local \& Global Frames]
  \

  A \textbf{frame} on an open subset $U \subset M^m$ of a manifold $M$ is an ordered $n$-tuple of (possibly discontinuous) vector fields $(X_1,\dots, X_m)$ such that at each $p \in M$, $(X_{1,p}, \dots, X_{m,p})$ forms a basis for $T_pM$. A \textbf{global frame} is defined for all $M$, whereas a \textbf{local frame} is defined for some neighborhood about $p$.

  We say two frames $(X^1, \dots, X_m)$, $(Y_1, \dots, Y_m)$ are related if for all $p \in M$, $(X_{1,p}, \dots, X_{m,p}) \sim (Y_{1,p}, \dots, Y_{m,p})$ (as ordered bases of the vector space $T_pM$)
\end{definition}

\begin{definition}[Orientations on Manifolds]
  \

  A \textbf{pointwise orientation} on a manifold $M$ assigns each $p \in M$ an orientation $\mu_p$ of $T_pM$. It is then, an equivalence class of possibly discontinuous frames on $M$.

  A pointwise orientation $\mu$ is \textbf{continuous} at $p \in M$ if $p$ has a neighborhood $U$ on which $\mu$ can be represented by some continuous frame.

  A pointwise orientation which is continuous everywhere is called a \textbf{orientation} on $M$. If a manifold $M$ has an orientation, it is \textbf{orientable}.
\end{definition}

\begin{example}
  \

  $\mathbb{R}^n$ has orientation $(\frac{\partial}{\partial r^1}, \dots, \partial{\partial r^n})$
\end{example}

\begin{proposition}
  \

  A connected orientable manifold $M$ has exactly two orientations
\end{proposition}

\begin{customproof}
  \

  Let $\mu$, $\nu$ be two orientations on $M$. At any $p \in M$, $\mu_p$ and $\nu_p$ are orientations of $T_pM$, and since every vector space admits exactly two orientations, they can either be the same or opposite.

  Let $$f(p) =
  \begin{cases}
    1 & \mu_p = \nu_p\\
    -1 & \mu_p \neq \nu_p
  \end{cases}$$Fix a point $p \in M$. Since both $\mu$, $\nu$ are continuous, on some neighborhood $p \in U$, they are represented by frames $\mu \sim (X_1, \dots, X_m)$ and $\nu \sim (Y_1, \dots, Y_m)$ respectively, and these frames are continuous (each $X_i$ and $ Y_j$ are smooth vector fields). But then, there is a function $A: U \to \text{GL}(n, \mathbb{R})$ which takes a point $p \in U$, and outputs $A(p)$ is the transition matrix between $(X_{1,p}, \dots, X_{m,p})$ and $(Y_{1,p}, \dots, Y_{m,p})$ as bases of $T_pM$.

  Each coordinate function in $A(p) = [a^i_j(p)]$ is also continuous, by a prior proposition, so $\det A: U \to \mathbb{R}$ is continuous as well. Then, by IVT, since $\det A \neq 0$ (since it outputs a matrix in $\text{GL}(n, \mathbb{R})$), it must either be all positive or all negative, so either $\mu = \nu$ or $\mu = -\nu$ on $U$, so $f$ is locally constant, and since a locally constant function on a connected set is constant, $\mu = \nu$ or $\mu = -\nu$ on $M$, so there are exactly two orientations.
\end{customproof}

\subsection{Orientations and Differential Functions}
\

\begin{lemma}
  \

  A pointwise orientation represented by $\mu \sim (X_1, \dots, X_m)$ is continuous iff at each point $p \in M$, there is a coordinate neighborhood $(U, x^i)$ where $(dx^1 \wedge \dots \wedge dx^m)(X_1, \dots, X_m)$ is everywhere positive.
\end{lemma}

\begin{customproof}
  \

  ($\implies$) Assume that $\mu \sim (X_1, \dots, X_m)$ is a continuous orientation, meaning that for all $p \in M$, there is some $p \in W \subset M$ where $\mu$ is represented by a continuous frame $(Y_1, \dots, Y_m)$. Choose a connected coordinate neighborhood $(U, x^i)$, $U \subset W$. And let $[b^i_j(p)]$, for fixed $p$, be a change of bases matrix between $\partial_i$ and $Y_i$. Then,

  $$(dx^1 \wedge \dots \wedge dx^m)(Y_1,\dots,Y_m) = (\det[b^i_j(p)])(dx^1 \wedge \dots \wedge dx^m)(\partial_1, \dots, \partial_m) = \det[b^i_j(p)]$$and $\det[b^i_j(p)]$ is a continuous, never zero function, so it is either everywhere positive or everywhere negative. If it is everywhere negative, flip $dx^1$ and you get coordinates where it is everwhere positive.

  Now if you let $[c^i_j(p)]$ be the change of basis matrix between $X_i$ and $Y_i$, we get that $$(dx^1 \wedge \dots \wedge dx^m)(X_1, \dots, X_m) = (\det[c^i_j(p)])(dx^1 \wedge \dots \wedge dx^n)(Y_1, \dots, Y_n) > 0$$

  ($\impliedby$) Let $[a^i_j(p)]$ change between $\partial_i$ and $X_i$. Then, $$(dx^1 \wedge \dots\wedge dx^n)(X_1, \dots, X_n) = (\det[a^i_j(p)])(dx^1 \wedge \dots \wedge dx^n)(\partial_1, \dots, \partial_n) = \det[a^i_j(p)]$$Since the left is positive, the right is positive, and so the orientation represented by $X_i$ is equivalent to that of $\partial_i$ locally, and $\partial_i$ is continous, so so is $X_i$, (on chart $U$), and since this is true for all $p \in M$, it is true globally.
\end{customproof}

\begin{theorem}
  \

  A manifold $M^m$ is orientable iff there exists a smooth nowhere vanishing $m$-form on $M$
\end{theorem}

\begin{customproof}
  \

  $(\implies)$ Suppose $[(X_1, \dots, X_m)]$ is an orientation on $M$. Then, for each $p \in M$, there is a chart $(U, x^i)$ where $(dx^1 \wedge \dots \wedge dx^m)(X_1, \dots, X_m) > 0$. Let $\{U_\alpha, x^i_\alpha\}$ be a collection of such charts that covers $M$, and let $\{\rho_\alpha\}$ a partition subordinate to $\{U_\alpha\}$. Since each $\rho_\alpha(p) \geq 0$, and for any fixed $p$, at least one $\alpha$ has $\rho_\alpha(p) > 0$, $\omega = \sum_{\alpha} \rho_\alpha dx^1 \wedge \dots \wedge dx^m$ is a nowhere zero $m$-form.

  $(\impliedby)$ Let $\omega$ be a nowhere vanishing $m$-form on $M$. For each $p \in M$, let $(X_{1,p}, \dots, X_{m,p})$ be an ordered basis for $T_pM$ such that $\omega_p(X_{1,p}, \dots, X_{m,p}) >0$, then, $\omega = f dx^1 \wedge \dots \wedge dx^m$, and since $f$ continuous, nowhere vanishing, it is either all positive or all negative. If $f > 0$, then $(dx^1, \dots, dx^m)(X_1, \dots, X_m) > 0$, else just flip the sign of $x^1$. In either case, $(X_1, \dots, X_m)$ is an orientation.
\end{customproof}

\begin{definition}[Oriented Manifolds]
  \

  If $\omega$ is a nowhere vanishing $m$ form on $M$, then $\omega$ \textbf{determines} the orientation $[(X_1, \dots, X_m)]$ if $\omega(X_1, \dots, X_n) > 0$, and orientations are associated with equivalence classes of nowhere vanishing $m$ forms by $\omega \sim \omega'$ if they both determine the same orientation.

  A tuple of $(M, [\omega])$ is an \textbf{oriented manifold}, where $[\omega]$ is an orientation, $\omega$ a nowhere vanishing $m$-form on $M$

  A diffeomorphism $F: (M, [\omega_M]) \to (N, [\omega_N])$ is \textbf{orientation preserving} if $[F^*\omega_N] = [\omega_M]$
\end{definition}

\begin{proposition}
  \

  Let $U,V$ be open subsets of $\mathbb{R}^n$, with standard orientation inherited from $\mathbb{R}^n$. A diffeomorphism is orientation preserving iff $\det[\frac{\partial F^i}{\partial x^j}]$ is everywhere positive on $U$.
\end{proposition}

\begin{customproof}
  \

  Let $x^i$, $y^i$ be the standard coordinates on $U, V$ respectively. Then,
  \begin{align*}
    F^*(dy^1 \wedge \dots \wedge d^n) &= d(F^*y^1) \wedge \dots \wedge d(F^*y^n)\\
    &= d(y^1 \circ F) \wedge \circ \wedge d(y^n \circ F)\\
    &= dF^1 \wedge \dots \wedge dF^n\\
    &= \text{det}[\frac{\partial F^i}{\partial x^j}] dx^1 \wedge \dots \wedge dx^n
  \end{align*}And so acting on frames, we can see that this is orientation preserving iff the Jacobian is everywhere positive.
\end{customproof}

\subsection{Orientations and Atlases}
\

\begin{definition}[Oriented Atlases]
  \

  An atlas on $M$ is \textbf{oriented} if for any two overlapping charts $(U, x^i)$, $(V, y^i)$, $\det [\frac{\partial y^i}{\partial x^i}]$ is everywhere positive on $U \cap V$
\end{definition}

\begin{theorem}
  \

  A manifold is orientable iff it has an oriented atlas.
\end{theorem}

\begin{customproof}
  \

  $(\implies)$ If $[(X_1, \dots, X_m)]$ is an orientation on $M$, each point $p \in M$ has a coordinate neighborhood $(U, x^i)$ on which $$(dx^1 \wedge \dots \wedge dx^m)(X_1, \dots, X_m) > 0$$For any two overlapping coordinate neighborhoods $(U, x^i)$, $(V, y^i)$ with the above property, $$(dx^1 \wedge \dots \wedge dx^m)(X_1, \dots, X_m) > 0$$ and $$(dy^1 \wedge \dots \wedge dy^m)(X_1, \dots, X_m) > 0$$ so $$dy^1 \wedge \dots \wedge dy^m = \det[\frac{\partial y^i}{\partial x^j}] dx^1 \wedge \dots \wedge dx^m$$and $\det[\frac{\partial y^i}{\partial x^j}] > 0$

  $(\impliedby)$ Suppose $\{U, x^i\}$ is an oriented atlas. For each $p \in U$, let $\mu_p = [(\frac{\partial}{\partial x^1}, \dots, \frac{\partial}{\partial x^m})]$.
\end{customproof}

\section{Manifolds With Boundary}

\subsection{Smooth Invariance of the Domain in $\mathbb{R}^n$}
\

\begin{definition}[Upper Half Space]
  \

  The \textbf{closed upper half space} is given by $$\mathcal{H}^n = \{(x^1, \dots, x^n) \in \mathbb{R} \mid x^n \geq 0\}$$Points where $x^n > 0$ are called \textbf{interior points}, the set of these is denoted $(\mathcal{H}^n)^\circ$, and points where $x^n = 0$ are called the \textbf{boundary points}, denoted $\partial(x^n)$.
\end{definition}

\begin{definition}[Smooth at a Point]
  \

  Let $S \subset \mathbb{R}^n$ be an arbitrary subset. A function $f: S \to \mathbb{R}^m$ is \textbf{smooth at $p$} in $S$ if there exists a neighborhood $U$ of $p$, a smooth function $\tilde{f}: U \to \mathbb{R}^m$, and $\tilde{f} = f$ on $U \cap S$. The function is \textbf{smooth on $S$} if it is smooth at each point of $S$.
\end{definition}

\begin{theorem}
  \

  Let $U \subset \mathbb{R}^n$ be an open subset, $S \subset \mathbb{R}^n$ an arbitrary subset, and $f: U \to S$ a diffeomorphism. Then, $S$ is open in $\mathbb{R}^n$
\end{theorem}

\begin{customproof}
  \

  Let $f(p)$ be some point in $S$, there exists some $g: S \to \mathbb{R}^n$ such that $g\circ f = 1_U$. By chain rule, $$g_{*, f(p)} \circ f_{*,p} = 1_{T_pU}$$so $f_{*,p}$ is injective, and since $U, V$ have the same dimension, it is also invertible. By inverse function theorem, $f$ is locally invertible, so there are neighborhoods $U_p$ of $p$ and $V_{f(p)}$ of $f(p)$ such that $f: U_P \to V_{f(p)}$ is a diffeomorphism, and $f(p) \in V_{f(p)} = f(U_p) \subset f(U) = S$ and since $V$ is open in $\mathbb{R}^n$, and $V_{f(p)}$ is open in $V$, $V_{f(p)}$ is open in $\mathbb{R}^n$. By the local criterion for openness, $S$ open in $\mathbb{R}^n$.
\end{customproof}

\begin{proposition}
  \

  Let $U$, $V$ be open subsets of the upper half space $\mathcal{H}^n$ and $f: U \to V$ a diffeomorphism. Then, $f$ maps interior points to interior points,and exterior points to exterior points.
\end{proposition}

\begin{customproof}
  \

  If $p$ is an interior point, there exists a neighborhood $p \in U$ around $p$. But since $f$ a diffeomorphism, it maps open sets to open sets, so $f(U)$ is open, and so there exists an open set containing $f(p)$, hence $f(p)$ is an interior point.

  If $p$ is an exterior point, and $f(p)$ was not an exterior point, then take some $f(p) \in V$, and consider $f ^{-1}(V)$ also open, since $f$ a diffeomorphism. But then, there is an open set containing $p$, so $p$ is not exterior, a contradiction.
\end{customproof}

\subsection{Manifolds With Boundary}
\

\begin{definition}[Manifolds with Boundary]
  \

  A \textbf{topological $n$-manifold with boundary} is a second countable, Hausdorff, topological space that is locally $\mathcal{H}^n$
\end{definition}

\begin{definition}[Charts]
  \

  For $n \geq 2$, a \textbf{chart} on a topological manifold $M$ is a pair $(U, \varphi)$, where $U \subset M$ is an open set, and $\varphi: U \to \mathcal{H}^n$ is a homeomorphism where $\varphi(U)$ is open.

  For $n = 1$, a chart $(U,\varphi)$ is instead a homeomorphism onto $\mathcal{H}$ or $\mathcal{L} = \{x \in \mathbb{R} \mid x \leq 0\}$, where $\varphi(U)$ is open.

  A collection of charts $\{(U, \varphi)\}$ is an \textbf{atlas} if for any two charts $(U, \varphi)$ and $(V, \psi)$, the transition map $\psi \circ \varphi ^{-1}$ is a diffeomorphism.
\end{definition}

\begin{definition}[Smooth Manifolds with Boundary]
  \

  A \textbf{smooth manifold with boundary} is a topological manifold with boundary equipped with a maximal smooth atlas.

  A point $p \in M$ a smooth manifold with boundary is an \textbf{interior point} if $\varphi(p)$ is an interior point. If $\varphi(p)$ is a boundary point, then $p$ is a \textbf{boundary point}.

\end{definition}

\begin{proposition}
  \

  Show that if $p$ is an interior point with respect to some chart $(U, \varphi)$, it is an interior point with respect to some other chart $(V, \psi)$.
\end{proposition}

\begin{customproof}
  \

  For two charts $(U, \varphi)$, $(V,\psi)$, if $\varphi(p)$ is an interior point, then since $\psi \circ \varphi ^{-1}$ is a diffeomorphism, which carries interior points to interior points and boundary points to boundary points, $\psi(p) = \psi \circ \varphi ^{-1}(\varphi(p))$ is also an interior point. Same goes for boundary points.
\end{customproof}

\begin{example}
  \

  The manifold boundary, defined above, and the topological boundary of a set don't always coincide.

  Let $A$ be the open unit disk in $\mathbb{R}^2$. Then, the topological boundary is the unit circle, the set of all points such that every neighborhood contains points in $A$ and $\mathbb{R}^2 \setminus A$. The manifold boundary, however, is the empty set, since as a manifold, the open disk is diffeomorphic to $\mathbb{R}^2$, without boundary.
\end{example}

\subsection{The Boundary of a Manifold with Boundary}
\

\begin{example}
  \

  Let $M$ be a manifold of dimension $n$, with boundary $\partial M$. Let $(U,\varphi)$ be a chart on $M$, and $\varphi' = \varphi \vert_{U \cap \partial M}$. Since $\varphi$ maps boundary points to boundary points, $$\varphi': U \cap \partial M \to \partial \mathcal{H}^n = \mathbb{R}^{n-1}$$and since $\varphi'$ is just a restriction, they form a smooth atlas on $\partial M$. Hence, $\partial M$ is a smooth manifold of dimension $n-1$ without boundary.
\end{example}

\subsection{Tangent Vectors, Differential Forms, and Orientations}
\

\begin{remark}
  \

  Tangent vectors/spaces are still derivations on germs of functions, $k$-forms are still smooth sections of $\bigwedge^k(T^*M)$, and orientations are still continuous pointwise orientations.

  One interesting thing happens at the boundary. Take for example $\mathcal{H}^2$, and consider a point $p$ on the boundary. $T_p\mathcal{H}^2$ still has basis $\left.\frac{\partial}{\partial x}\right\vert_p, \left.\frac{\partial}{\partial y}\right\vert_p$, and even though $\left.-\frac{\partial}{\partial y}\right\vert_p$ points out of the manifold, it is still a valid derivation, and hence a valid tangent vector. However, you will not find an integral curve with initial velocity $\left.-\frac{\partial}{\partial y}\right\vert_p$
\end{remark}

\subsection{Outward-Pointing Vector Fields}
\

\begin{definition}[Outward and Inward Pointing Vector Fields]
  \

  Let $M$ be a manifold with boundar, $p \in \partial M$. A tangent vector $X_p \in T_pM$ is \textbf{inward pointing} if $X_p \notin T_p(\partial M)$ and there is a positive real number $\varepsilon > 0$ and a curve $c: [0,\varepsilon) \to M$ such that $c(0) = p$, $c(]0,\varepsilon[) \subset M^\circ$ and $c'(0) = X_p$

  A vector $X_p$ is \textbf{outward pointing} if $-X_p$ is inward pointing.

  A vector field is \textbf{along} $\partial M$ if it is a vector field which assigns each $p \in M$ a vector in $T_pM$ (instead of $T_p(\partial M)$).
\end{definition}
\begin{proposition}
  \

  On a manifold $M$ with boundary $\partial M$, there is a smooth outward pointing vector field along $\partial M$.
\end{proposition}

\begin{customproof}
  \

  Cover $\partial M$ with coordinate open sets $(U_\alpha, x^i_\alpha)$ in $M$. On each $U_\alpha$ the vector field $X_\alpha = -\frac{\partial}{\partial x^n_\alpha}$ is smooth and outward pointing. Choose a PoU $\{\rho_\alpha\}$ suboordinate to $\{U_\alpha \cap \partial M\}$, and $\sum \rho_\alpha X_\alpha$ is a smooth outward pointing vector field along $\partial M$.
\end{customproof}

\subsection{Boundary Orientations}
\

\begin{proposition}
  \

  Let $M^m$ be an oriented manifold with boundary. If $\omega$ is an orientation form on $M$ and $X$ is a smooth outward-pointing vector field on $\delta M$, then $\iota_X \omega$ is a smooth, nowhere vanishing $(n-1)$ form on $\delta M$, so $\delta M$ is orientable.
\end{proposition}

\begin{customproof}
  \

  Since $\omega, X$ smooth on $\delta M$, so is $\iota_X \omega$. Suppose $\iota_X \omega$ vanishes at some point $p \in \partial M$. Then, $(\iota_X \omega)_p(v_1, \dots, v_{n-1}) = 0$ for all $v_1, \dots, v_{n-1} \in T_p(\partial M)$, but then if $e_1, \dots, e_{n-1}$ is a basis for $T_p(\partial M)$, then $$\omega_p(X_p, e_1, \dots, e_{n-1}) = (\iota_X \omega)_p(e_1, \dots, e_{n-1}) = 0$$And $\omega_p = 0$ on $T_pM$ as well. But then, this is a contradiction since orientation forms must be nowhere vanishing. Then, $\iota_X \omega$ is nowhere vanishing and so $\partial M$ is orientable.
\end{customproof}

\begin{proposition}
  \

  Suppose $M$ is an oriented $n$-manifold with boundary. Let $p$ be a point of the boundary $\partial M$ and let $X_p$ be an outward-pointing vector in $T_pM$. An ordered basis $(v_1, \dots, v_{n-1})$ tfor $T_p(\partial M)$ represents the boundary orientation at $p$ iff the ordered basis $(X_p, v_1, \dots, v_{n-1})$ for $T_pM$ represents the orientation on $M$ at $p$.

  "Outward vector first"
\end{proposition}

\begin{customproof}
  \

  For $p \in \partial M$, let $(v_1,\dots, v_{n-1})$ be an ordered basis for the tangent space $T_p(\partial M)$. Then, $(v_1, \dots, v_{n-1})$ represents the boundary orientation on $\partial M$ at $p$ iff $(\iota_{X_p} \omega_p)(v_1, \dots, v_{n-1}) > 0$ iff $\omega_p(X_p, v_1, \dots, v_{n-1}) > 0$ iff $(X_p, v_1, \dots, v_{n-1})$ represents the orientation on $M$ at $p$.
\end{customproof}

\section{Integration on Manifolds}

\subsection{The Riemann Integral over a Function on $\mathbb{R}^n$}
\

\begin{definition}[Closed Rectangles]
  \

  A \textbf{closed rectangle} in $\mathbb{R}^n$ is the Cartesian product $R = [a^1, b^1] \times \dots \times [a^n, b^n]$ of closed intervals in $\mathbb{R}$. The \textbf{volume} of a closed rectangle $R$ is given by $$\text{vol}(R) = \Pi_{i=1}^n (b_i-a_i)$$
\end{definition}

\begin{definition}[Partitions]
  \

  A \textbf{partition} of the closed interval $[a,b]$ is a set $P = \{p_0, \dots, p_n\}$ where $ a= p_0 < p_1 < \dots < p_{n-1} < p_n = b$

  A \textbf{partition of the rectangle} $R$ is a collection of partitions $P = \{P_i\}$ where each $P_i$ is a partition of $[a^i, b^i]$ of the rectangle. A partition $P$ divides $R$ into subrectangles, denoted $R_j$

  The \textbf{upper sum} of a function $f$ with respect to a partition $P$ is given by $$U(f,P) = \sum (\sup_{R_j} f)\text{vol}(R_j)$$

  The \textbf{lower sum} of a function $f$ with respect to a partition $P$ is given by $$L(f,P) = \sum (\inf_{R_j} f)\text{vol}(R_j)$$
\end{definition}

\begin{definition}[Refinements]
  \

  A partition $P' = \{P'_1, \dots, P'_n\}$ is a \textbf{refinement} of $P = \{P_1, \dots, P_n\}$ if each $P'_i \supseteq P_i$. Refinements divide each $R_j$ into subrectangles $R'_{jk}$
\end{definition}

\begin{proposition}
  \

  If $P'$ a refinement of $P$, then $$L(f, P) \leq L(f, P')$$$$U(f, P') \leq U(f,P)$$for any function $f: \mathbb{R}^n \to \mathbb{R}$
\end{proposition}

\begin{customproof}
  \

  If $R'_{jk} \subseteq R_j$ then $\inf_{R_j} \leq \inf_{R_{jk}}$, and the oppsite is true for $\sup$. Since the total sum of volume remains the same, the result follows.
\end{customproof}

\begin{definition}[Riemann Integrals]
  \

  Let $R \subseteq \mathbb{R}^n$ be a rectangle, $f: \mathbb{R}^n \to \mathbb{R}$ a function, then the \textbf{lower Riemann integral} is given by $$\underline{\int_R} = \sup_P L(f,P)$$and the \textbf{upper Riemann integral} is given by $$\overline{\int_R} = \inf_P U(f,P)$$

  $f$ is \textbf{Riemann integrable} if $\underline{\int_R} = \overline{\int_R}$, and is denoted $$\int_R dx^1\dots dx^n = \underline{\int_R} = \overline{\int_R}$$where $x^1, \dots, x^n$ are the standard coordinates on $\mathbb{R}^n$.
\end{definition}

\begin{definition}[Extension by Zero]
  \

  If $f: A \subset \mathbb{R}^n \to \mathbb{R}$, then the \textbf{extension of $f$ by zero} is the function $$\tilde{f}(x) =
  \begin{cases}
    f(x) & x \in A\\
    0 & x \notin A
  \end{cases}$$And we define $$\int_A f dx^1\dots dx^n = \int_R \tilde{f} dx^1 \dots dx^n$$if such an integral exists.
\end{definition}

\begin{definition}[Volume]
  \

  For $A \subset \mathbb{R}^n$, the \textbf{volume} of $A$ is given by $$\int_A 1 dx^1 \dots dx^n $$if such an integral exists.
\end{definition}

\subsection{Integrability Conditions}
\

\begin{definition}[Measure]
  \

  A set $A \subset \mathbb{R}^n$ has \textbf{measure zero} if for every $\varepsilon > 0$, there is a countable cover $\{R_i\}_{i=1}^\infty$ of $A$ by closed rectangles $R_i$ such that $$\sum_{i=1}^\infty \text{vol}(R_i) < \varepsilon$$
\end{definition}

\begin{theorem}[Lebesgue's Theorem]
  \

  A bounded function $f: A \to \mathbb{R}^n$ on a bounded subset $A \subset \mathbb{R}^n$ is Riemann integrable iff the set of discontinuities of $\tilde{f}$ forms a measure zero set.
\end{theorem}

\begin{proposition}
  \

  If a continuous function $f: U \to \mathbb{R}$ defined on an open subset $U$ of $\mathbb{R}^n$ has compact support, it is Riemann integrable.
\end{proposition}

\begin{customproof}
  \

  Since $f$ is continuous on a compact set, $f$ is bounded, and $\text{supp}(f)$ is closed and bounded in $\mathbb{R}^n$.

  Since $\tilde{f}$ agrees with $f$ on $U$, it is continuous for all $p \in U$. For any $p \notin U$, since $\mathbb{R}^n \setminus U$ is open, there is an open ball $p \in B$ and $B \cap U = \varnothing$, so $\tilde{f}\vert_B = 0$ and $\tilde{f}$ is continuous for all $p \notin U$, so $\tilde{f}$ is continuous everywhere, so it is continuous, and has no discontinuities, so by Lebesgue's theorem, it is Riemann integrable on $U$.
\end{customproof}

\begin{definition}[Domain of Integration]
  \

  A subset $A \subset \mathbb{R}^n$ is a \textbf{domain of integration} if it is bounded, and the topological boundary, $\text{bd}(A)$, has measure zero.
\end{definition}

\begin{proposition}
  \

  Every bounded continuous function $f$ defined on a domain of integration $A$ in $\mathbb{R}^n$ is Riemann integrable over $A$.
\end{proposition}

\subsection{The Integral of an $n$-Form over $\mathbb{R}^n$}
\

\begin{definition}[Integration of an $n$-Form]
  \

  Let $\omega = f(x) dx^1 \wedge \dots \wedge dx^n$ be a smooth $n$-form on an open subset $U \subset \mathbb{R}^n$ with standard coordinates $x^1, \dots, x^n$. The \textbf{integral} of $\omega$ over a subset $A \subset U$ is given by $$\int_A = \int_A f(x)dx^1\wedge\dots\wedge dx^n = \int_A f(x)dx^1\dots dx^n$$if such an integral exists.
\end{definition}

\begin{example}
  \

  Let $\omega = fdx^1 \wedge \dots \wedge dx^n$ be an $n$-form on $(U, x^i)$, $U\subseteq \mathbb{R}^n$, let $T: V \to U$ be a diffeomorphism, $(V, y^i)$ be another chart. Then, $dT^1 \wedge \dots \wedge dT^n = \det(J(T))dy^1 \wedge \dots \wedge dy^n$ where $J(T)$ is the Jacobian of $T$. Then, $$\int_V T^*\omega = \int_V (T^*f)T^*dx^1 \wedge \dots \wedge T^*dx^n = \int_V (f \circ T) dT^1 \wedge \dots \wedge dT^n = \int_V (f \circ T)\det(J(T))dy^1 \wedge \dots \wedge dy^n$$$$= \int_V(f \circ T)\det(J(T))dy^1\dots dy^n$$

  However, $$\int_U \omega = \int f dx^1\dots dx^n = \int_V (f \circ T) \vert \det(J(T)) dy^1 \dots dy^n\vert$$by a formula from calculus?

  In essence, $$\int_V T^* \omega = \pm \int_U \omega$$and is invariant under orientation-preserving diffeomorphisms.
\end{example}

\subsection{Integral of a Differential Form over a Manifold}
\

\begin{definition}[Integral over Manifolds]
  \

  Let $M$ be an orientable manifold, and let there be an oriented atlas $\{(U_\alpha, \varphi_\alpha)\}$. Denote $\Omega_c^k(M)$ the vector space of $k$-forms with compact support on $M$. Let $(U, \varphi)$ be a chart in the oriented atlas, $\omega \in \Omega^m_c(U)$. Then, since $\varphi$ is a diffeomorphism onto its image, $(\varphi ^{-1})^* \omega$ is an $m$-form with compact support on $\varphi(U) \subset \mathbb{R}^m$. Define the \textbf{integral of $\omega$} to be $$\int_U \omega = \int_U (\varphi ^{-1})^* \omega$$

  This is well defined, since for another chart $(U,\psi)$, we have $$\int_{\varphi(U)} (\varphi ^{-1})^* \omega = \int_{\psi(U)} (\psi ^{-1})^*\omega$$since $\psi \circ \varphi ^{-1}$ is an orientation preserving diffeomorphism (by definition of an oriented atlas).

  Let $\{\rho_\alpha\}$ be subordinate to the open cover $\{U_\alpha\}$ of the oriented atlas. Since $\omega$ has compact support, each $\rho_\alpha$ has locally finite support, all except finitely many $\rho_\alpha \omega$ are identically zero, so $\omega = \sum_\alpha \rho_\alpha \omega$ is a sum over finitely many terms. Since $\text{supp}(\rho_\alpha\omega) \subset \text{supp}\rho_\alpha\cap \text{supp} \omega$ is a closed subset of $\text{supp} \omega$ which is compact, $\text{supp}(\rho_\alpha \omega)$ is compact as well. Then, the \textbf{integral over the manifold} given by $$\int_M \omega = \sum_\alpha \int_{U_\alpha} \rho_\alpha \omega$$is a well defined sum.
\end{definition}

\begin{proposition}
  \

  Let $M$ be an orientable manifold, $(U, \varphi)$ a chart on $M$ in the oriented atlas, $\omega, \tau \in \Omega^m_c(U)$ $$\int_U \omega + \tau = \int_U \omega + \int_U \tau$$
\end{proposition}

\begin{customproof}
  \

  Consequence of the linearity of the integral in $\mathbb{R}^m$.
\end{customproof}

\begin{proposition}
  \

  Let $\omega$ be an $m$-form with compact support on an oriented manifold $M^m$. If $-M$ denotes the same manifold, but opposite orientation, then $$\int_{-M} = -\int_M \omega$$
\end{proposition}

\begin{definition}[Parameterized Sets]
  \

  Let $M^m$ be an oriented manifold. A \textbf{parameterized set} in $M$ is a subset $A \subset M$, and a subset $D \subset \mathbb{R}^n$, where $D$ a compact domain of integration (has measure zero boundary), along with a smooth map $F: D \to M$, such that $F(D) = A$, and when restricted to $\text{int}(D)$, $F \vert_{\text{int}(D)}$ becomes an orientation preserving diffeomorphism.

  The smooth map $F: D \to A$ is called a \textbf{parameterization} of $A$.
\end{definition}

\begin{example}[Integration Over a $0$-Dimensional Manifold]
  \

  A $0$-dimensional oriented manifold is a collection of disconnected points, each with orientatino $\pm 1$. If $M = \sum p_i - \sum q_j$, with $p_i$ being points with positive orientation, $q_j$ being points with negative orientation, then $$\int_M f = \sum f(p_i) - \sum f(q_j)$$
\end{example}

\subsection{Stokes's Theorem}
\

\begin{theorem}
  \

  For any smooth $(m-1)$ form $\omega$ with compact support on $M^m$ oriented, $$\int_M d\omega = \int_{\partial M} \omega$$
\end{theorem}

\begin{customproof}
  \

  Choose an atlas $\{(U_\alpha, \varphi_\alpha)\}$ for $M$ where each $U_\alpha$ is diffeomorphic to $\mathbb{R}^n$ or $\mathcal{H}^n$. If Stokes's Theorem is true for these two spaces, it is true for any manifold by the following:

  Suppose Stokes's Theorem is true for $\mathbb{R}^n$ and $\mathcal{H}^n$. Then, choose a partition of unity $\{\rho_\alpha\}$ subordinate to $\{U_\alpha\}$, and
  \begin{align*}
    \int_{\partial M} \omega & = \int_{\partial M} \sum_\alpha (\rho_\alpha\omega) && \text{since $\sum_\alpha \rho_\alpha = 1$}\\
    &= \sum_\alpha \int_{\partial M} \rho_\alpha \omega && \text{by linearity, we can interchange the two since $\sum_{\alpha }\rho_\alpha \omega$ is a finite sum}\\
    &= \sum_{\alpha} \int_{\partial U_\alpha} \rho_\alpha \omega && \text{we can restrict since $\text{supp}(\rho_\alpha \omega) \subset U_\alpha$}\\
    &= \sum_\alpha \int_{U_\alpha} d(\rho_\alpha \omega) && \text{Stokes's theorem for $U_\alpha$}\\
    &= \sum_{\alpha} \int_M d(\rho_\alpha \omega) && \text{supp}(d(\rho_\alpha \omega)) \subset U_\alpha\\
    &= \int_M d(\sum \rho_\alpha \omega) && \text{linearity}\\
    &= \int_M d\omega
  \end{align*}

  Now it remains to show Stokes's Theorem on $\mathcal{H}^n$ and $\mathbb{R}^n$. We just show for $\mathcal{H}^2$, the others are analagous.

  Let $\mathcal{H}^2$ have orientation represented by the form $dx \wedge dy$. Then, on the boundary, the orientation is given by $\iota_{-\frac{\partial}{\partial y}}(dx \wedge dy) = dx$.

  Any $1$-form has the form $\omega = f(x,y) dx + g(x,y) dy$. Since $\omega$ is compactly supported, so are $f,g$, so we can choose a real number $a>0$ large enough such that $\text{supp}(f), \text{supp}(g) \subset [-a,a] \times [0,a]$. Then, $$d\omega = (g_x-f_y)dx \wedge dy$$where $g_x = \frac{\partial g}{\partial x}$.

  Then, $$\int_{\mathcal{H}^2} d\omega = \sum_{\mathcal{H}^2} g_x dxdy - \int_{\mathcal{H}^2} f_y dxdy = \int_0^a \int_{-a}^a g_xdxdy - \int_{-a}^a \int_0^a f_y dydx$$But note that $\int_{-a}^a g_x dx = g(x,y) \vert^a_{-a} = 0 - 0 = 0$ since $\text{supp}(g) \subset [-a,a] \times [0,a]$, and $\int_0^a f_y dy = f(x,y) \vert_0^a = -f(x,0)$. So $$\int_{\mathcal{H}^2} d\omega = \int_{-a}^a f(x,0) dx$$

  However, since $dy = 0$ on $\partial \mathcal{H}^2$, $$\int_{\partial \mathcal{H}^2} \omega = \int{\partial \mathcal{H}^2} f dx + g dy = \int_{\partial \mathcal{H}^2} f dx = \int_{-a}^a f(x,0) dx$$that last equality is because $\partial \mathcal{H}^2$ is only when $y = 0$.
\end{customproof}

\subsection{Line Integrals and Green's Theorem}
\

\begin{remark}
  \

  Note that for $F: \mathbb{R}^3 \to \mathbb{R}^3$, $F = \langle P, Q, R \rangle$, and $r = (x,y,z)$, we have $F \cdot dr = P dx + Q dy + R dz$.
\end{remark}

\begin{theorem}[Fundamental Theorem for Line Integrals]
  \

  Let $C$ be a curve in $\mathbb{R}^3$, parameterized by $r(t) = (x(t), y(t), z(t))$, $a \leq t \leq b$, and $F$ a vector field in $\mathbb{R}^3$. If $F = \text{grad}(f)$ for some scalar function $f$, then $$\int_C F \cdot dr = f(r(b)) - f(r(a))$$
\end{theorem}

\begin{customproof}
  \

  Take $M$ to be a curve $C$ with parameterization $r(t)$, $a \leq t \leq b$, and $\omega$ to be $f$ on $C$. By Stokes's, $$\int_C d\omega = \int_C df = \int_C \frac{\partial f}{\partial x} dx + \frac{\partial f}{\partial y} dy + \frac{\partial f}{\partial z} dz = \int_C \text{grad}(f) \cdot dr$$and $$\int_{\partial C} = f \vert_{r(a)}^{r(b)} = f(r(b)) - f(r(a))$$and so $$\int_C \text{grad}(f) \cdot dr = \int_C F \cdot dr = f(r(b)) - f(r(a))$$
\end{customproof}

\begin{theorem}
  \

  If $D$ is a plane region with boundary $\partial D$, and $P, Q$ are smooth functions on $D$, then $$\int_{\partial D} P dx + Q dy = \int_D (\frac{\partial Q}{\partial x} - \frac{\partial P}{\partial y} )dA$$(where $dA = dxdy$)
\end{theorem}

\begin{customproof}
  \

  Let $M$ be $D$ of dimension $2$, let $\omega = Pdx + Qdy$ on $D$ a $1$-form. Then, $$\int_{\partial D} \omega = \int_{\partial D} Pdx + Qdy$$and$$\int_D d\omega = \int_D P_y dy\wedge dx + Q_xdx\wedge dy = \int_D (Q_x-P_y) dx\wedge dy = \int_D (Q_x - P_y) dxdy$$And we have our proof by Stokes's.
\end{customproof}

\section{De Rham Cohomology}

\subsection{De Rham Cohomology}
\

\begin{remark}
  \

  Let $F(x,y) = \langle P(x,y), Q(x,y) \rangle$ be a vector field in $\mathbb{R}^2$, and suppose we want to find its line integral along some curve $C$. If $F$ is the gradient of some scalar function, $F = \text{grad}(f) = \langle f_x, f_y\rangle$, we can use Stokes's theorem to say $$\int_C Pdx + Qdy = \int_C f_xdx + f_ydy = f(q) - f(p)$$if $C$ has endpoints $p, q$.

  A necessary condition for $F = \langle P, Q \rangle $ is that $P_y = Q_x$, or $$f_{xy} = f_{yx}$$But if the above is true, how can we know that $F = \langle P, Q \rangle$ is the gradient of some scalar function on $U$?

  It turns out on subsets $U \subseteq \mathbb{R}^n$, there exists a correspondence
  \[
    \begin{tabular}{c | c}
      vector fields & differential $1$-forms\\
      $F = \langle P, Q \rangle$ & $\omega = P dx + Q dy$\\
      $\text{grad}(f) = \langle f_x, f_y \rangle$ & $df = f_x dx + f_ydy$\\
      $Q_x - P_y = 0$ & $d\omega = (Q_x-P_y)dx\wedge dy = 0$
  \end{tabular}\]And now we want to measure if closed forms are exact
\end{remark}

\begin{definition}[Closed and Exact Forms]
  \

  A differential form $\omega \in \Omega^k(M)$ for a manifold $M$ is \textbf{closed} if $d\omega = 0$

  A differential $k$-form $\omega$ is \textbf{exact} if there is some $k-1$ form $\tau$ such that $d\tau = \omega$.

  Let $Z^k(M)$ denote the vector space of all the closed $k$-forms on $M$, $B^k(M)$ denote all of the exact $k$-forms on $M$. Since $d^2 = 0$, every exact form is closed, and $B^k(M) \subseteq Z^k(M)$.
\end{definition}

\begin{definition}[De Rham Cohomology]
  \

  The \textbf{de Rham cohomology} is given by $$H^k(M) = Z^k(M)/B^k(M)$$And this quotient induces an equivalence relation on $Z^k(M)$, $\omega, \omega' \in Z^k(M)$ are related (and hence represent the same equivalence class in $H^k(M)$) iff $\omega - \omega' \in B^k(M)$, or only if their difference is exact.

  The equivalence class of $\omega$ is called its \textbf{cohomology class}. Two forms are \textbf{cohomologous} if $\omega' = \omega + d\tau$ for some $k-1$-form $\tau$.
\end{definition}

\begin{theorem}
  \

  If $M$ has $r$ connected components, then $H^0(M) \cong \mathbb{R}^r$. An element of $H^0(M)$ is specified by an ordered $r$-tuple of real numbers, each real number representing a constant function on a connected component of $M$.
\end{theorem}

\begin{customproof}
  \

  No $0$-forms $\omega$ are exact, since there is no form $\tau$ where $d\tau = \omega$, this isn't even defined. Hence, $B^0(M) = \{0\}$ and so $$H^k(M) = Z^k(M)/B^k(M) \cong Z^k(M)$$$0$-forms are just functions $f: M \to \mathbb{R}$, but they are closed iff $df = \sum \frac{\partial f}{\partial x^i} x^i = 0$ iff $f$ is constant on every connected component of the manifold. Then, we have our result.
\end{customproof}

\begin{proposition}
  \

  On a manifold $M^m$, the deRham cohomology $H^k(M)$ vanishes for $k > m$.
\end{proposition}

\begin{customproof}
  \

  There are no $k$-forms for $k>m$ besides the $0$ form, so we just have $H^k(M) = \{0\}/\{0\} = \{0\}$.
\end{customproof}

\subsection{Examples of the deRham Cohomology}
\

\begin{example}[$\mathbb{R}$]
  \

  Let's find the cohomologies of $\mathbb{R}$.

  Since $\mathbb{R}$ has one connected component, $H^0(\mathbb{R}) = \mathbb{R}$. There are no nonzero $2$-forms on $\mathbb{R}^1$, so every $1$-form is closed. Then, for some $1$-form $\omega = f(x) dx$, suppose $\tau = g(x)$ a $0$-form, then we want $\omega = d\tau = g'(x) dx$, but then $$g(x) = \int_0^x f(x) dx$$so every $1$-form is exact as well, so $H^1(\mathbb{R}) = \{0\}$

  Hence, the cohomologies are $$H^k(\mathbb{R}) =
  \begin{cases}
    \mathbb{R} & k = 0\\
    0 & k \geq 1
  \end{cases}$$
\end{example}

\begin{lemma}
  \

  Suppose $\bar{f}$ is a smooth periodic function of period $2\pi$ on $\mathbb{R}$, and $\int_0^{2\pi} \bar{f}(u) du = 0$. Then, $\bar{f} dt = d\bar{g}$ for a smooth periodic function $\bar{g}$ of period $2\pi$ on $\mathbb{R}$.
\end{lemma}

\begin{customproof}
  \

  Let $g \in \Omega^0(\mathbb{R})$ by $$\bar(g) = \int_0^t \bar{f}(u) du$$Then, $$g(t + 2\pi) = \int_0^{2\pi} \bar{f}(u) du + \int_{2\pi}^{t + 2\pi} \bar{f}(u) du = 0 + \int_{2\pi}^{t + 2\pi} \bar{f}(u) du = \int_0^t \bar{f}(u) du = g(t)$$So $\bar{g}(t)$ is also $2\pi$ periodic, and $d\bar{g} = \bar{f}(t) dt$.
\end{customproof}

\begin{example}[$S^1$]
  \

  Since $S^1$ is connected, $H^0(S^1) = \mathbb{R}$, and $H^k(S^1) = 0$ for all $k \geq 2$.

  Since a circle has dimension $1$, all $1$-forms are closed, so $\Omega^1(S^1) = Z^1(S^1)$.

  Note before, we found a nowhere vanishing $1$-form on $S^1$, being $\omega = -ydx + x dy$. If $F = h \circ \iota: [0, 2\pi] \to S^1$, where $h(t) = (\cos t, \sin t)$, and $\iota: [0,2\pi] \to \mathbb{R}$ the inclusion map, then $F^* \omega = dt$, and so $$\int_{S^1} \omega = \int_{[0,2\pi]} F^*\omega = \int_0^{2\pi} dt = 2\pi$$

  Consider $\varphi: Z^1(S^1) = \Omega^1(S^1) \to \mathbb{R}$, mapping $\varphi(\alpha) = \int_{S^1} \alpha$. Since this map is linear, and there is a nonzero value $\varphi(\omega) = 2\pi$, it is surjective.

  By Stokes's theorem, the exact $1$-forms are in $\ker \varphi$, and we want to show every $1$-form in $\ker\varphi $ is exact.

  Let $\alpha = f\omega$ a smooth $1$-form such that $\varphi(\alpha) = 0$. Then, $\bar{f} = f \circ h$ must be $2\pi$ periodic, with $\int_0^{2\pi} \bar{f} = 0$.

  But then, by the prior lemma, there exists a $\bar{g}$ that is $2\pi$ periodic on $\mathbb{R}$, and so this induces a function $g$ on $S^1$, and $dg = \alpha$. Hence, $\ker\varphi$ has only exact forms, and so $H^1(S^1) = Z^1(S^1)/B^1(S^1) \cong \mathbb{R}$ by first isomorphism theorem of vector spaces.
\end{example}

\subsection{Diffeomorphism Invariance}
\

\begin{lemma}
  \

  If $F: M \to N$ is a smooth map, $F^*: \Omega^*(N) \to \Omega^*(M)$ sends closed forms to closed forms, and exact forms to exact forms.
\end{lemma}

\begin{customproof}
  \

  Since the pullback commutes with the derivative, if $\omega$ is closed in $\Omega^*(N)$, $$dF^* \omega = F^* d\omega = F^* 0 = 0$$

  If instead, $\omega = d\tau$ is exact, then $F^* \omega = F^*d\tau = dF^*\tau$ so $F^*\omega$ exact as well.
\end{customproof}

\begin{corollary}
  \

  $F^*$ induces a linear map of quotient spaces $$F^\#: \frac{Z^k(N)}{B^k(N)} \to \frac{Z^k(M)}{B^k(M)}$$that sends $$F^\#: [\omega] \mapsto [F^* \omega]$$This map is called the \textbf{pullback map in cohomologies}
\end{corollary}

\subsection{Ring Structure on Cohomologies}
\

\begin{definition}[Graded Cohomology Algebra]
  \

  Given a manifold $M^m$, the \textbf{cohomology algebra} $$H^*(M) = \bigoplus_{k=0} H^k(M)$$forms an anticommutative graded $\mathbb{R}$-algebra. Given $[\omega] \in H^k(M)$, $[\tau] \in H^l(M)$, we define multiplication by $$[\omega] \wedge [\tau] = [\omega \wedge \tau] \in H^{k+l}(M)$$This product is well defined, since $[\omega \wedge \tau]$ is indeed closed, and it is independent of choice of representative.
\end{definition}

\subsection{Exercises}
\

\begin{exercise}
  \

  Prove that a nowhere vanishing $1$-form on a compact manifold $M$ cannot be exact.
\end{exercise}

\begin{solution}
  \

  Suppose that $\omega$ is a nowhere vanishing $1$-form, and $\omega = d\tau$ is exact, for some $0$-form $\tau$. Since $0$-forms are just real valued functions, for a chart $(U, x^i)$, $$\omega = d\tau = \sum_{i} \frac{\partial f}{\partial x^i} dx^i$$But the extreme value theorem states that any continuous function from a compact set to the real line must attain both a minimum and a maximum, hence at those points, $\omega = f'(x) dx = 0$ hence $\omega$ cannot be nowhere vanishing.
\end{solution}

\begin{exercise}
  \

  Suppose that $M$ has an infinite number of connected components. What is $H^0(M)$?
\end{exercise}

\begin{solution}
  \

  If $M^m$ has an infinite number of connected components, then since the basis of $M$ must be second countable, the number of components must be countable. Then, by a prior theorem, $H^0(M) = \mathbb{R}^k$ where $k$ is the number of connected components in $M$, so $$H^0(M) = \Pi_{i=0}^\infty \mathbb{R}$$
\end{solution}

\section{The Long Exact Sequence in Cohomology}

\subsection{Exact Sequences}
\

\begin{definition}[Cochain Complex]
  \

  A \textbf{cochain complex} is a collection of vector spaces $\{C^k\}_{k \in \mathbb{Z}}$ with a sequence of linear maps $d_k: C^k \to C^{k+1}$ \[
    \begin{tikzcd}
      \dots & C^{-1} & C^0 & C^1 & \dots
      \arrow[from=1-1, to=1-2, "d_{-2}"]
      \arrow[from=1-2, to=1-3, "d_{-1}"]
      \arrow[from=1-3, to=1-4, "d_0"]
      \arrow[from=1-4, to=1-5, "d_1"]
  \end{tikzcd} \]such that $d_k \circ d_{k-1} = 0$ for all $k$. The morphisms $d_k$ are called \textbf{differentials}.
\end{definition}

\begin{example}
  \

  The vector space $\Omega^*(M)$ with the exterior derivative $d$ forms a cochain complex called the \textbf{de Rham complex}:

  \[
    \begin{tikzcd}
      0 & \Omega^0(M) & \Omega^1(M) & \dots
      \arrow[from=1-1, to=1-2]
      \arrow[from=1-2, to=1-3, "d"]
      \arrow[from=1-3, to=1-4, "d"]
  \end{tikzcd}\]
\end{example}

\begin{definition}[Exact Sequences]
  \

  A sequence of homomorphisms of vector spaces \[
    \begin{tikzcd}
      A & B & C
      \arrow[from=1-1, to=1-2, "f"]
      \arrow[from=1-2, to=1-3, "g"]
  \end{tikzcd}\]is \textbf{exact} if $\text{im}(f) = \text{ker}(g)$.

  A sequence of vector spaces and homomorphisms \[
    \begin{tikzcd}
      A^0 & A^1 & A^2 & \dots & A^n
      \arrow[from=1-1, to=1-2, "f_0"]
      \arrow[from=1-2, to=1-3, "f_1"]
      \arrow[from=1-3, to=1-4, "f_2"]
      \arrow[from=1-4, to=1-5, "f_{n-1}"]
  \end{tikzcd}\] is an \textbf{exact sequence} if $\text{im}(f_{i-1}) = \ker(f_i)$ for all $i$.

  A five term exact sequence of the form \[
    \begin{tikzcd}
      0 & A & B & C & 0
      \arrow[from=1-1, to=1-2]
      \arrow[from=1-2, to=1-3]
      \arrow[from=1-3, to=1-4]
      \arrow[from=1-4, to=1-5]
  \end{tikzcd}\]is called a \textbf{short exact sequence}
\end{definition}

\begin{proposition}
  \

  Suppose \[
    \begin{tikzcd}
      A & B & C
      \arrow[from=1-1, to=1-2, "f"]
      \arrow[from=1-2, to=1-3, "g"]
  \end{tikzcd}\]is exact. Then,

  \begin{enumerate}
    \item $f$ is surjective iff $g$ is the zero map
    \item $g$ is injective iff $f$ is the zero map
  \end{enumerate}
\end{proposition}

\begin{customproof}
  \

  \begin{enumerate}
    \item $(\implies)$ if $f$ is surjective, $\text{im} f = \ker g = B$, so $g$ maps every vector in $B$ into $0_C$, so $g$ is the zero map.

      $(\impliedby)$ if $g$ is the zero map, $\ker g = B = \text{im} f$, so $f$ is surjective

    \item $(\implies)$ If $g$ is injective, then for vector space morphisms, this only happens iff $\ker g = 0 = \text{im} f$, so $f$ maps all vectors in $A$ to $0_B$, so it is the zero map.

      $(\impliedby)$ If $f$ is the zero map, $\text{im} f = 0 = \ker g$, but this implies that $g$ is injective.
  \end{enumerate}
\end{customproof}

\begin{proposition}
  \

  \begin{enumerate}
    \item The four term short exact sequence \[
        \begin{tikzcd}
          0 & A & B & 0
          \arrow[from=1-1, to=1-2, "f_1"]
          \arrow[from=1-2, to=1-3, "f"]
          \arrow[from=1-3, to=1-4, "f_2"]
      \end{tikzcd}\]is exact iff $f: A \to B$ is an isomorphism.
    \item If \[
        \begin{tikzcd}
          A & B & C & 0
          \arrow[from=1-1, to=1-2, "f"]
          \arrow[from=1-2, to=1-3, "f_1"]
          \arrow[from=1-3, to=1-4, "f_2"]
      \end{tikzcd}\]is an exact sequence of vector spaces, then there is a vector space isomorphism $$C \cong \text{coker}(f) = \frac{B}{\text{im} f}$$
  \end{enumerate}
\end{proposition}

\begin{customproof}
  \

  \begin{enumerate}
    \item $(\implies)$ Suppose the sequence was exact. Then, $\text{im}(f_1) = \text{ker}(f)$ but since $f_1$ is a vector space morphism, and $0$ has only one vector, the zero vector, $f_1$ must send the zero vector to the zero vector, so $\text{im}(f_1) = \ker(f) = 0$ so $f$ injective. Also, $f_2$ maps $C$ to the zero vector, so $\ker(f_2) = \text{im}(f) = B$, so $f$ surjective. Since $f$ is a surjective and injective vector space morphism, it is an isomorphism.

      $(\impliedby)$ Suppose that $f$ is an isomorphism. Then, $\ker(f) = 0$ and $\text{im}(f) = B$. But $\ker(f_2) = B$, since it is the zero map, and $\text{im}(f_1) = 0$, hence we fulfill the criteria for this sequence to be exact.
    \item We have $\text{im}(f) = \ker(f_1)$. But for some surjective vector space morphism $g: B \to C$, $C \cong \frac{B}{\ker g}$ by first isomorphism theorem.

      But since the sequence is exact, $\text{ker}(f_2) = \text{im}(f_1) = C$, so $f_1$ is a surjective vector space morphism, so $C \cong \frac{B}{\ker g}$
  \end{enumerate}
\end{customproof}

\subsection{Cohomology of Cochain Complexes}
\

\begin{definition}[Cohomology Vector Space]
  \

  If $\mathcal{C}$ is a cochain complex, then since $d_k \circ d_{k-1} = 0$, we have $\text{im} (d_{k-1}) \subset \ker (d_k)$, and we can form the quotient vector space $$H^k(\mathcal{C}) = \frac{\ker d_k}{\text{im} d_{k-1}}$$called the \textbf{$k$th cohomology vector space}, and which basically "measures" the extent to which $\mathcal{C}$ fails to be exact at $C^k$

  Elements of the $k$th vector space $C^k$ are called \textbf{$k$-cochains}. A cochain in $\ker(d_k)$ is called a \textbf{$k$-cocycle} and a cochain in $\text{im}(d_{k-1})$ is called a \textbf{$k$-coboundary}.

  For $c \in \text{ker}(d_k)$, the equivalence class $[c] \in H^k(\mathcal{C}) = \frac{\ker(d_k)}{\text{im}(d_{k-1})}$ is called its \textbf{cohomology class}.

  We denote the subspaces of $k$-cocyles as $Z^k(\mathcal{C}) = \ker(d_k)$ and $k$-coboundaries as $B^k(\mathcal{C}) = \text{im}(d_{k-1})$
\end{definition}

\begin{example}
  \

  In the de Rham complex, a cocycle is a closed form, and a coboundary is an exact form.
\end{example}

\begin{definition}[Cochain Map]
  \

  Let $\mathcal{A}$ and $\mathcal{B}$ be two cochain complexes with differentials $d_i$, $d'_j$ respectively. A \textbf{cochain map} is a map $\varphi: \mathcal{A} \to \mathcal{B}$ such that at every vector space $A^k \in \mathcal{A}$, $\varphi_k: A^k \to B^k$ is a vector space morphism, and the differentials commute: $$d'_k \circ \varphi_k = \varphi_{k+1} \circ d_k$$or in other words, the following diagram commutes: \[
    \begin{tikzcd}
      \dots & A^{k-1} & A^k & A^{k+1} & \dots\\
      \dots & B^{k-1} & B^k & B^{k+1} & \dots
      \arrow[from=1-1, to=1-2, "d_{k-2}"]
      \arrow[from=1-2, to=1-3, "d_{k-1}"]
      \arrow[from=1-3, to=1-4, "d_k"]
      \arrow[from=1-4, to=1-5, "d_{k+1}"]
      \arrow[from=2-1, to=2-2, "d'_{k-2}"']
      \arrow[from=2-2, to=2-3, "d'_{k-1}"']
      \arrow[from=2-3, to=2-4, "d'_k"']
      \arrow[from=2-4, to=2-5, "d'_{k+1}"']
      \arrow[from=1-2, to=2-2, "\varphi_{k-1}"]
      \arrow[from=1-3, to=2-3, "\varphi_k"]
      \arrow[from=1-4, to=2-4, "\varphi_{k+1}"]
  \end{tikzcd}\]A cochain map $\varphi: \mathcal{A} \to \mathcal{B}$ induces a vector space morphism between cohomologies: $$\varphi^*_k: H^k(\mathcal{A}) \to H^k(\mathcal{B})$$$$\varphi^*_k: [a] \mapsto [\varphi_k(a)]$$
\end{definition}

\begin{proposition}
  \

  Show that $\varphi^*_k$ is well defined, that is, show $\varphi_k$ takes cocycles to cocycles and coboundaries to coboundaries
\end{proposition}

\begin{customproof}
  \

  For a cocyle $a \in Z^k(\mathcal{A})$, we have $d'_k\varphi_k(a) = \varphi_{k+1}(da) = \varphi_{k+1}(0) = 0$ so $\varphi_k(a) \in Z^k(\mathcal{B})$

  For a coboundary $b = d_{k-1}x \in B^k(\mathcal{A})$ for some $x \in A^{k-1}$, we have $\varphi_k(b) = \varphi_k(d_{k-1}x) = d_k\varphi_{k-1}(x)$ so $\varphi_k(b) \in B^k(\mathcal{B})$
\end{customproof}

\subsection{The Connecting Homomorphism}
\

\begin{definition}[Short Exact Sequence of Cochain Complexes]
  \

  A sequence of cochain complexes \[
    \begin{tikzcd}
      0 & \mathcal{A} & \mathcal{B} & \mathcal{C} & 0
      \arrow[from=1-1, to=1-2]
      \arrow[from=1-2, to=1-3, "i"]
      \arrow[from=1-3, to=1-4, "j"]
      \arrow[from=1-4, to=1-5]
  \end{tikzcd}\] is \textbf{short exact} if $i$ and $j$ are cochain maps, and for a fixed $k$, \[
    \begin{tikzcd}
      0 & A^k & B^k & C^k & 0
      \arrow[from=1-1, to=1-2]
      \arrow[from=1-2, to=1-3, "i_k"]
      \arrow[from=1-3, to=1-4, "i_j"]
      \arrow[from=1-4, to=1-5]
  \end{tikzcd}\]is a short exact sequence of vector spaces.

\end{definition}

\begin{proposition}
  \

  Given a short exact sequence of cochain complexes \[
    \begin{tikzcd}
      0 & \mathcal{A} & \mathcal{B} & \mathcal{C} & 0
      \arrow[from=1-1, to=1-2]
      \arrow[from=1-2, to=1-3, "f"]
      \arrow[from=1-3, to=1-4, "g"]
      \arrow[from=1-4, to=1-5]
  \end{tikzcd}\]

  There is a vector space morphism $d^*: H^k(\mathcal{C}) \to H^{k+1}(\mathcal{A})$, called the \textbf{connecting homomorphism}
\end{proposition}

\begin{customproof}
  \

  Consider the expanded view of the short exact sequence of cochain complexes:

  \[
    \begin{tikzcd}
      0 & A^{k+1} & B^{k+1} & C^{k+1} & 0\\
      0 & A^k & B^k & C^k & 0
      \arrow[from=1-1, to=1-2]
      \arrow[from=1-2, to=1-3, "f^{k+1}"]
      \arrow[from=1-3, to=1-4, "g^{k+1}"]
      \arrow[from=1-4, to=1-5]
      \arrow[from=2-1, to=2-2]
      \arrow[from=2-2, to=2-3, "f^{k}"]
      \arrow[from=2-3, to=2-4, "g^{k}"]
      \arrow[from=2-4, to=2-5]
      \arrow[from=2-2, to=1-2, "d_A"]
      \arrow[from=2-3, to=1-3, "d_B"]
      \arrow[from=2-4, to=1-4, "d_C"]
  \end{tikzcd}\]and note that the degree on each differential is implied here. Then, start with some $[c] \in H^k(\mathcal{C})$, and representative $c$ which is a cocycle. Since $g^k$ is surjective, $c = g^k(b)$ for some $b \in B^k$. Then, consider $$g^{k+1}(d_Bb) = d_C(g^k(b)) = d_Cc = 0$$So $d_Bb \in \text{ker}(g^{k+1}) = \text{im}(f^{k+1})$, so there exists some $a \in A^{k+1}$ so that $f^{k+1}(a) = d_Bb$, and since $f^{k+1}$ injective, this $a$ is unique. But $$f^{k+1}(d_A a) = d_Bf^{k+1}(a) = d_Bd_Bb = 0$$so $da = 0$ so $a$ is a cocycle and defines a cohomology class $[a]$.

  Now, just let $$d^*: [c] \mapsto [a]$$

  To show this map is well defined, we need to show that it is independent of choice of representative for $[c]$, and choice of $b$ in the preimage of $c$ under $g^k$. First, we show the former:

  Let $c, c' \in [c]$. Then, $c' = c + dx$ for some coboundary $dx$. Let $c = g^k(b)$ and $c' = g^k(b') $ for some $b, b' \in B^k$, and by linearity, this implies that $b' = b + dy$ for some coboundary $dy$. But then, $d_Bb' = d_B(b+dy) = d_Bb$, so their derivation is the same, so in our construction they will map to the same cohomology class in $H^{k+1}(\mathcal{A})$

  Now, we show the latter:

  Let $[c] \in H^k(\mathcal{C})$. Then, there exist $b, b' \in B^k$ such that $g^k(b) = c$ and $g^k(b') = c$. Then, since $g^k(b-b') = c-c = 0$, $b-b' \in \ker(g^k)= \text{im}(f^k)$ and by injectivity, there exists some $a'' \in A^k$ such that $f^k(a'') = b-b'$. Now suppose that $d^*[c] = [a]$ such that $f^{k+1}(a) = db$ and $d^*[c] = [a']$ such that $f^{k+1}(a') = db'$. Then, $$f^{k+1}(a-a') = d(b-b') = df^k(a'') = f^{k+1}(da'')$$Since $f^{k+1}$ is injective, we know that $a-a' = da''$. But then, the two representative differ by a coboundary (which we modded out), so they are equal.
\end{customproof}

\subsection{The Zig-Zag Lemma}
\

\begin{theorem}
  \

  A short exact sequence of cochain complexes \[
    \begin{tikzcd}
      0 & \mathcal{A} & \mathcal{B} & \mathcal{C} & 0
      \arrow[from=1-1, to=1-2]
      \arrow[from=1-2, to=1-3, "f"]
      \arrow[from=1-3, to=1-4, "g"]
      \arrow[from=1-4, to=1-5]
  \end{tikzcd}\]gives rise to a long exact sequence in cohomology

  \[
    \begin{tikzcd}
      H^{k+1}(\mathcal{A}) & \dots & \\
      H^k(\mathcal{A}) & H^{k}(\mathcal{B}) & H^k(\mathcal{C}) \\
      & \dots & H^{k-1}(\mathcal{C})
      \arrow[from=1-1, to=1-2, "f^*"]
      \arrow[from=2-1, to=2-2, "f^*"]
      \arrow[from=2-2, to=2-3, "g^*"]
      \arrow[from=3-2, to=3-3, "g^*"]
      \arrow[from=2-3, to=1-1, snake left, "d^*"']
      \arrow[from=3-3, to=2-1, snake left, "d^*"']
  \end{tikzcd}\]where $f^*$ and $g^*$ are the maps in cohomology induced from the cochain maps $i$ and $j$, and $d^*$ is the connecting homomorphism.
\end{theorem}

\begin{customproof}
  \

  It's a lot of diagram chasing.
\end{customproof}

\subsection{Exercises}
\

\begin{exercise}
\

Show that a commutative diagram of vector spaces with exact rows \[\begin{tikzcd}
  0 & A^1 & B^1 & C^1 & 0\\
  0 & A^0 & B^0 & C^0 & 0
  \arrow[from=1-1, to=1-2]
  \arrow[from=1-2, to=1-3, "i_1"]
  \arrow[from=1-3, to=1-4, "j_1"]
  \arrow[from=1-4, to=1-5]
  \arrow[from=2-1, to=2-2]
  \arrow[from=2-2, to=2-3, "i_0"]
  \arrow[from=2-3, to=2-4, "j_0"]
  \arrow[from=2-4, to=2-5]
  \arrow[from=2-2, to=1-2, "\alpha"]
  \arrow[from=2-3, to=1-3, "\beta"]
  \arrow[from=2-4, to=1-4, "\gamma"]
\end{tikzcd}\]incudes a long exact sequence \[\begin{tikzcd}
  & \text{coker}(\alpha) & \text{coker}(\beta) & \text{coker}(\gamma) & 0\\
  0 & \ker(\alpha) & \ker(\beta) & \ker(\gamma) & 
  \arrow[from=1-2, to=1-3]
  \arrow[from=1-3, to=1-4]
  \arrow[from=1-4, to=1-5]
  \arrow[from=2-4, to=1-2, snake left, "d"']
  \arrow[from=2-1, to=2-2]
  \arrow[from=2-2, to=2-3, "f_1"]
  \arrow[from=2-3, to=2-4, "f_2"]
\end{tikzcd}\]
\end{exercise}

\begin{customproof}
\

We want to show that $\text{im}(f_1) = \ker(f_2)$. By letting $f_1 = i_0$, we get this automatically. 

Let some element $c \in \ker(\gamma) \subseteq C^0$. Since $j_0$ surjective, there is some $b \in B^0$ such that $j_0(b) = c$. By commutativity, $$j_1 \circ \beta(b) = \gamma \circ j_0(b) = \gamma (c) = 0$$so $\beta(b) \in \ker(j_1) = \text{im}(i_1)$, so there is some $a \in A^1$ such that $i_1(a) = \beta(b)$ and it is unique by injectivity of $i_1$. By mapping $c$ to $[a] \in \text{coker}(\alpha)$ we get a morphism that is well defined, and it is exact because for any $c \in \text{im}(f_2)$, this means the associated $b$ is in $\ker(\beta)$, so $\beta(b) = 0$, so by injectivity of $i_1$, $a = 0$, hence $[a] \in \ker(d)$.
\end{customproof}

\section{The Mayer-Vietoris Sequence}

\subsection{The Mayer-Vietoris Sequence}
\

\begin{remark}
\

For a manifold $M$ want to construt a short exact sequence of de Rham complexes \[\begin{tikzcd}
  0 & \Omega^*(M) & \Omega^*(U) \oplus \Omega^*(V) & \Omega^*(U \cap V) & 0
  \arrow[from=1-1, to=1-2]
  \arrow[from=1-2, to=1-3, "i"]
  \arrow[from=1-3, to=1-4, "j"]
  \arrow[from=1-4, to=1-5]
\end{tikzcd}\] called the Mayer-Vietoris sequence, which by the zig-zag lemma, extends to a very useful result about cohomologies that will help us find the cohomologies on various manifolds.

We first want to define $i$, $j$, and show they are cochain maps, so they commute with the differential.
\end{remark}

\begin{definition}[Restriction and Difference Map]
\

Let $\{U, V\}$ be an open cover of some manifold $M$ (so $M = U \cup V$). Consider the map $\iota_U: U \to M$ the inclusion map, whose pullback is the restriction map $\iota^*_U: \Omega^k(M) \to \Omega^k(U)$, by $\omega \mapsto \omega \vert_U$.

The \textbf{restriction map} is a vector space morphism which restricts a $k$-form to both $U$ and $V$, by $$i: \Omega^k(M) \to \Omega^k(U) \oplus \Omega^k(V)$$$$i: \sigma\mapsto (\sigma\vert_U,\sigma\vert_V)$$

The \textbf{difference map} is a vector space morphism as follows: $$k: \Omega^k(U) \oplus \Omega^k(V) \mapsto \Omega^k(U \cap V)$$$$j: (\omega, \tau) \mapsto \tau\vert_{U \cap V} - \omega \vert_{U \cap V}$$ 
\end{definition}

\begin{proposition}
\

Both $i$ and $j$ commute with the exterior derivative $d$ (which is the differential on the de Rham complex)
\end{proposition}

\begin{customproof}
\

Note that the exterior derivative on $\Omega^*(U) \oplus \Omega^*(V)$ is given by $d(\omega, \tau) = (d\omega, d\tau)$.

Let $\sigma \in \Omega^k(M)$. Then, $$d(i(\sigma)) = (d\iota^*_U(\sigma), d\iota_V^*(\sigma)) = (\iota^*_U d\sigma, \iota^*_V d\sigma) = i(d\sigma)$$

Let $(\omega, \tau) \in \Omega^k(U)\oplus \Omega^k(V)$. Then, $$dj(\omega, \tau) = d(\iota_{U\cap V}^*(\omega) - \iota_{U \cap V}(\tau)) = d\iota_{U\cap V}(\omega - \tau) = \iota_{U \cap V}(d\omega - d\tau) = j(d\omega, d\tau)$$Hence they commute with the differential, and $i$, and $j$ are cochain maps.
\end{customproof}

\begin{proposition}
\

For each integer $k \geq 0$, the sequence \[\begin{tikzcd}
  0 & \Omega^k(M) & \Omega^k(U) \oplus \Omega^k(V) & \Omega^k(U \cap V) & 0
  \arrow[from=1-1, to=1-2]
  \arrow[from=1-2, to=1-3, "i"]
  \arrow[from=1-3, to=1-4, "j"]
  \arrow[from=1-4, to=1-5]
\end{tikzcd}\]is exact.
\end{proposition}

\begin{customproof}
\

WTS $i$ is injective. Suppose that there were $\omega, \tau$ such that they disagreed at at least one point $p$. Then, $i(\omega) = (\omega \vert_U, \omega \vert_v)$ and $i(\tau) = (\tau \vert_U, \tau \vert_V)$ but since $U \cup V = M$, $p$ must be in at least one $U$ or $V$, WLOG assume $U$, but then $\omega \vert_U \neq \omega \vert_V$.

WTS $j$ surjective. Let $\omega$ be a smooth $k$-form on $U \cap V$, and let $\{\rho_U, \rho_V\}$ be a partition of unity subordinate to $\{U, V\}$, and consider $$\omega_V = \begin{cases}
  \rho_U(x)\omega & x \in U \cap V\\
  0 & x \in V \setminus (U \cap V)
\end{cases}$$you can verify that this is indeed a smooth form, but then we have $\omega_V \in \Omega^k(V)$ and $\omega_U \in \Omega^k(U)$, and $\omega_V \vert_{U \cap V} - (-\omega_U \vert_{U \cap V}) = \omega$ (We want to sum them because the partition of unity equals $1$ when we sum).

Since $i$ injective, $j$ surjective, the sequence is exact. 
\end{customproof}

\begin{corollary}
\

\[\begin{tikzcd}
  0 & \Omega^*(M) & \Omega^*(U) \oplus \Omega^*(V) & \Omega^*(U \cap V) & 0
  \arrow[from=1-1, to=1-2]
  \arrow[from=1-2, to=1-3, "i"]
  \arrow[from=1-3, to=1-4, "j"]
  \arrow[from=1-4, to=1-5]
\end{tikzcd}\]is a short exact sequence of cochain complexes.
\end{corollary}

\begin{corollary}[Mayer-Vietoris Sequence]
\

By the zig-zag lemma, we have a long exact sequence of cohomologies: \[
    \begin{tikzcd}
      H^{k+1}(M) & \dots & \\
      H^k(M) & H^{k}(U) \oplus H^k(V) & H^k(U\cap V) \\
      & \dots & H^{k-1}(U \cap V)
      \arrow[from=1-1, to=1-2, "i^*"]
      \arrow[from=2-1, to=2-2, "i^*"]
      \arrow[from=2-2, to=2-3, "j^*"]
      \arrow[from=3-2, to=3-3, "j^*"]
      \arrow[from=2-3, to=1-1, snake left, "d^*"']
      \arrow[from=3-3, to=2-1, snake left, "d^*"']
  \end{tikzcd}\]called the \textbf{Mayer-Vietoris sequence} 
\end{corollary}

\begin{proposition}
\

In the Mayer-Vietoris sequence, if $U$, $V$, and $U \cap V$ are connected and nonempty, then

\begin{enumerate}
  \item $M$ is connected and \[\begin{tikzcd}
    0 & H^0(M) & H^0(U) \oplus H^0(V) & H^0(U \cap V) & 0
    \arrow[from=1-1, to=1-2]
    \arrow[from=1-2, to=1-3]
    \arrow[from=1-3, to=1-4]
    \arrow[from=1-4, to=1-5]
  \end{tikzcd}\]is exact
  \item We may start the Mayer-Vietoris sequence with \[\begin{tikzcd}
    0 & H^1(M) & H^1(U) \oplus H^1(V) & H^1(U \cap V) & \dots
    \arrow[from=1-1, to=1-2]
    \arrow[from=1-2, to=1-3, "i^*"]
    \arrow[from=1-3, to=1-4, "j^*"]
    \arrow[from=1-4, to=1-5]
  \end{tikzcd}\]
\end{enumerate}
\end{proposition}

\begin{customproof}
\


\begin{enumerate}
  \item If $U, V, U \cap V$ are connected and nonempty and $M = U \cup V$, then $M$ connected. On a connected space $C$, $H^0(C) = \mathbb{R}$, so $j^*: \mathbb{R} \oplus \mathbb{R} \to \mathbb{R}$ by $(u,v) \mapsto v-u$ is surjective, and $i^*: \mathbb{R} \to \mathbb{R} \oplus \mathbb{R}$ by $a \mapsto (a,a)$ is also injective, hence the sequence is exact.
  \item Since the Mayer-Vietoris sequence is exact, $\text{im} j^* = \ker d^*$, but $\text{im} j^* = \mathbb{R} = \ker d^*$, and $d^*: \mathbb{R} \to H^1(M)$, so $d^*$ must map its entire domain into $0$, hence it is the zero map. Then, this is indistinguishable from just starting with $0 \to H^1(M)$.
\end{enumerate}
\end{customproof}

\subsection{The Cohomology of the Circle}
\

\begin{proposition}
\

Let \[\begin{tikzcd}
  0 & A^0 & A^1 & A^2 & \dots & A^m & 0
  \arrow[from=1-1, to=1-2]
  \arrow[from=1-2, to=1-3, "d_0"]
  \arrow[from=1-3, to=1-4, "d_1"]
  \arrow[from=1-4, to=1-5, "d_2"]
  \arrow[from=1-5, to=1-6]
  \arrow[from=1-6, to=1-7]
\end{tikzcd}\]be a short exact sequence of finite-dimensional vector spaces. Then, $$\sum_{k=0}^m (-1)^k \dim A^k = 0$$
\end{proposition}

\begin{customproof}
\

The rank nullity theorem states that $\dim A^k = \dim \ker(d_k) + \dim \text{im}(d_k)$. Then, consider $$\sum_{k=0}^m (-1)^k \dim A^k = \sum_{k=0}^m (-1)^k (\dim \ker(d_K) + \dim \text{im}(d_k)) = \sum_{k=0}^m (-1)^k (\dim \text{im}(d_{k-1}) + \dim \text{im}(d_k)) = \dim \text{im}(d_m) = 0$$
\end{customproof}

\begin{example}
\

We will illustrate the power of the Mayer-Vietoris sequence. Before, we had to do a lot of integration and differential equations to find $H^1(S^1)$, but it is much easier with this tool:

Start with $S^1$, cover the circle with two open arcs, $U$, $V$. They will overlap into two disjoint open arcs, $U \cap V$. Then, we have the Mayer-Vietoris sequence \[\begin{tabular}{c | c c c}
  & $S^1$ & $U \sqcup V$ & $U \cap V$\\
  \hline
  $H^2$ & $0$ & $0$ & $0$\\
  $H^1$ & $H^1(S^1)$ & $0$ & $0$\\
  $H^0$ & $\mathbb{R}$ & $\mathbb{R} \oplus \mathbb{R}$ & $\mathbb{R}^2$
\end{tabular}\]If we just look at the segment \[\begin{tikzcd}
  \mathbb{R} \oplus \mathbb{R} & \mathbb{R}^2 & H^1(S^1) & 0
  \arrow[from=1-1, to=1-2, "f"]
  \arrow[from=1-2, to=1-3, "g"]
  \arrow[from=1-3, to=1-4, "h"]
\end{tikzcd}\]We have $$H^1(S^1) = \text{im}(g) \cong \frac{\mathbb{R}^2}{\ker(g)} = \frac{\mathbb{R}^2}{\text{im}(f)}$$but since the mapping before $f$ in the sequence took $\mathbb{R} \to \mathbb{R}^2$ and had image $\mathbb{R}$, the kernel of $f$ is $\mathbb{R}$, so the image of $f$ is $\mathbb{R}$ as well, so we have $$H^1(S^1) = \frac{\mathbb{R}^2}{\mathbb{R}} = \mathbb{R}$$
\end{example}

\subsection{Euler Characteristic}
\

\begin{definition}[Euler Characteristic]
\

If $M^m$ is a manifold and every $H^k(M)$ is finite-dimensional for all $k$, the \textbf{Euler characteristic} is given by $$\chi(M) = \sum_{k=0}^m (-1)^k \dim H^k(M)$$
\end{definition}

\section{Homotopy Invariance}

\subsection{Smooth Homotopy}
\

\begin{definition}[Smooth Homotopy]
\

Let $M$ and $N$ be manifolds. Two smooth maps $f, g: M \to N$ are \textbf{smoothly homotopic} if there is a smooth map $F: M \times \mathbb{R} \to N$ such that $F(x, 0) = f(x)$ and $F(x,1) = g(x)$ for all $x \in M$. The map $F$ is called a \textbf{homotopy} from $f$ to $g$.

If $f$ and $g$ are homotopic, we write $f \sim g$.
\end{definition}

\subsection{Homotopy Type}
\

\begin{definition}[Homotopy Equivalence]
\

A map $f: M \to N$ is a \textbf{homotopy equivalence} if it has \textbf{homotopy inverse}, so there exists a map $g: N \to M$ such that $g\circ f \sim 1_M$ is homotopic to the identity on $M$, and $f \circ g \sim 1_N$.

Then, we say $M$ is \textbf{homotopy equivalent} to $N$, and they have the same \textbf{homotopy type}  
\end{definition}

\begin{example}
\

A diffeomorphism is a (albeit trivial) homotopy equivalence
\end{example}

\begin{definition}[Contractible Manifolds]
\

A manifold is \textbf{contractible} if it is homotopy equivalent to a point. 
\end{definition}

\begin{example}
\

$\mathbb{R}^n$ is contractible. Let $p \in \mathbb{R}^n$, $\iota: \{p \} \to \mathbb{R}^n$ the inclusion map, and $r: \mathbb{R}^n \to\{p\}$. Then, $r \circ i = 1_{\{p\}}$ and the straight line homotopy provides a homotopy between $\iota \circ r: \mathbb{R}^n \to \mathbb{R}^n$ and the identity map on $\mathbb{R}^n$, by $$F(x,t) = (1-t)x + tr(x) = (1-t)x + tp$$
\end{example}

\subsection{Deformation Retractions}
\

\begin{definition}[Retractions]
\

Let $S \subseteq M$ and $\iota: S \to M$ be the inclusion map. Then, a \textbf{retraction} from $M$ to $S$ is a map $r: M \to S$ that restricts to the identity on $S$, so $r \circ \iota = 1_S$. 

If there is a retraction from $M$ to $S$, we say $S$ is a \textbf{retract} of $M$ 
\end{definition}

\begin{definition}[Deformation Retractions]
\

A \textbf{deformation retraction} from $M$ to $S$ is a map $F: M \times \mathbb{R} \to M$ such that for all $x \in M$, \begin{enumerate}
  \item $F(x,0) = x$
  \item there is a retraction $r: M \to S$ such that $F(x,1) = r(x)$
  \item for all $s \in S$, $t \in \mathbb{R}$, $F(s,t) = s$.
\end{enumerate} 

If there is a deformation retraction from $M$ to $S$, we say $S$ is a \textbf{deformation retract} of $M$.  
\end{definition}

\begin{proposition}
\

If $S \subseteq M$ is a deformation retraction of $M$, then $S$ and $M$ have the same homotopy type.
\end{proposition}

\begin{customproof}
\

Let $F: M \times \mathbb{R} \to M$ be a deformation retraction, $r(x) = F(x,1)$ be the retraction. Then, $r \circ \iota_S = 1_S$, and $\iota \circ r = r$, but the deformation retraction gives a homotopy between $r$ and $1_M$. Hence, $r: M \to S$ is a homotopy equivalence with homotopy inverse $\iota$
\end{customproof}

\subsection{The Homotopy Axiom for de Rham Cohomology}
\

\begin{theorem}[Homotopy Axiom for de Rham Cohomology]
\

Homotopic maps $f_0, f_1: M \to N$ induce the same map $f_0^* = f_1^*: H^*(N)\to H^*(M)$ in cohomology.
\end{theorem}

\begin{corollary}
\

If $f: M \to N$ is a homotopy equivalence, then the induced map in cohomology, $f^*: H^*(N) \to H^*(M)$ is an isomorphism.
\end{corollary}

\begin{customproof}
\

Let $g: N \to M$ be homotopy inverse to $f$, so $g \circ f \sim 1_M$ and $f \circ g \sim 1_N$. But then, $(g \circ f)^* = 1_{H^*(M)}$ and $(f \circ g)^* = 1_{H^*(N)}$. But then, $$f^* \circ g^* = (g \circ f)^* = 1_{H^*(M)}$$and $$g^* \circ f^* = 1_{H^*(N)}$$so we have a set of inverses, hence they are isomorphic.
\end{customproof}

\begin{corollary}
\

Suppose $S$ is a submanifold of $M$ and $F$ a deformation retraction from $M$ to $S$. Let $F(x,1)=r: M \to S$, then $r$ induces an isomorphism in cohomology $$r^*: H^*(S) \to H^*(M)$$
\end{corollary}

\begin{corollary}[Poincare Lemma]
\

Since $\mathbb{R}^n$ has the homotopy type of a point, the cohomology of $\mathbb{R}^n$ is $$H^k(\mathbb{R}^n) = \begin{cases}
  \mathbb{R} & k=0\\
  0 & k > 0
\end{cases}$$
\end{corollary}

\section{Computation of de Rham Cohomology}

\subsection{Cohomology Vector Space of a Torus}
\

\begin{example}
\



Consider a torus, cover it with two open macaroni-shaped subsets, $U, V$. Then, $U \cap V$ is two disjoint cylinders. Then, $U, V$ diffeomorphic to a cylinder, which is homotopy equivalent to a circle, and $U \cap V$ is diffeomorphic to two circles.

We then have the Mayer-Vietoris Sequence \[\begin{tabular}{c | c c c}
  & $M$ & $U \sqcup V$ & $U \cap V$\\
  \hline
  $H^2$ & $H^2(M)$ & $0$ & $0$\\
  $H^1$ & $H^1(M)$ & $ \mathbb{R}^2 $ & $\mathbb{R}^2$\\
  $H^0$ & $\mathbb{R}$ & $\mathbb{R} \oplus \mathbb{R}$ & $\mathbb{R}^2$
\end{tabular}\]and expanded we have \[\begin{tikzcd}
  0 & \mathbb{R} & \mathbb{R} \oplus \mathbb{R} & \mathbb{R}^2 & H^1(M) & \mathbb{R} \oplus \mathbb{R} & \mathbb{R}^2 & H^2(M) & 0
  \arrow[from=1-1, to=1-2]
  \arrow[from=1-2, to=1-3, "i^*_0"]
  \arrow[from=1-3, to=1-4, "\alpha"]
  \arrow[from=1-4, to=1-5, "d^*_0"]
  \arrow[from=1-5, to=1-6, "i^*_1"]
  \arrow[from=1-6, to=1-7, "\beta"]
  \arrow[from=1-7, to=1-8, "d_1^*"]
  \arrow[from=1-8, to=1-9]
\end{tikzcd}\]

Lets try to characterize $\alpha$, which is the difference map $j^*_0$. Let $a \in H^0(U)$ be some constant function. When restricted to $U \cap V$ each function on each of the two components now takes on the same value, so $a \vert_{U \cap V} = (a,a)$, and same for $b \in H^0(V)$. Then, the difference map gives us $$\alpha(a,b) = j^*_0(a,b) = b\vert_{U \cap V} - a\vert_{U \cap V} = (b-a, b-a)$$

For $\beta: H^1(U) \oplus H^1(V) \to H^1(U \cap V) = H^1(A) \oplus H^1(B)$, where $A, B$ are the two cylinders. Since $A$ is diffeomorphic to $U$ (they are both cylinders), $A$ is a deformation retract of $U$, so by the homotopy axioms, $H^*(U) \to H^*(A)$ is an isomorphism. Then, if $\omega_U$ generates $H^(U)$, then $j^*_{U, 1} \omega_U$ is a generator of $H^1$ on $A$ and $B$.

Since $H^1(U \cap V) \cong \mathbb{R} \oplus \mathbb{R}$, we can write $j^*_{U, 1} = (1,1)$. And if $\omega_V$ generates $H^1(V)$, we have that $(a,b)$ stands for $(a \omega_U, b \omega_V)$, and so $$\beta(a,b) = (b,b) - (a,a) = (b-a, b-a)$$


Since the Mayer Vietoris sequence is exact, $$H^2(M) = \text{im}(d_1^*) \cong \frac{H^1(U \cap V)}{\ker d_1^*} \cong \frac{(\mathbb{R} \oplus \mathbb{R})}{\text{im}(\beta)} \cong \frac{\mathbb{R} \oplus \mathbb{R}}{\mathbb{R}} \cong \mathbb{R}$$

Similarly, $$H^1(M) \cong \text{ker}(i^*_1) \oplus \text{im}(i_1^*) \cong \text{im}(d_0^*) \oplus \ker \beta \cong \frac{H^0(U \cap V)}{\ker (d_0^*)} \oplus \ker(\beta) \cong \mathbb{R} \oplus \mathbb{R}$$
\end{example}

\subsection{The Cohomology Ring of a Torus}
\

\begin{example}
\

If $\Lambda = \mathbb{Z}^2$ the integer lattice, then a torus is $\mathbb{R}^2/\Lambda$, and there is a projection map $$\pi: \mathbb{R}^2 \to \mathbb{R}^2/\Lambda$$and a corresponding pullback $$\pi^*: \Omega^*(\mathbb{R}^2/\Lambda) \to \Omega^*(\mathbb{R}^2)$$Since $\pi$ is a local diffeomorphism, the differential $\pi_*: T_q(\mathbb{R}^2) \to T_{\pi(q)}(\mathbb{R}^2/\Lambda)$ is an isomorphism at each $q \in \mathbb{R}^2$, and in particular, $\pi$ is a submersion, and so by a prior problem, $\pi^*: \Omega^*(\mathbb{R}^2/\Lambda) \to \Omega^*(\mathbb{R}^2)$ is an injection. 
\end{example}

\begin{definition}[Translational Invariance]
\

For some $\lambda \in \Lambda$, let $\ell_\lambda: \mathbb{R}^2 \to \mathbb{R}^2$ to be translation by $\lambda$, so $$\ell_\lambda(q) = q + \lambda$$A differential form $\bar{\omega}$ is \textbf{invariant under translation} if $\ell_\lambda^*\bar{\omega} = \bar{\omega}$. 
\end{definition}

\begin{proposition}
\

The image of the injection $\pi^*: \Omega^*(\mathbb{R}^2/\Lambda) \to \Omega^*(\mathbb{R}^2)$ is the subspace of differential forms invariant under translation by elements of $\lambda$.
\end{proposition}

\begin{customproof}
\

For all $q \in \mathbb{R}^2$, $$(\pi \circ \ell_\lambda) = \pi(q + \lambda) = \pi(q)$$since we modded out by elements of $\lambda$. Then, $\pi \circ \ell_\lambda = \pi$, so $(\pi \circ \ell_\lambda)^* = \pi^*$ so $\ell_\lambda^* \circ \pi^* = \pi^*$ so for any $\omega \in \Omega^k(\mathbb{R}^2/\Lambda)$, we have $\pi^*\omega = \ell_\lambda^*\pi^*\omega$, hence the image of $\Omega^*(\mathbb{R}^2/\Lambda)$ under $\pi^*$ is the space of translationally invariant forms.
\end{customproof}

\begin{proposition}
\

Let $M$ be the torus $\mathbb{R}^2/\mathbb{Z}^2$. A basis for the cohomology vector space $H^*(M)$ is represented by the forms $1, \alpha, \beta, \alpha \wedge \beta$, where $\alpha$ is the preimage of $dx$ under $\pi^*$, and $\beta$ is the preimage of $dy$.
\end{proposition}

\subsection{The Cohomology of a Surface of Genus $g$}
\

\begin{lemma}
\

Suppose $p$ is a point in a compact oriented surface $M$ without boundary, and $i: C \to M \setminus \{p\}$ is the inclusion of a small circle around the puncutre. Then, the restriction map $$i^* H^1(M \setminus \{p\}) \to H^1(C)$$is the zero map.
\end{lemma}

\begin{customproof}
\

An element $[\omega] \in H^1(M \setminus \{p\})$ is represented by a closed $1$-form $\omega$ on $M \setminus \{p\}$, but $H^1(C) \cong H^1(S^1) \cong \mathbb{R}$ is given by integration over $C$, so to find $i^*[\omega]$, we just need to compute the integral. Let $D$ be an open disk in $M$ bounded by $C$, $M\setminus D$ a compact oriented surface $$\int_C i^*\omega = \int_{\partial (M \setminus D)} i^* \omega = \int_{M-D} d\omega = 0$$so $i^*$ is the zero map.
\end{customproof}

\begin{proposition}
\

Let $M$ be a torus, $p$ a point in $M$, and $A$ the punctured torus $M \setminus \{p\}$. The cohomology of $A$ is $$H^k(A) = \begin{cases}
  \mathbb{R} & k=0\\
  \mathbb{R}^2 & k=1\\
  0 & k > 1
\end{cases}$$
\end{proposition}


\begin{proposition}
\

The cohomology of a compact orientable surface $\Sigma_2$ of genus $2$ is $$H^k(\Sigma_2) = \begin{cases}
  \mathbb{R} & k = 0,2\\
  \mathbb{R}^4 & k = 1\\
  0 & k > 2
\end{cases}$$
\end{proposition}

\subsection{Exercises}
\

\begin{exercise}[1]
\

Compute the cohomology of $\mathbb{RP}^2$
\end{exercise}

\begin{solution}
\

$\mathbb{RP}^2$ is given by $\mathbb{R}^3$ mod the equivalence relation that $(a,b,c) \sim \lambda(a,b,c)$ for some real number $\lambda \in \mathbb{R}$. Denote equivalence classes of points as $[x,y,z]$. Let $U = \mathbb{RP}^2 \setminus \{[1,0,1]\}$ and $V = \mathbb{RP}^2 \setminus \{-1, 0, 1\}$. Then, \[\begin{tabular}{c | c c c}
  & $\mathbb{RP}^2$ & $U \sqcup V$ & $U \cap V$\\
  \hline
  $H^2$ & $H^2(\mathbb{RP}^2)$ & $0$ & $0$\\
  $H^1$ & $H^1(\mathbb{RP}^2)$ & $ \mathbb{R}^2 $ & $\mathbb{R}^2$\\
  $H^0$ & $\mathbb{R}$ & $\mathbb{R} \oplus \mathbb{R}$ & $\mathbb{R}$
\end{tabular}\]Looking at $\alpha: \mathbb{R} \to H^1(\mathbb{RP}^2)$, $\beta: H^1(\mathbb{RP}^2) \to \mathbb{R}^2$, since we know that the image of the prior map is $\mathbb{R}$ (by rank nullity, since the kernel of the prior map is $\mathbb{R}$, the image must also be $\mathbb{R}$), we get that $\text{ker}(\alpha) = \mathbb{R}$, so $\text{im}(\alpha) = \ker(\beta) = \mathbb{R}$. 
\end{solution}


\section{Proof of Homotopy Invariance}

\subsection{Reduction to Two Sections}
\

\begin{remark}
\

In this section, if $f: M \to N$ is a smooth map, let $f^*: \Omega^k(N) \to \Omega^k(M)$ and $f^\#: H^k(N) \to H^k(M)$, so that $f^\#[\omega] = [f^*\omega]$
\end{remark}

\begin{theorem}[Homotopy Axiom for the de Rham Cohomology]
\

Two smoothly homotopic maps $f,g : M \to N$ of manifolds induce the same map in cohomology, $$f^\# = g^\#: H^k(N) \to H^k(M)$$
\end{theorem}

\begin{customproof}
\

Suppose $f, g: M \to N$ are smoothly homotopic maps. Let $F: M \times \mathbb{R} \to N$ be a smooth homotopy from $f$ to $g$, so $F(x,0) = f(x)$ and $F(x,1) = g(x)$ for all $x \in M$. Let $i_t: M \to M \times \mathbb{R}$ take $i_t(x) = (x,t)$, so $F \circ i_0 = f$ and $F \circ i_1 = g$.

Then, $f^\# = i_0^\# \circ F^\#$ and $g^\# = i^\#_1 \circ F^\#$. Then, we just want to prove that $i_0^\# = i_1\#$. Note that $i_0 \sim i_1$ just by virtue of the identity map $1_{M \times \mathbb{R}}$, since at $t = 0$, we get $i_0$, and at $t=1$, we get $i_1$.
\end{customproof}

\subsection{Cochain Homotopies}
\

\begin{definition}[Cochain Homotopies]
\

The usual method for showing two cochain maps $\varphi, \psi: \mathcal{A} \to \mathcal{B}$ induce the same map in cohomologies is to find some $K: \mathcal{A} \to \mathcal{B}$ of degree $-1$ such that $\varphi - \psi = d \circ K + K \circ d$, and this map $K$ is called the \textbf{cochain homotopy}.

THis is because if $a \in \mathcal{A}$ is a cocycle, then $\varphi(a) - \psi(a) = dKa + Kda = dKa$ is a coboundary, so it gets modded out. 
\end{definition}

\subsection{Differential Forms on $M \times \mathbb{R}$}
\

\begin{definition}[Locally Finite Sums]
\

A sum $\sum_{\alpha} \omega_\alpha$ of smooth differential forms is \textbf{locally finite} if the collection $\{\text{supp}(\omega_\alpha)\}$ has only finitely many nonzero elements at each fixed $p$. 
\end{definition}

\begin{lemma}
\

Let $U$ be an open subset of a manifold $M$. If a smooth $k$-form $\tau \in \Omega^k(U)$ has support in a closed subset of $M$ contained in $U$, then $\tau$ can be extended by zero to a smooth $k$-form on $M$.
\end{lemma}

\begin{proposition}
\

Let $\pi: M \times \mathbb{R} \to M$ be the projection map. Then, every smooth dfferential form $\omega$ on $M \times \mathbb{R}$ is a locally finite sum of the following two types of forms:

\begin{enumerate}
  \item $f(x,t)\pi^*\eta$
  \item $f(x,t) dt \wedge \pi^*\eta$
\end{enumerate}for some $\eta \in \Omega^*(M)$ and $f(x,t): M \times \mathbb{R} \to \mathbb{R}$ a smooth function.
\end{proposition}

\begin{customproof}
\

Let $\{(U_\alpha, \varphi_\alpha)\}$ be an atlas on $M$, let $\{\rho_\alpha\}$ be a partition of unity subordinate to $U_\alpha$, and let $g_\alpha$ be a collection of smooth functions on $M$ such that $g_\alpha = 1$ on $\text{supp}(\rho_\alpha)$ and $\text{supp}(g_\alpha) \subseteq U_\alpha$. Such a $g_\alpha$ exists via the smooth Uryshon lemma. Then, $\{\pi U_\alpha\}$ is a cover of $M \times \mathbb{R}$ and $\{\pi^* \rho_\alpha\}$ is a partition of unity subordinate to it. 

Let $\omega$ be any smooth $k$-form on $M \times \mathbb{R}$ and let $\omega_\alpha = (\pi^*\rho_\alpha) \omega$. Since $\sum \pi^* \rho_\alpha = 1$, $\omega = \sum_{\alpha} \omega_\alpha$. Since $\{\text{supp} (\rho_\alpha)\}$ is locally finite, the above is a locally finite sum, so $$\text{supp}(\omega_\alpha) \subset \text{supp}(\pi^*\rho_\alpha) \cap \text{supp}(\omega) \subset \text{supp}(\pi^*\rho_\alpha) \subset \pi ^{-1}U_\alpha$$On $\pi ^{-1}U_\alpha$, we have coordinates $\pi^*x^1, \dots, \pi^*x^n, t$, just call them $x^i$. On $\pi ^{-1}U_\alpha$, the $k$-form $\omega_\alpha$ can be written uniquely as a linear combination $$\omega_\alpha = \sum_I a_Idx^I + \sum_J b_jdt \wedge dx^J$$so $\omega_\alpha$ can be written as a finite sum of the two types of forms above, on $\pi ^{-1}U_\alpha$.

To extend the decomposition to $M \times \mathbb{R}$, we multiply by $\pi^*g_alpha$, we have equality since $\text{supp}(g_\alpha) \subset U_\alpha$. We get $$\omega_\alpha = \sum a_I \pi^*(g_\alpha dx^I) + \sum_J b_jdt\wedge \pi^*(g_\alpha dx^J)$$
\end{customproof}

\subsection{A Cochain Homotopy Between $i_0^*$ and $i_1^*$}
\

\begin{proposition}
\

If $\omega \in \Omega^k(M \times \mathbb{R})$, by the prior lemma, let $$\omega = \sum_{\alpha, I} a_I^\alpha \pi^*(g_\alpha dx^I_\alpha) + \sum_{\alpha, J} b_J^\alpha dt\wedge \pi^*(g_\alpha dx^J_\alpha)$$and$$K(\omega) = K(\sum_{\alpha} \omega_\alpha) = \sum_{\alpha_J}(\int_0^1 b_j^\alpha(x,t) dt) g_\alpha dx^J_\alpha$$is a cochain homotopy between $i_0^*$ and $i_1^*$
\end{proposition}

\begin{customproof}
\

We want $K$ to obey the following rules: 

\begin{enumerate}
  \item $K(f\pi^*\eta) = 0$
  \item $K(fdt\wedge \pi^*\eta) = (\int_0^1f(x,t)dt)\eta$
\end{enumerate}

And just applying them gives us our answer.
\end{customproof}

\subsection{Verification of Cochain Homotopy}
\

\begin{lemma}
\

\begin{enumerate}
  \item The exterior derivative $d$ is $\mathbb{R}$-linear over locally finite sums
  \item Pullback by a smooth map is $\mathbb{R}$-linear over locally finite sums
\end{enumerate}
\end{lemma}

\begin{remark}
\

The two above verifications should be sufficient to verify that $K(\omega)$ is indeed a cochain homotopy.
\end{remark}

\subsection{Exercises}
\

\begin{exercise}[1]
\

Let $U$ be an open subset of a manifold $M$. If a smooth $k$-form $\tau \in \Omega^k(U)$ defined on $U$ has support in a closed subset of $M$ contained in $U$, then $\tau$ can be extended by zero to a smooth $k$-form on $M$
\end{exercise}

\begin{solution}
\

Let $\rho_U$ be a bump function, $\text{supp}(\rho_U) \subset U$. Then, $$\omega' = \begin{cases}
  \omega_p\rho(p) & p \in U\\
  0 & p \notin U
\end{cases}$$Now, we want to show $\omega'$ is smooth. Let $(V, x^i)$ be any chart on $M$. If $ \text{supp}(\rho) \cap V = \varnothing$, then on $V$, $\omega' = \sum_I 0 dx^I$, and $0$ is smooth, so it is smooth on $V$. If $\text{supp}(\rho) \cap V \neq \varnothing$, then on points not supported by $\rho$, it is identically zero anyways. On point supported by $\rho$, both $\rho$ and $a^I$ the $I$th component function are smooth, and the product of two smooth functions is smooth, so all the components everwhere are smooth, and $\omega'$ is a smooth $k$-form on $M$.
\end{solution}

\begin{exercise}[2]
\

Let $h: M \to N$ be a smooth map, $\sum \omega_\alpha$ a locally finite sum of smooth $k$ forms on $N$. Then, show that $$h^*(\sum \omega_\alpha) = \sum h^*\omega_\alpha$$
\end{exercise}

\begin{solution}
\

Fix some $p \in M$, then $$h^*(\sum \omega_\alpha)_p = h^*(\sum \omega_{\alpha, f(p)})$$but then, for fixed $p$, the space of $\omega_{f(p)}$ is a vector space, and $h^*$ is a vector space morphism, so it is linear, so we can commute to get $$h^*(\sum \omega_{\alpha, f(p)}) = \sum h^*\omega_{\alpha, f(p)} = \sum h^*(\omega_\alpha)_p$$and if this is true of all points $p \in M$, then $$h^*(\sum \omega_\alpha) = \sum h^*\omega_\alpha$$
\end{solution}

\begin{exercise}[3]
\

\begin{enumerate}
  \item Check that if $\omega = \sum_{\alpha, I} a_I^\alpha \pi^*(g_\alpha dx^I_\alpha) + \sum_{\alpha, J} b_J^\alpha dt \wedge \pi^*(g_\alpha dx^J_\alpha)$, then $$K(\omega) = \sum_{\alpha, J} (\int_0^1 b_J^\alpha(x,t) dt)g_\alpha dx_\alpha^J$$satisfies the following three rules: $K(f\pi^*\eta) = 0$ on type $1$-forms, $K(f dt \wedge \pi^*\eta) = (\int_0^1 f(x,t)dt)\eta$ on type $2$ forms, and $K$ is linear over locally finite sums.
  \item Prove that the linear operator that satisfies the above criteria is unique, if it exists.
\end{enumerate}
\end{exercise}

\begin{solution}
\

\begin{enumerate}
  \item A type $1$-form is a form of formula: $$\omega_1 = f(x,t)\pi^*\eta$$but this form has formula of $a_I^\alpha \pi^*(g_\alpha dx^*_\alpha)$, with no $b_J^\alpha$ term, hence $\sum_{\alpha, J}(\int_0^1 b_J^\alpha(x,t)dt)g_\alpha dx^J_\alpha = 0$. By matching coefficients for a type $2$ form, we get that $f = f$, $\eta = g_\alpha dx^J_\alpha$, and so $K(f dt \wedge \pi^*\eta) = (\int_0^1 f(x,t) dt)\eta$. Finally, for real number $r$ and two forms $\omega, \tau$, we have $$K(r\omega + \tau) = \sum_{\alpha, J}(\int_0^1 rb_{J, \omega}^\alpha + b_{J, \tau}^\alpha dt)g_\alpha dx^J_\alpha$$but since the integral is linear, this is linear as well.
  \item There is only one way that for a type 2 form, $K(f dt \wedge \pi^* \eta) = (\int_0^1 f(x,t) dt)\eta$, and likewise, there is only one way for $K(f\pi^*\eta) = 0$ for all type $1$ forms, and being linear forces everything else, hence if this formula exists, it is unique, since all forms are a linear combination of type $1$ and $2$ forms.
\end{enumerate}

\end{solution}

\end{document}